\documentclass[11pt,openany]{book}
\usepackage[T1]{fontenc}
\usepackage[utf8]{inputenc}
\usepackage{lmodern}
\usepackage{amsmath,amssymb,amsthm,mathtools}
\usepackage{graphicx}
\usepackage[margin=1in]{geometry}
\usepackage[hidelinks]{hyperref}

\newtheorem{theorem}{Theorem}[section]
\newtheorem{lemma}[theorem]{Lemma}
\newtheorem{corollary}[theorem]{Corollary}
\newtheorem{proposition}[theorem]{Proposition}

% Match the source's section/theorem numbering inside each chapter.
\renewcommand{\thesection}{\arabic{section}}

\title{Real Analysis\\Chapters 1, 2, 3, and 6}
\author{Elias M. Stein and Rami Shakarchi}
\date{}

\begin{document}
\maketitle
\tableofcontents
\clearpage

% Stein & Shakarchi, Real Analysis (Princeton Lectures in Analysis III)
% Chapter 1: Measure Theory
% Transcribed from the user-provided PDF.
% Exercises and Problems are intentionally omitted.
% This file is intended to be \input{} into a book/report document.
% Suggested packages in the parent document:
%   amsmath, amssymb, amsthm, mathtools
% Suggested theorem declarations:
%   \newtheorem{theorem}{Theorem}[chapter]
%   \newtheorem{lemma}[theorem]{Lemma}
%   \newtheorem{corollary}[theorem]{Corollary}
%   \newtheorem{proposition}[theorem]{Proposition}
%
% The original chapter labels many statements as ``Property'',
% ``Observation'', and ``Example''; these are kept as bold run-in headings.

\chapter{Measure Theory}

\begin{quote}
The sets whose measure we can define by virtue of the preceding ideas we will call measurable sets; we do this without intending to imply that it is not possible to assign a measure to other sets.

\hfill E. Borel, 1898
\end{quote}

This chapter is devoted to the construction of Lebesgue measure in $\mathbb{R}^d$ and the study of the resulting class of measurable functions. After some preliminaries we pass to the first important definition, that of exterior measure for any subset $E$ of $\mathbb{R}^d$. This is given in terms of approximations by unions of cubes that cover $E$. With this notion in hand we can define measurability and thus restrict consideration to those sets that are measurable. We then turn to the fundamental result: the collection of measurable sets is closed under complements and countable unions, and the measure is additive if the subsets in the union are disjoint.

The concept of measurable functions is a natural outgrowth of the idea of measurable sets. It stands in the same relation as the concept of continuous functions does to open (or closed) sets. But it has the important advantage that the class of measurable functions is closed under pointwise limits.

\section{Preliminaries}

We begin by discussing some elementary concepts which are basic to the theory developed below.

The main idea in calculating the ``volume'' or ``measure'' of a subset of $\mathbb{R}^d$ consists of approximating this set by unions of other sets whose geometry is simple and whose volumes are known. It is convenient to speak of ``volume'' when referring to sets in $\mathbb{R}^d$; but in reality it means ``area'' in the case $d=2$ and ``length'' in the case $d=1$. In the approach given here we shall use rectangles and cubes as the main building blocks of the theory: in $\mathbb{R}$ we use intervals, while in $\mathbb{R}^d$ we take products of intervals. In all dimensions rectangles are easy to manipulate and have a standard notion of volume that is given by taking the product of the length of all sides.

Next, we prove two simple theorems that highlight the importance of these rectangles in the geometry of open sets: in $\mathbb{R}$ every open set is a countable union of disjoint open intervals, while in $\mathbb{R}^d$, $d\ge 2$, every open set is ``almost'' the disjoint union of closed cubes, in the sense that only the boundaries of the cubes can overlap. These two theorems motivate the definition of exterior measure given later.

We shall use the following standard notation. A point $x\in\mathbb{R}^d$ consists of a $d$-tuple of real numbers
\[
 x=(x_1,x_2,\ldots,x_d),\qquad x_i\in\mathbb{R},\quad i=1,\ldots,d.
\]
Addition of points is componentwise, and so is multiplication by a real scalar. The norm of $x$ is denoted by $|x|$ and is defined to be the standard Euclidean norm given by
\[
 |x|=(x_1^2+\cdots+x_d^2)^{1/2}.
\]
The distance between two points $x$ and $y$ is then simply $|x-y|$.

The complement of a set $E$ in $\mathbb{R}^d$ is denoted by $E^c$ and defined by
\[
 E^c=\{x\in\mathbb{R}^d:x\notin E\}.
\]
If $E$ and $F$ are two subsets of $\mathbb{R}^d$, we denote the complement of $F$ in $E$ by
\[
 E-F=\{x\in\mathbb{R}^d:x\in E\text{ and }x\notin F\}.
\]
The distance between two sets $E$ and $F$ is defined by
\[
 d(E,F)=\inf |x-y|,
\]
where the infimum is taken over all $x\in E$ and $y\in F$.

\subsection*{Open, closed, and compact sets}

The open ball in $\mathbb{R}^d$ centered at $x$ and of radius $r$ is defined by
\[
 B_r(x)=\{y\in\mathbb{R}^d:|y-x|<r\}.
\]
A subset $E\subset\mathbb{R}^d$ is open if for every $x\in E$ there exists $r>0$ with $B_r(x)\subset E$. By definition, a set is closed if its complement is open.

We note that any (not necessarily countable) union of open sets is open, while in general the intersection of only finitely many open sets is open. A similar statement holds for the class of closed sets, if one interchanges the roles of unions and intersections.

A set $E$ is bounded if it is contained in some ball of finite radius. A bounded set is compact if it is also closed. Compact sets enjoy the Heine--Borel covering property:
\begin{itemize}
\item Assume $E$ is compact, $E\subset\bigcup_\alpha O_\alpha$, and each $O_\alpha$ is open. Then there are finitely many of the open sets, $O_{\alpha_1},O_{\alpha_2},\ldots,O_{\alpha_N}$, such that
\[
 E\subset\bigcup_{j=1}^N O_{\alpha_j}.
\]
\end{itemize}
In words, any covering of a compact set by a collection of open sets contains a finite subcovering.

A point $x\in\mathbb{R}^d$ is a limit point of the set $E$ if for every $r>0$, the ball $B_r(x)$ contains points of $E$. This means that there are points in $E$ which are arbitrarily close to $x$. An isolated point of $E$ is a point $x\in E$ such that there exists an $r>0$ where $B_r(x)\cap E$ is equal to $\{x\}$.

A point $x\in E$ is an interior point of $E$ if there exists $r>0$ such that $B_r(x)\subset E$. The set of all interior points of $E$ is called the interior of $E$. Also, the closure $\overline E$ of $E$ consists of the union of $E$ and all its limit points. The boundary of a set $E$, denoted by $\partial E$, is the set of points which are in the closure of $E$ but not in the interior of $E$.

Note that the closure of a set is a closed set; every point in $E$ is a limit point of $E$; and a set is closed if and only if it contains all its limit points. Finally, a closed set $E$ is perfect if $E$ does not have any isolated points.

\subsection*{Rectangles and cubes}

A (closed) rectangle $R$ in $\mathbb{R}^d$ is given by the product of $d$ one-dimensional closed and bounded intervals
\[
 R=[a_1,b_1]\times[a_2,b_2]\times\cdots\times[a_d,b_d],
\]
where $a_j\le b_j$ are real numbers, $j=1,2,\ldots,d$. In other words,
\[
 R=\{(x_1,\ldots,x_d)\in\mathbb{R}^d:a_j\le x_j\le b_j\text{ for all }j=1,2,\ldots,d\}.
\]
We remark that in our definition, a rectangle is closed and has sides parallel to the coordinate axis. In $\mathbb{R}$, the rectangles are precisely the closed and bounded intervals, while in $\mathbb{R}^2$ they are the usual four-sided rectangles. In $\mathbb{R}^3$ they are the closed parallelepipeds.

We say that the lengths of the sides of the rectangle $R$ are $b_1-a_1,\ldots,b_d-a_d$. The volume of the rectangle $R$ is denoted by $|R|$, and is defined to be
\[
 |R|=(b_1-a_1)\cdots(b_d-a_d).
\]
Of course, when $d=1$ the ``volume'' equals length, and when $d=2$ it equals area.

\begin{figure}[ht]
\centering
\fbox{\parbox{0.75\textwidth}{\centering Figure from the original text omitted in this transcription.}}
\caption{Rectangles in $\mathbb{R}^d$, $d=1,2,3$.}
\end{figure}

An open rectangle is the product of open intervals, and the interior of the rectangle $R$ is then
\[
 (a_1,b_1)\times(a_2,b_2)\times\cdots\times(a_d,b_d).
\]
Also, a cube is a rectangle for which $b_1-a_1=b_2-a_2=\cdots=b_d-a_d$. So if $Q\subset\mathbb{R}^d$ is a cube of common side length $\ell$, then $|Q|=\ell^d$.

A union of rectangles is said to be almost disjoint if the interiors of the rectangles are disjoint.

In this chapter, coverings by rectangles and cubes play a major role, so we isolate here two important lemmas.

\begin{lemma}
If a rectangle is the almost disjoint union of finitely many other rectangles, say
\[
 R=\bigcup_{k=1}^N R_k,
\]
then
\[
 |R|=\sum_{k=1}^N |R_k|.
\]
\end{lemma}

\begin{proof}
We consider the grid formed by extending indefinitely the sides of all rectangles $R_1,\ldots,R_N$. This construction yields finitely many rectangles $\widetilde R_1,\ldots,\widetilde R_M$, and a partition $J_1,\ldots,J_N$ of the integers between $1$ and $M$, such that the unions
\[
 R=\bigcup_{j=1}^M \widetilde R_j
 \qquad\text{and}\qquad
 R_k=\bigcup_{j\in J_k}\widetilde R_j,\quad k=1,\ldots,N,
\]
are almost disjoint.

\begin{figure}[ht]
\centering
\fbox{\parbox{0.75\textwidth}{\centering Figure from the original text omitted in this transcription.}}
\caption{The grid formed by the rectangles $R_k$.}
\end{figure}

For the rectangle $R$, for example, we see that $|R|=\sum_{j=1}^M|\widetilde R_j|$, since the grid actually partitions the sides of $R$ and each $\widetilde R_j$ consists of taking products of the intervals in these partitions. Thus when adding the volumes of the $\widetilde R_j$ we are summing the corresponding products of lengths of the intervals that arise. Since this also holds for the other rectangles $R_1,\ldots,R_N$, we conclude that
\[
 |R|=\sum_{j=1}^M|\widetilde R_j|
 =\sum_{k=1}^N\sum_{j\in J_k}|\widetilde R_j|
 =\sum_{k=1}^N|R_k|.
\]
\end{proof}

A slight modification of this argument then yields the following.

\begin{lemma}
If $R,R_1,\ldots,R_N$ are rectangles, and
\[
 R\subset\bigcup_{k=1}^N R_k,
\]
then
\[
 |R|\le \sum_{k=1}^N |R_k|.
\]
\end{lemma}

The main idea consists of taking the grid formed by extending all sides of the rectangles $R,R_1,\ldots,R_N$, and noting that the sets corresponding to the $J_k$ (in the above proof) need not be disjoint any more.

We now proceed to give a description of the structure of open sets in terms of cubes. We begin with the case of $\mathbb{R}$.

\begin{theorem}
Every open subset $O$ of $\mathbb{R}$ can be written uniquely as a countable union of disjoint open intervals.
\end{theorem}

\begin{proof}
For each $x\in O$, let $I_x$ denote the largest open interval containing $x$ and contained in $O$. More precisely, since $O$ is open, $x$ is contained in some small (non-trivial) interval, and therefore if
\[
 a_x=\inf\{a<x:(a,x)\subset O\},
 \qquad
 b_x=\sup\{b>x:(x,b)\subset O\},
\]
we must have $a_x<x<b_x$ (with possibly infinite values for $a_x$ and $b_x$). If we now let $I_x=(a_x,b_x)$, then by construction we have $x\in I_x$ as well as $I_x\subset O$. Hence
\[
 O=\bigcup_{x\in O}I_x.
\]
Now suppose that two intervals $I_x$ and $I_y$ intersect. Then their union (which is also an open interval) is contained in $O$ and contains $x$. Since $I_x$ is maximal, we must have $(I_x\cup I_y)\subset I_x$, and similarly $(I_x\cup I_y)\subset I_y$. This can happen only if $I_x=I_y$; therefore, any two distinct intervals in the collection $\mathcal I=\{I_x\}_{x\in O}$ must be disjoint. The proof will be complete once we have shown that there are only countably many distinct intervals in the collection $\mathcal I$. This, however, is easy to see, since every open interval $I_x$ contains a rational number. Since different intervals are disjoint, they must contain distinct rationals, and therefore $\mathcal I$ is countable, as desired.
\end{proof}

Naturally, if $O$ is open and $O=\bigcup_{j=1}^{\infty}I_j$, where the $I_j$'s are disjoint open intervals, the measure of $O$ ought to be $\sum_{j=1}^{\infty}|I_j|$. Since this representation is unique, we could take this as a definition of measure; we would then note that whenever $O_1$ and $O_2$ are open and disjoint, the measure of their union is the sum of their measures. Although this provides a natural notion of measure for an open set, it is not immediately clear how to generalize it to other sets in $\mathbb{R}$. Moreover, a similar approach in higher dimensions already encounters complications even when defining measures of open sets, since in this context the direct analogue of Theorem 1.3 is not valid (see Exercise 12). There is, however, a substitute result.

\begin{theorem}
Every open subset $O$ of $\mathbb{R}^d$, $d\ge1$, can be written as a countable union of almost disjoint closed cubes.
\end{theorem}

\begin{proof}
We must construct a countable collection $\mathcal Q$ of closed cubes whose interiors are disjoint, and so that
\[
 O=\bigcup_{Q\in\mathcal Q}Q.
\]
As a first step, consider the grid in $\mathbb{R}^d$ formed by taking all closed cubes of side length $1$ whose vertices have integer coordinates. In other words, we consider the natural grid of lines parallel to the axes, that is, the grid generated by the lattice $\mathbb{Z}^d$. We shall also use the grids formed by cubes of side length $2^{-N}$ obtained by successively bisecting the original grid.

We either accept or reject cubes in the initial grid as part of $\mathcal Q$ according to the following rule: if $Q$ is entirely contained in $O$ then we accept $Q$; if $Q$ intersects both $O$ and $O^c$ then we tentatively accept it; and if $Q$ is entirely contained in $O^c$ then we reject it.

As a second step, we bisect the tentatively accepted cubes into $2^d$ cubes with side length $1/2$. We then repeat our procedure, by accepting the smaller cubes if they are completely contained in $O$, tentatively accepting them if they intersect both $O$ and $O^c$, and rejecting them if they are contained in $O^c$.

\begin{figure}[ht]
\centering
\fbox{\parbox{0.75\textwidth}{\centering Figure from the original text omitted in this transcription.}}
\caption{Decomposition of $O$ into almost disjoint cubes.}
\end{figure}

This procedure is then repeated indefinitely, and (by construction) the resulting collection $\mathcal Q$ of all accepted cubes is countable and consists of almost disjoint cubes. To see why their union is all of $O$, we note that given $x\in O$ there exists a cube of side length $2^{-N}$ (obtained from successive bisections of the original grid) that contains $x$ and that is entirely contained in $O$. Either this cube has been accepted, or it is contained in a cube that has been previously accepted. This shows that the union of all cubes in $\mathcal Q$ covers $O$.
\end{proof}

Once again, if $O=\bigcup_{j=1}^{\infty}R_j$ where the rectangles $R_j$ are almost disjoint, it is reasonable to assign to $O$ the measure $\sum_{j=1}^{\infty}|R_j|$. This is natural since the volume of the boundary of each rectangle should be $0$, and the overlap of the rectangles should not contribute to the volume of $O$. We note, however, that the above decomposition into cubes is not unique, and it is not immediate that the sum is independent of this decomposition. So in $\mathbb{R}^d$, with $d\ge2$, the notion of volume or area, even for open sets, is more subtle.

The general theory developed in the next section actually yields a notion of volume that is consistent with the decompositions of open sets of the previous two theorems, and applies to all dimensions. Before we come to that, we discuss an important example in $\mathbb{R}$.

\subsection*{The Cantor set}

The Cantor set plays a prominent role in set theory and in analysis in general. It and its variants provide a rich source of enlightening examples.

We begin with the closed unit interval $C_0=[0,1]$ and let $C_1$ denote the set obtained from deleting the middle third open interval from $[0,1]$, that is,
\[
 C_1=[0,1/3]\cup[2/3,1].
\]
Next, we repeat this procedure for each sub-interval of $C_1$; that is, we delete the middle third open interval. At the second stage we get
\[
 C_2=[0,1/9]\cup[2/9,1/3]\cup[2/3,7/9]\cup[8/9,1].
\]
We repeat this process for each sub-interval of $C_2$, and so on. This procedure yields a sequence $C_k$, $k=0,1,2,\ldots$, of compact sets with
\[
 C_0\supset C_1\supset C_2\supset\cdots\supset C_k\supset C_{k+1}\supset\cdots.
\]

\begin{figure}[ht]
\centering
\fbox{\parbox{0.75\textwidth}{\centering Figure from the original text omitted in this transcription.}}
\caption{Construction of the Cantor set.}
\end{figure}

The Cantor set $C$ is by definition the intersection of all $C_k$'s:
\[
 C=\bigcap_{k=0}^{\infty}C_k.
\]
The set $C$ is not empty, since all end-points of the intervals in $C_k$ (all $k$) belong to $C$.

Despite its simple construction, the Cantor set enjoys many interesting topological and analytical properties. For instance, $C$ is closed and bounded, hence compact. Also, $C$ is totally disconnected: given any $x,y\in C$ there exists $z\notin C$ that lies between $x$ and $y$. Finally, $C$ is perfect: it has no isolated points (Exercise 1).

Next, we turn our attention to the question of determining the ``size'' of $C$. This is a delicate problem, one that may be approached from different angles depending on the notion of size we adopt. For instance, in terms of cardinality the Cantor set is rather large: it is not countable. Since it can be mapped to the interval $[0,1]$, the Cantor set has the cardinality of the continuum (Exercise 2).

However, from the point of view of ``length'' the size of $C$ is small. Roughly speaking, the Cantor set has length zero, and this follows from the following intuitive argument: the set $C$ is covered by sets $C_k$ whose lengths go to zero. Indeed, $C_k$ is a disjoint union of $2^k$ intervals of length $3^{-k}$, making the total length of $C_k$ equal to $(2/3)^k$. But $C\subset C_k$ for all $k$, and $(2/3)^k\to0$ as $k$ tends to infinity. We shall define a notion of measure and make this argument precise in the next section.

\section{The exterior measure}

The notion of exterior measure is the first of two important concepts needed to develop a theory of measure. We begin with the definition and basic properties of exterior measure. Loosely speaking, the exterior measure $m^*$ assigns to any subset of $\mathbb{R}^d$ a first notion of size; various examples show that this notion coincides with our earlier intuition. However, the exterior measure lacks the desirable property of additivity when taking the union of disjoint sets. We remedy this problem in the next section, where we discuss in detail the other key concept of measure theory, the notion of measurable sets.

The exterior measure, as the name indicates, attempts to describe the volume of a set $E$ by approximating it from the outside. The set $E$ is covered by cubes, and if the covering gets finer, with fewer cubes overlapping, the volume of $E$ should be close to the sum of the volumes of the cubes.

The precise definition is as follows: if $E$ is any subset of $\mathbb{R}^d$, the exterior measure\footnote{Some authors use the term outer measure instead of exterior measure.} of $E$ is
\begin{equation}
 m^*(E)=\inf\sum_{j=1}^{\infty}|Q_j|,
 \tag{1}
\end{equation}
where the infimum is taken over all countable coverings
\[
 E\subset\bigcup_{j=1}^{\infty}Q_j
\]
by closed cubes. The exterior measure is always non-negative but could be infinite, so that in general we have $0\le m^*(E)\le\infty$, and therefore takes values in the extended positive numbers.

We make some preliminary remarks about the definition of the exterior measure given by (1).
\begin{enumerate}
\item[(i)] It is important to note that it would not suffice to allow finite sums in the definition of $m^*(E)$. The quantity that would be obtained if one considered only coverings of $E$ by finite unions of cubes is in general larger than $m^*(E)$. (See Exercise 14.)
\item[(ii)] One can, however, replace the coverings by cubes, with coverings by rectangles; or with coverings by balls. That the former alternative yields the same exterior measure is quite direct. (See Exercise 15.) The equivalence with the latter is more subtle. (See Exercise 26 in Chapter 3.)
\end{enumerate}

We begin our investigation of this new notion by providing examples of sets whose exterior measures can be calculated, and we check that the latter matches our intuitive idea of volume (length in one dimension, area in two dimensions, etc.).

\paragraph{Example 1.} The exterior measure of a point is zero. This is clear once we observe that a point is a cube with volume zero, and which covers itself. Of course the exterior measure of the empty set is also zero.

\paragraph{Example 2.} The exterior measure of a closed cube is equal to its volume. Indeed, suppose $Q$ is a closed cube in $\mathbb{R}^d$. Since $Q$ covers itself, we must have $m^*(Q)\le |Q|$. Therefore, it suffices to prove the reverse inequality. We consider an arbitrary covering $Q\subset\bigcup_{j=1}^{\infty}Q_j$ by cubes, and note that it suffices to prove that
\begin{equation}
 |Q|\le\sum_{j=1}^{\infty}|Q_j|.
 \tag{2}
\end{equation}
For a fixed $\epsilon>0$ we choose for each $j$ an open cube $S_j$ which contains $Q_j$, and such that $|S_j|\le(1+\epsilon)|Q_j|$. From the open covering $\bigcup_{j=1}^{\infty}S_j$ of the compact set $Q$, we may select a finite subcovering which, after possibly renumbering the rectangles, we may write as $Q\subset\bigcup_{j=1}^N S_j$. Taking the closure of the cubes $S_j$, we may apply Lemma 1.2 to conclude that $|Q|\le\sum_{j=1}^N|S_j|$. Consequently,
\[
 |Q|\le(1+\epsilon)\sum_{j=1}^N|Q_j|
 \le(1+\epsilon)\sum_{j=1}^{\infty}|Q_j|.
\]
Since $\epsilon$ is arbitrary, we find that the inequality (2) holds; thus $|Q|\le m^*(Q)$, as desired.

\paragraph{Example 3.} If $Q$ is an open cube, the result $m^*(Q)=|Q|$ still holds. Since $Q$ is covered by its closure $\overline Q$, and $|\overline Q|=|Q|$, we immediately see that $m^*(Q)\le |Q|$. To prove the reverse inequality, we note that if $Q_0$ is a closed cube contained in $Q$, then $m^*(Q_0)\le m^*(Q)$, since any covering of $Q$ by a countable number of closed cubes is also a covering of $Q_0$ (see Observation 1 below). Hence $|Q_0|\le m^*(Q)$, and since we can choose $Q_0$ with a volume as close as we wish to $|Q|$, we must have $|Q|\le m^*(Q)$.

\paragraph{Example 4.} The exterior measure of a rectangle $R$ is equal to its volume. Indeed, arguing as in Example 2, we see that $|R|\le m^*(R)$. To obtain the reverse inequality, consider a grid in $\mathbb{R}^d$ formed by cubes of side length $1/k$. Then, if $\mathcal Q$ consists of the (finite) collection of all cubes entirely contained in $R$, and $\mathcal Q'$ the (finite) collection of all cubes that intersect the complement of $R$, we first note that
\[
 R\subset\bigcup_{Q\in(\mathcal Q\cup\mathcal Q')}Q.
\]
Also, a simple argument yields
\[
 \sum_{Q\in\mathcal Q}|Q|\le |R|.
\]
Moreover, there are $O(k^{d-1})$ cubes\footnote{We remind the reader of the notation $f(x)=O(g(x))$, which means that $|f(x)|\le C|g(x)|$ for some constant $C$ and all $x$ in a given range. In this particular example, there are fewer than $Ck^{d-1}$ cubes in question, as $k\to\infty$.} in $\mathcal Q'$, and these cubes have volume $k^{-d}$, so that $\sum_{Q\in\mathcal Q'}|Q|=O(1/k)$. Hence
\[
 \sum_{Q\in(\mathcal Q\cup\mathcal Q')}|Q|\le |R|+O(1/k),
\]
and letting $k$ tend to infinity yields $m^*(R)\le |R|$, as desired.

\paragraph{Example 5.} The exterior measure of $\mathbb{R}^d$ is infinite. This follows from the fact that any covering of $\mathbb{R}^d$ is also a covering of any cube $Q\subset\mathbb{R}^d$, hence $|Q|\le m^*(\mathbb{R}^d)$. Since $Q$ can have arbitrarily large volume, we must have $m^*(\mathbb{R}^d)=\infty$.

\paragraph{Example 6.} The Cantor set $C$ has exterior measure $0$. From the construction of $C$, we know that $C\subset C_k$, where each $C_k$ is a disjoint union of $2^k$ closed intervals, each of length $3^{-k}$. Consequently, $m^*(C)\le(2/3)^k$ for all $k$, hence $m^*(C)=0$.

\subsection*{Properties of the exterior measure}

The previous examples and comments provide some intuition underlying the definition of exterior measure. Here, we turn to the further study of $m^*$ and prove five properties of exterior measure that are needed in what follows.

First, we record the following remark that is immediate from the definition of $m^*$:
\begin{itemize}
\item For every $\epsilon>0$, there exists a covering $E\subset\bigcup_{j=1}^{\infty}Q_j$ with
\[
 \sum_{j=1}^{\infty}m^*(Q_j)\le m^*(E)+\epsilon.
\]
\end{itemize}

The relevant properties of exterior measure are listed in a series of observations.

\paragraph{Observation 1 (Monotonicity).} If $E_1\subset E_2$, then $m^*(E_1)\le m^*(E_2)$. This follows once we observe that any covering of $E_2$ by a countable collection of cubes is also a covering of $E_1$.

In particular, monotonicity implies that every bounded subset of $\mathbb{R}^d$ has finite exterior measure.

\paragraph{Observation 2 (Countable sub-additivity).} If $E=\bigcup_{j=1}^{\infty}E_j$, then
\[
 m^*(E)\le\sum_{j=1}^{\infty}m^*(E_j).
\]
First, we may assume that each $m^*(E_j)<\infty$, for otherwise the inequality clearly holds. For any $\epsilon>0$, the definition of the exterior measure yields for each $j$ a covering
\[
 E_j\subset\bigcup_{k=1}^{\infty}Q_{k,j}
\]
by closed cubes with
\[
 \sum_{k=1}^{\infty}|Q_{k,j}|\le m^*(E_j)+\frac{\epsilon}{2^j}.
\]
Then $E\subset\bigcup_{j,k=1}^{\infty}Q_{k,j}$ is a covering of $E$ by closed cubes, and therefore
\[
\begin{aligned}
 m^*(E)
 &\le \sum_{j,k}|Q_{k,j}|
 =\sum_{j=1}^{\infty}\sum_{k=1}^{\infty}|Q_{k,j}|\\
 &\le\sum_{j=1}^{\infty}\left(m^*(E_j)+\frac{\epsilon}{2^j}\right)
 =\sum_{j=1}^{\infty}m^*(E_j)+\epsilon.
\end{aligned}
\]
Since this holds true for every $\epsilon>0$, the second observation is proved.

\paragraph{Observation 3.} If $E\subset\mathbb{R}^d$, then
\[
 m^*(E)=\inf m^*(O),
\]
where the infimum is taken over all open sets $O$ containing $E$.

By monotonicity, it is clear that the inequality $m^*(E)\le\inf m^*(O)$ holds. For the reverse inequality, let $\epsilon>0$ and choose cubes $Q_j$ such that $E\subset\bigcup_{j=1}^{\infty}Q_j$, with
\[
 \sum_{j=1}^{\infty}|Q_j|\le m^*(E)+\frac{\epsilon}{2}.
\]
Let $Q'_j$ denote an open cube containing $Q_j$, and such that
\[
 |Q'_j|\le |Q_j|+\frac{\epsilon}{2^{j+1}}.
\]
Then $O=\bigcup_{j=1}^{\infty}Q'_j$ is open, and by Observation 2,
\[
\begin{aligned}
 m^*(O)
 &\le\sum_{j=1}^{\infty}m^*(Q'_j)
 =\sum_{j=1}^{\infty}|Q'_j|\\
 &\le\sum_{j=1}^{\infty}\left(|Q_j|+\frac{\epsilon}{2^{j+1}}\right)
 \le\sum_{j=1}^{\infty}|Q_j|+\frac{\epsilon}{2}
 \le m^*(E)+\epsilon.
\end{aligned}
\]
Hence $\inf m^*(O)\le m^*(E)$, as was to be shown.

\paragraph{Observation 4.} If $E=E_1\cup E_2$, and $d(E_1,E_2)>0$, then
\[
 m^*(E)=m^*(E_1)+m^*(E_2).
\]
By Observation 2, we already know that $m^*(E)\le m^*(E_1)+m^*(E_2)$, so it suffices to prove the reverse inequality. To this end, we first select $\delta$ such that $d(E_1,E_2)>\delta>0$. Next, we choose a covering $E\subset\bigcup_{j=1}^{\infty}Q_j$ by closed cubes, with
\[
 \sum_{j=1}^{\infty}|Q_j|\le m^*(E)+\epsilon.
\]
We may, after subdividing the cubes $Q_j$, assume that each $Q_j$ has a diameter less than $\delta$. In this case, each $Q_j$ can intersect at most one of the two sets $E_1$ or $E_2$. If we denote by $J_1$ and $J_2$ the sets of those indices $j$ for which $Q_j$ intersects $E_1$ and $E_2$, respectively, then $J_1\cap J_2$ is empty, and we have
\[
 E_1\subset\bigcup_{j\in J_1}Q_j,
 \qquad
 E_2\subset\bigcup_{j\in J_2}Q_j.
\]
Therefore,
\[
 m^*(E_1)+m^*(E_2)
 \le\sum_{j\in J_1}|Q_j|+\sum_{j\in J_2}|Q_j|
 \le\sum_{j=1}^{\infty}|Q_j|
 \le m^*(E)+\epsilon.
\]
Since $\epsilon$ is arbitrary, the proof of Observation 4 is complete.

\paragraph{Observation 5.} If a set $E$ is the countable union of almost disjoint cubes
\[
 E=\bigcup_{j=1}^{\infty}Q_j,
\]
then
\[
 m^*(E)=\sum_{j=1}^{\infty}|Q_j|.
\]
Let $\widetilde Q_j$ denote a cube strictly contained in $Q_j$ such that
\[
 |Q_j|\le |\widetilde Q_j|+\frac{\epsilon}{2^j},
\]
where $\epsilon$ is arbitrary but fixed. Then, for every $N$, the cubes $\widetilde Q_1,\widetilde Q_2,\ldots,\widetilde Q_N$ are disjoint, hence at a finite distance from one another, and repeated applications of Observation 4 imply
\[
 m^*\left(\bigcup_{j=1}^N\widetilde Q_j\right)
 =\sum_{j=1}^N|\widetilde Q_j|
 \ge\sum_{j=1}^N\left(|Q_j|-\frac{\epsilon}{2^j}\right).
\]
Since $\bigcup_{j=1}^N\widetilde Q_j\subset E$, we conclude that for every integer $N$,
\[
 m^*(E)\ge\sum_{j=1}^N|Q_j|-\epsilon.
\]
In the limit as $N$ tends to infinity we deduce $\sum_{j=1}^{\infty}|Q_j|\le m^*(E)+\epsilon$ for every $\epsilon>0$, hence $\sum_{j=1}^{\infty}|Q_j|\le m^*(E)$. Therefore, combined with Observation 2, our result proves that we have equality.

This last property shows that if a set can be decomposed into almost disjoint cubes, its exterior measure equals the sum of the volumes of the cubes. In particular, by Theorem 1.4 we see that the exterior measure of an open set equals the sum of the volumes of the cubes in a decomposition, and this coincides with our initial guess. Moreover, this also yields a proof that the sum is independent of the decomposition.

One can see from this that the volumes of simple sets that are calculated by elementary calculus agree with their exterior measure. This assertion can be proved most easily once we have developed the requisite tools in integration theory. (See Chapter 2.) In particular, we can then verify that the exterior measure of a ball (either open or closed) equals its volume.

Despite Observations 4 and 5, one cannot conclude in general that if $E_1\cup E_2$ is a disjoint union of subsets of $\mathbb{R}^d$, then
\begin{equation}
 m^*(E_1\cup E_2)=m^*(E_1)+m^*(E_2).
 \tag{3}
\end{equation}
In fact (3) holds when the sets in question are not highly irregular or ``pathological'' but are measurable in the sense described below.

\section{Measurable sets and the Lebesgue measure}

The notion of measurability isolates a collection of subsets in $\mathbb{R}^d$ for which the exterior measure satisfies all our desired properties, including additivity (and in fact countable additivity) for disjoint unions of sets.

There are a number of different ways of defining measurability, but these all turn out to be equivalent. Probably the simplest and most intuitive is the following: A subset $E$ of $\mathbb{R}^d$ is \emph{Lebesgue measurable}, or simply \emph{measurable}, if for any $\epsilon>0$ there exists an open set $O$ with $E\subset O$ and
\[
 m^*(O-E)\le\epsilon.
\]
This should be compared to Observation 3, which holds for all sets $E$.

If $E$ is measurable, we define its \emph{Lebesgue measure} (or measure) $m(E)$ by
\[
 m(E)=m^*(E).
\]
Clearly, the Lebesgue measure inherits all the features contained in Observations 1--5 of the exterior measure.

Immediately from the definition, we find:

\paragraph{Property 1.} Every open set in $\mathbb{R}^d$ is measurable.

Our immediate goal now is to gather various further properties of measurable sets. In particular, we shall prove that the collection of measurable sets behave well under the various operations of set theory: countable unions, countable intersections, and complements.

\paragraph{Property 2.} If $m^*(E)=0$, then $E$ is measurable. In particular, if $F$ is a subset of a set of exterior measure $0$, then $F$ is measurable.

By Observation 3 of the exterior measure, for every $\epsilon>0$ there exists an open set $O$ with $E\subset O$ and $m^*(O)\le\epsilon$. Since $(O-E)\subset O$, monotonicity implies $m^*(O-E)\le\epsilon$, as desired.

As a consequence of this property, we deduce that the Cantor set $C$ in Example 6 is measurable and has measure $0$.

\paragraph{Property 3.} A countable union of measurable sets is measurable.

Suppose $E=\bigcup_{j=1}^{\infty}E_j$, where each $E_j$ is measurable. Given $\epsilon>0$, we may choose for each $j$ an open set $O_j$ with $E_j\subset O_j$ and
\[
 m^*(O_j-E_j)\le\frac{\epsilon}{2^j}.
\]
Then the union $O=\bigcup_{j=1}^{\infty}O_j$ is open, $E\subset O$, and
\[
 (O-E)\subset\bigcup_{j=1}^{\infty}(O_j-E_j),
\]
so monotonicity and sub-additivity of the exterior measure imply
\[
 m^*(O-E)\le\sum_{j=1}^{\infty}m^*(O_j-E_j)\le\epsilon.
\]

\paragraph{Property 4.} Closed sets are measurable.

First, we observe that it suffices to prove that compact sets are measurable. Indeed, any closed set $F$ can be written as the union of compact sets, say
\[
 F=\bigcup_{k=1}^{\infty}(F\cap B_k),
\]
where $B_k$ denotes the closed ball of radius $k$ centered at the origin; then Property 3 applies.

So, suppose $F$ is compact (so that in particular $m^*(F)<\infty$), and let $\epsilon>0$. By Observation 3 we can select an open set $O$ with $F\subset O$ and $m^*(O)\le m^*(F)+\epsilon$. Since $F$ is closed, the difference $O-F$ is open, and by Theorem 1.4 we may write this difference as a countable union of almost disjoint cubes
\[
 O-F=\bigcup_{j=1}^{\infty}Q_j.
\]
For a fixed $N$, the finite union $K=\bigcup_{j=1}^NQ_j$ is compact; therefore $d(K,F)>0$ (we isolate this little fact in a lemma below). Since $(K\cup F)\subset O$, Observations 1, 4, and 5 of the exterior measure imply
\[
 m^*(O)\ge m^*(F)+m^*(K)
 =m^*(F)+\sum_{j=1}^N m^*(Q_j).
\]
Hence
\[
 \sum_{j=1}^N m^*(Q_j)\le m^*(O)-m^*(F)\le\epsilon,
\]
and this also holds in the limit as $N$ tends to infinity. Invoking the sub-additivity property of the exterior measure finally yields
\[
 m^*(O-F)\le\sum_{j=1}^{\infty}m^*(Q_j)\le\epsilon,
\]
as desired.

We digress briefly to complete the above argument by proving the following.

\begin{lemma}
If $F$ is closed, $K$ is compact, and these sets are disjoint, then $d(F,K)>0$.
\end{lemma}

\begin{proof}
Since $F$ is closed, for each point $x\in K$, there exists $\delta_x>0$ so that $d(x,F)>3\delta_x$. Since $\bigcup_{x\in K}B_{2\delta_x}(x)$ covers $K$, and $K$ is compact, we may find a subcover, which we denote by $\bigcup_{j=1}^N B_{2\delta_j}(x_j)$. If we let $\delta=\min(\delta_1,\ldots,\delta_N)$, then we must have $d(K,F)\ge\delta>0$. Indeed, if $x\in K$ and $y\in F$, then for some $j$ we have $|x_j-x|\le2\delta_j$, and by construction $|y-x_j|\ge3\delta_j$. Therefore
\[
 |y-x|\ge |y-x_j|-|x_j-x|\ge3\delta_j-2\delta_j\ge\delta,
\]
and the lemma is proved.
\end{proof}

\paragraph{Property 5.} The complement of a measurable set is measurable.

If $E$ is measurable, then for every positive integer $n$ we may choose an open set $O_n$ with $E\subset O_n$ and $m^*(O_n-E)\le1/n$. The complement $O_n^c$ is closed, hence measurable, which implies that the union
\[
 S=\bigcup_{n=1}^{\infty}O_n^c
\]
is also measurable by Property 3. Now we simply note that $S\subset E^c$, and
\[
 (E^c-S)\subset(O_n-E),
\]
such that $m^*(E^c-S)\le1/n$ for all $n$. Therefore, $m^*(E^c-S)=0$, and $E^c-S$ is measurable by Property 2. Therefore $E^c$ is measurable since it is the union of two measurable sets, namely $S$ and $(E^c-S)$.

\paragraph{Property 6.} A countable intersection of measurable sets is measurable.

This follows from Properties 3 and 5, since
\[
 \bigcap_{j=1}^{\infty}E_j
 =\left(\bigcup_{j=1}^{\infty}E_j^c\right)^c.
\]

In conclusion, we find that the family of measurable sets is closed under the familiar operations of set theory. We point out that we have shown more than simply closure with respect to finite unions and intersections: we have proved that the collection of measurable sets is closed under countable unions and intersections. This passage from finite operations to infinite ones is crucial in the context of analysis. We emphasize, however, that the operations of uncountable unions or intersections are not permissible when dealing with measurable sets!

\begin{theorem}
If $E_1,E_2,\ldots$ are disjoint measurable sets, and $E=\bigcup_{j=1}^{\infty}E_j$, then
\[
 m(E)=\sum_{j=1}^{\infty}m(E_j).
\]
\end{theorem}

\begin{proof}
First, we assume further that each $E_j$ is bounded. Then, for each $j$, by applying the definition of measurability to $E_j^c$, we can choose a closed subset $F_j$ of $E_j$ with
\[
 m^*(E_j-F_j)\le\frac{\epsilon}{2^j}.
\]
For each fixed $N$, the sets $F_1,\ldots,F_N$ are compact and disjoint, so that
\[
 m\left(\bigcup_{j=1}^N F_j\right)=\sum_{j=1}^N m(F_j).
\]
Since $\bigcup_{j=1}^N F_j\subset E$, we must have
\[
 m(E)\ge\sum_{j=1}^N m(F_j)
 \ge\sum_{j=1}^N m(E_j)-\epsilon.
\]
Letting $N$ tend to infinity, since $\epsilon$ was arbitrary we find that
\[
 m(E)\ge\sum_{j=1}^{\infty}m(E_j).
\]
Since the reverse inequality always holds (by sub-additivity in Observation 2), this concludes the proof when each $E_j$ is bounded.

In the general case, we select any sequence of cubes $\{Q_k\}_{k=1}^{\infty}$ that increases to $\mathbb{R}^d$, in the sense that $Q_k\subset Q_{k+1}$ for all $k\ge1$ and $\bigcup_{k=1}^{\infty}Q_k=\mathbb{R}^d$. We then let $S_1=Q_1$ and $S_k=Q_k-Q_{k-1}$ for $k\ge2$. If we define measurable sets by $E_{j,k}=E_j\cap S_k$, then
\[
 E=\bigcup_{j,k}E_{j,k}.
\]
The union above is disjoint and every $E_{j,k}$ is bounded. Moreover $E_j=\bigcup_{k=1}^{\infty}E_{j,k}$, and this union is also disjoint. Putting these facts together, and using what has already been proved, we obtain
\[
 m(E)=\sum_{j,k}m(E_{j,k})
 =\sum_j\sum_k m(E_{j,k})
 =\sum_j m(E_j),
\]
as claimed.
\end{proof}

With this, the countable additivity of the Lebesgue measure on measurable sets has been established. This result provides the necessary connection between the following:
\begin{itemize}
\item our primitive notion of volume given by the exterior measure,
\item the more refined idea of measurable sets, and
\item the countably infinite operations allowed on these sets.
\end{itemize}

We make two definitions to state succinctly some further consequences. If $E_1,E_2,\ldots$ is a countable collection of subsets of $\mathbb{R}^d$ that increases to $E$ in the sense that $E_k\subset E_{k+1}$ for all $k$, and $E=\bigcup_{k=1}^{\infty}E_k$, then we write $E_k\nearrow E$.

Similarly, if $E_1,E_2,\ldots$ decreases to $E$ in the sense that $E_k\supset E_{k+1}$ for all $k$, and $E=\bigcap_{k=1}^{\infty}E_k$, we write $E_k\searrow E$.

\begin{corollary}
Suppose $E_1,E_2,\ldots$ are measurable subsets of $\mathbb{R}^d$.
\begin{enumerate}
\item[(i)] If $E_k\nearrow E$, then
\[
 m(E)=\lim_{N\to\infty}m(E_N).
\]
\item[(ii)] If $E_k\searrow E$ and $m(E_k)<\infty$ for some $k$, then
\[
 m(E)=\lim_{N\to\infty}m(E_N).
\]
\end{enumerate}
\end{corollary}

\begin{proof}
For the first part, let $G_1=E_1$, $G_2=E_2-E_1$, and in general $G_k=E_k-E_{k-1}$ for $k\ge2$. By their construction, the sets $G_k$ are measurable, disjoint, and $E=\bigcup_{k=1}^{\infty}G_k$. Hence
\[
 m(E)=\sum_{k=1}^{\infty}m(G_k)
 =\lim_{N\to\infty}\sum_{k=1}^N m(G_k)
 =\lim_{N\to\infty}m\left(\bigcup_{k=1}^N G_k\right),
\]
and since $\bigcup_{k=1}^N G_k=E_N$ we get the desired limit.

For the second part, we may clearly assume that $m(E_1)<\infty$. Let $G_k=E_k-E_{k+1}$ for each $k$, so that
\[
 E_1=E\cup\bigcup_{k=1}^{\infty}G_k
\]
is a disjoint union of measurable sets. As a result, we find that
\[
\begin{aligned}
 m(E_1)
 &=m(E)+\lim_{N\to\infty}\sum_{k=1}^{N-1}\bigl(m(E_k)-m(E_{k+1})\bigr)\\
 &=m(E)+m(E_1)-\lim_{N\to\infty}m(E_N).
\end{aligned}
\]
Hence, since $m(E_1)<\infty$, we see that $m(E)=\lim_{N\to\infty}m(E_N)$, and the proof is complete.
\end{proof}

The reader should note that the second conclusion may fail without the assumption that $m(E_k)<\infty$ for some $k$. This is shown by the simple example when $E_n=(n,\infty)\subset\mathbb{R}$, for all $n$.

What follows provides an important geometric and analytic insight into the nature of measurable sets, in terms of their relation to open and closed sets. Its thrust is that, in effect, an arbitrary measurable set can be well approximated by the open sets that contain it, and alternatively, by the closed sets it contains.

\begin{theorem}
Suppose $E$ is a measurable subset of $\mathbb{R}^d$. Then, for every $\epsilon>0$:
\begin{enumerate}
\item[(i)] There exists an open set $O$ with $E\subset O$ and $m(O-E)\le\epsilon$.
\item[(ii)] There exists a closed set $F$ with $F\subset E$ and $m(E-F)\le\epsilon$.
\item[(iii)] If $m(E)$ is finite, there exists a compact set $K$ with $K\subset E$ and $m(E-K)\le\epsilon$.
\item[(iv)] If $m(E)$ is finite, there exists a finite union $F=\bigcup_{j=1}^NQ_j$ of closed cubes such that
\[
 m(E\triangle F)\le\epsilon.
\]
\end{enumerate}
\end{theorem}

The notation $E\triangle F$ stands for the symmetric difference between the sets $E$ and $F$, defined by
\[
 E\triangle F=(E-F)\cup(F-E),
\]
which consists of those points that belong to only one of the two sets $E$ or $F$.

\begin{proof}
Part (i) is just the definition of measurability. For the second part, we know that $E^c$ is measurable, so there exists an open set $O$ with $E^c\subset O$ and $m(O-E^c)\le\epsilon$. If we let $F=O^c$, then $F$ is closed, $F\subset E$, and $E-F=O-E^c$. Hence $m(E-F)\le\epsilon$ as desired.

For (iii), we first pick a closed set $F$ so that $F\subset E$ and $m(E-F)\le\epsilon/2$. For each $n$, we let $B_n$ denote the ball centered at the origin of radius $n$, and define compact sets $K_n=F\cap B_n$. Then $E-K_n$ is a sequence of measurable sets that decreases to $E-F$, and since $m(E)<\infty$, we conclude that for all large $n$ one has $m(E-K_n)\le\epsilon$.

For the last part, choose a family of closed cubes $\{Q_j\}_{j=1}^{\infty}$ so that
\[
 E\subset\bigcup_{j=1}^{\infty}Q_j
 \qquad\text{and}\qquad
 \sum_{j=1}^{\infty}|Q_j|\le m(E)+\frac{\epsilon}{2}.
\]
Since $m(E)<\infty$, the series converges and there exists $N>0$ such that $\sum_{j=N+1}^{\infty}|Q_j|<\epsilon/2$. If $F=\bigcup_{j=1}^NQ_j$, then
\[
\begin{aligned}
 m(E\triangle F)
 &=m(E-F)+m(F-E)\\
 &\le m\left(\bigcup_{j=N+1}^{\infty}Q_j\right)
 +m\left(\bigcup_{j=1}^{\infty}Q_j-E\right)\\
 &\le\sum_{j=N+1}^{\infty}|Q_j|
 +\sum_{j=1}^{\infty}|Q_j|-m(E)\\
 &\le\epsilon.
\end{aligned}
\]
\end{proof}

\subsection*{Invariance properties of Lebesgue measure}

A crucial property of Lebesgue measure in $\mathbb{R}^d$ is its translation-invariance, which can be stated as follows: if $E$ is a measurable set and $h\in\mathbb{R}^d$, then the set $E_h=E+h=\{x+h:x\in E\}$ is also measurable, and $m(E+h)=m(E)$. With the observation that this holds for the special case when $E$ is a cube, one passes to the exterior measure of arbitrary sets $E$, and sees from the definition of $m^*$ given in Section 2 that $m^*(E_h)=m^*(E)$. To prove the measurability of $E_h$ under the assumption that $E$ is measurable, we note that if $O$ is open, $O\supset E$, and $m^*(O-E)<\epsilon$, then $O_h$ is open, $O_h\supset E_h$, and $m^*(O_h-E_h)<\epsilon$.

In the same way one can prove the relative dilation-invariance of Lebesgue measure. Suppose $\delta>0$, and denote by $\delta E$ the set $\{\delta x:x\in E\}$. We can then assert that $\delta E$ is measurable whenever $E$ is, and
\[
 m(\delta E)=\delta^d m(E).
\]
One can also easily see that Lebesgue measure is reflection-invariant. That is, whenever $E$ is measurable, so is $-E=\{-x:x\in E\}$ and $m(-E)=m(E)$.

Other invariance properties of Lebesgue measure are in Exercises 7 and 8, and Problem 4 of Chapter 2.

\subsection*{$\sigma$-algebras and Borel sets}

A $\sigma$-algebra of sets is a collection of subsets of $\mathbb{R}^d$ that is closed under countable unions, countable intersections, and complements.

The collection of all subsets of $\mathbb{R}^d$ is of course a $\sigma$-algebra. A more interesting and relevant example consists of all measurable sets in $\mathbb{R}^d$, which we have just shown also forms a $\sigma$-algebra.

Another $\sigma$-algebra, which plays a vital role in analysis, is the Borel $\sigma$-algebra in $\mathbb{R}^d$, denoted by $\mathcal B_{\mathbb{R}^d}$, which by definition is the smallest $\sigma$-algebra that contains all open sets. Elements of this $\sigma$-algebra are called Borel sets.

The definition of the Borel $\sigma$-algebra will be meaningful once we have defined the term ``smallest,'' and shown that such a $\sigma$-algebra exists and is unique. The term ``smallest'' means that if $\mathcal S$ is any $\sigma$-algebra that contains all open sets in $\mathbb{R}^d$, then necessarily $\mathcal B_{\mathbb{R}^d}\subset\mathcal S$. Since we observe that any intersection (not necessarily countable) of $\sigma$-algebras is again a $\sigma$-algebra, we may define $\mathcal B_{\mathbb{R}^d}$ as the intersection of all $\sigma$-algebras that contain the open sets. This shows the existence and uniqueness of the Borel $\sigma$-algebra.

Since open sets are measurable, we conclude that the Borel $\sigma$-algebra is contained in the $\sigma$-algebra of measurable sets. Naturally, we may ask if this inclusion is strict: do there exist Lebesgue measurable sets which are not Borel sets? The answer is ``yes.'' (See Exercise 35.)

From the point of view of the Borel sets, the Lebesgue sets arise as the completion of the $\sigma$-algebra of Borel sets, that is, by adjoining all subsets of Borel sets of measure zero. This is an immediate consequence of Corollary 3.5 below.

Starting with the open and closed sets, which are the simplest Borel sets, one could try to list the Borel sets in order of their complexity. Next in order would come countable intersections of open sets; such sets are called $G_\delta$ sets. Alternatively, one could consider their complements, the countable union of closed sets, called the $F_\sigma$ sets.\footnote{The terminology $G_\delta$ comes from German ``Gebiete'' and ``Durschnitt''; $F_\sigma$ comes from French ``ferm\'e'' and ``somme.''}

\begin{corollary}
A subset $E$ of $\mathbb{R}^d$ is measurable
\begin{enumerate}
\item[(i)] if and only if $E$ differs from a $G_\delta$ by a set of measure zero,
\item[(ii)] if and only if $E$ differs from an $F_\sigma$ by a set of measure zero.
\end{enumerate}
\end{corollary}

\begin{proof}
Clearly $E$ is measurable whenever it satisfies either (i) or (ii), since the $F_\sigma$, $G_\delta$, and sets of measure zero are measurable.

Conversely, if $E$ is measurable, then for each integer $n\ge1$ we may select an open set $O_n$ that contains $E$, and such that $m(O_n-E)\le1/n$. Then
\[
 S=\bigcap_{n=1}^{\infty}O_n
\]
is a $G_\delta$ that contains $E$, and $(S-E)\subset(O_n-E)$ for all $n$. Therefore $m(S-E)\le1/n$ for all $n$; hence $S-E$ has exterior measure zero, and is therefore measurable.

For the second implication, we simply apply part (ii) of Theorem 3.4 with $\epsilon=1/n$, and take the union of the resulting closed sets.
\end{proof}

\subsection*{Construction of a non-measurable set}

Are all subsets of $\mathbb{R}^d$ measurable? In this section, we answer this question when $d=1$ by constructing a subset of $\mathbb{R}$ which is not measurable.\footnote{The existence of such a set in $\mathbb{R}$ implies the existence of corresponding non-measurable subsets of $\mathbb{R}^d$ for each $d$, as a consequence of Proposition 3.4 in the next chapter.} This justifies the conclusion that a satisfactory theory of measure cannot encompass all subsets of $\mathbb{R}$.

The construction of a non-measurable set $N$ uses the axiom of choice, and rests on a simple equivalence relation among real numbers in $[0,1]$. We write $x\sim y$ whenever $x-y$ is rational, and note that this is an equivalence relation since the following properties hold:
\begin{itemize}
\item $x\sim x$ for every $x\in[0,1]$;
\item if $x\sim y$, then $y\sim x$;
\item if $x\sim y$ and $y\sim z$, then $x\sim z$.
\end{itemize}
Two equivalence classes either are disjoint or coincide, and $[0,1]$ is the disjoint union of all equivalence classes, which we write as
\[
 [0,1]=\bigcup_\alpha E_\alpha.
\]
Now we construct the set $N$ by choosing exactly one element $x_\alpha$ from each $E_\alpha$, and setting $N=\{x_\alpha\}$. This (seemingly obvious) step requires further comment, which we postpone until after the proof of the following theorem.

\begin{theorem}
The set $N$ is not measurable.
\end{theorem}

\begin{proof}
The proof is by contradiction, so we assume that $N$ is measurable. Let $\{r_k\}_{k=1}^{\infty}$ be an enumeration of all the rationals in $[-1,1]$, and consider the translates
\[
 N_k=N+r_k.
\]
We claim that the sets $N_k$ are disjoint, and
\begin{equation}
 [0,1]\subset\bigcup_{k=1}^{\infty}N_k\subset[-1,2].
 \tag{4}
\end{equation}
To see why these sets are disjoint, suppose that the intersection $N_k\cap N_{k'}$ is non-empty. Then there exist rationals $r_k\ne r_{k'}$ and $\alpha$ and $\beta$ with $x_\alpha+r_k=x_\beta+r_{k'}$; hence
\[
 x_\alpha-x_\beta=r_{k'}-r_k.
\]
Consequently $\alpha\ne\beta$ and $x_\alpha-x_\beta$ is rational; hence $x_\alpha\sim x_\beta$, which contradicts the fact that $N$ contains only one representative of each equivalence class.

The second inclusion is straightforward since each $N_k$ is contained in $[-1,2]$ by construction. Finally, if $x\in[0,1]$, then $x\sim x_\alpha$ for some $\alpha$, and therefore $x-x_\alpha=r_k$ for some $k$. Hence $x\in N_k$, and the first inclusion holds.

Now we may conclude the proof of the theorem. If $N$ were measurable, then so would be $N_k$ for all $k$, and since the union $\bigcup_{k=1}^{\infty}N_k$ is disjoint, the inclusions in (4) yield
\[
 1\le\sum_{k=1}^{\infty}m(N_k)\le3.
\]
Since $N_k$ is a translate of $N$, we must have $m(N_k)=m(N)$ for all $k$. Consequently,
\[
 1\le\sum_{k=1}^{\infty}m(N)\le3.
\]
This is the desired contradiction, since neither $m(N)=0$ nor $m(N)>0$ is possible.
\end{proof}

\subsection*{Axiom of choice}

That the construction of the set $N$ is possible is based on the following general proposition.
\begin{itemize}
\item Suppose $E$ is a set and $\{E_\alpha\}$ is a collection of non-empty subsets of $E$. (The indexing set of $\alpha$'s is not assumed to be countable.) Then there is a function $\alpha\mapsto x_\alpha$ (a ``choice function'') such that $x_\alpha\in E_\alpha$, for all $\alpha$.
\end{itemize}
In this general form this assertion is known as the axiom of choice. This axiom occurs (at least implicitly) in many proofs in mathematics, but because of its seeming intuitive self-evidence, its significance was not at first understood. The initial realization of the importance of this axiom was in its use to prove a famous assertion of Cantor, the well-ordering principle. This proposition (sometimes referred to as ``transfinite induction'') can be formulated as follows.

A set $E$ is linearly ordered if there is a binary relation $\le$ such that:
\begin{enumerate}
\item[(a)] $x\le x$ for all $x\in E$.
\item[(b)] If $x,y\in E$ are distinct, then either $x\le y$ or $y\le x$ (but not both).
\item[(c)] If $x\le y$ and $y\le z$, then $x\le z$.
\end{enumerate}
We say that a set $E$ can be well-ordered if it can be linearly ordered in such a way that every non-empty subset $A\subset E$ has a smallest element in that ordering (that is, an element $x_0\in A$ such that $x_0\le x$ for any other $x\in A$).

A simple example of a well-ordered set is $\mathbb{Z}^+$, the positive integers with their usual ordering. The fact that $\mathbb{Z}^+$ is well-ordered is an essential part of the usual (finite) induction principle. More generally, the well-ordering principle states:
\begin{itemize}
\item Any set $E$ can be well-ordered.
\end{itemize}
It is in fact nearly obvious that the well-ordering principle implies the axiom of choice: if we well-order $E$, we can choose $x_\alpha$ to be the smallest element in $E_\alpha$, and in this way we have constructed the required choice function. It is also true, but not as easy to show, that the converse implication holds, namely that the axiom of choice implies the well-ordering principle. (See Problem 6 for another equivalent formulation of the Axiom of Choice.)

We shall follow the common practice of assuming the axiom of choice (and hence the validity of the well-ordering principle).\footnote{It can be proved that in an appropriate formulation of the axioms of set theory, the axiom of choice is independent of the other axioms; thus we are free to accept its validity.} However, we should point out that while the axiom of choice seems self-evident the well-ordering principle leads quickly to some baffling conclusions: one only needs to spend a little time trying to imagine what a well-ordering of the reals might look like!

\section{Measurable functions}

With the notion of measurable sets in hand, we now turn our attention to the objects that lie at the heart of integration theory: measurable functions.

The starting point is the notion of a characteristic function of a set $E$, which is defined by
\[
 \chi_E(x)=
 \begin{cases}
 1,&\text{if }x\in E,\\
 0,&\text{if }x\notin E.
 \end{cases}
\]
The next step is to pass to the functions that are the building blocks of integration theory. For the Riemann integral it is in effect the class of step functions, with each given as a finite sum
\begin{equation}
 f=\sum_{k=1}^N a_k\chi_{R_k},
 \tag{5}
\end{equation}
where each $R_k$ is a rectangle, and the $a_k$ are constants.

However, for the Lebesgue integral we need a more general notion, as we shall see in the next chapter. A simple function is a finite sum
\begin{equation}
 f=\sum_{k=1}^N a_k\chi_{E_k},
 \tag{6}
\end{equation}
where each $E_k$ is a measurable set of finite measure, and the $a_k$ are constants.

\subsection{Definition and basic properties}

We begin by considering only real-valued functions $f$ on $\mathbb{R}^d$, which we allow to take on the infinite values $+\infty$ and $-\infty$, so that $f(x)$ belongs to the extended real numbers
\[
 -\infty\le f(x)\le\infty.
\]
We shall say that $f$ is finite-valued if $-\infty<f(x)<\infty$ for all $x$. In the theory that follows, and the many applications of it, we shall almost always find ourselves in situations where a function takes on infinite values on at most a set of measure zero.

A function $f$ defined on a measurable subset $E$ of $\mathbb{R}^d$ is measurable, if for all $a\in\mathbb{R}$, the set
\[
 f^{-1}([ -\infty,a))=\{x\in E:f(x)<a\}
\]
is measurable. To simplify our notation, we shall often denote the set $\{x\in E:f(x)<a\}$ simply by $\{f<a\}$ whenever no confusion is possible.

First, we note that there are many equivalent definitions of measurable functions. For example, we may require instead that the inverse image of closed intervals be measurable. Indeed, to prove that $f$ is measurable if and only if
\[
 \{x:f(x)\le a\}=\{f\le a\}
\]
is measurable for every $a$, we note that in one direction, one has
\[
 \{f\le a\}=\bigcap_{k=1}^{\infty}\{f<a+1/k\},
\]
and recall that the countable intersection of measurable sets is measurable. For the other direction, we observe that
\[
 \{f<a\}=\bigcup_{k=1}^{\infty}\{f\le a-1/k\}.
\]
Similarly, $f$ is measurable if and only if $\{f\ge a\}$ (or $\{f>a\}$) is measurable for every $a$. In the first case this is immediate from our definition and the fact that $\{f\ge a\}$ is the complement of $\{f<a\}$, and in the second case this follows from what we have just proved and the fact that $\{f\le a\}=\{f>a\}^c$. A simple consequence is that $-f$ is measurable whenever $f$ is measurable.

In the same way, one can show that if $f$ is finite-valued, then it is measurable if and only if the sets $\{a<f<b\}$ are measurable for every $a,b\in\mathbb{R}$. Similar conclusions hold for whichever combination of strict or weak inequalities one chooses. For example, if $f$ is finite-valued, then it is measurable if and only if $\{a\le f<b\}$ is measurable for all $a,b\in\mathbb{R}$. By the same arguments one sees the following.

\paragraph{Property 1.} The finite-valued function $f$ is measurable if and only if $f^{-1}(O)$ is measurable for every open set $O$, and if and only if $f^{-1}(F)$ is measurable for every closed set $F$.

Note that this property also applies to extended-valued functions, if we make the additional hypothesis that both $f^{-1}(\infty)$ and $f^{-1}(-\infty)$ are measurable sets.

\paragraph{Property 2.} If $f$ is continuous on $\mathbb{R}^d$, then $f$ is measurable. If $f$ is measurable and finite-valued, and $\Phi$ is continuous, then $\Phi\circ f$ is measurable.

In fact, $\Phi$ is continuous, so $\Phi^{-1}(( -\infty,a))$ is an open set $O$, and hence
\[
 (\Phi\circ f)^{-1}(( -\infty,a))=f^{-1}(O)
\]
is measurable.

It should be noted, however, that in general it is not true that $f\circ\Phi$ is measurable whenever $f$ is measurable and $\Phi$ is continuous. See Exercise 35.

\paragraph{Property 3.} Suppose $\{f_n\}_{n=1}^{\infty}$ is a sequence of measurable functions. Then
\[
 \sup_n f_n(x),\qquad \inf_n f_n(x),\qquad
 \limsup_{n\to\infty}f_n(x),\qquad
 \liminf_{n\to\infty}f_n(x)
\]
are measurable.

Proving that $\sup_n f_n$ is measurable requires noting that
\[
 \left\{\sup_n f_n>a\right\}=\bigcup_n\{f_n>a\}.
\]
This also yields the result for $\inf_n f_n(x)$, since this quantity equals $-\sup_n(-f_n(x))$.

The result for the limsup and liminf also follows from the two observations
\[
 \limsup_{n\to\infty}f_n(x)
 =\inf_k\left\{\sup_{n\ge k}f_n(x)\right\}
\]
and
\[
 \liminf_{n\to\infty}f_n(x)
 =\sup_k\left\{\inf_{n\ge k}f_n(x)\right\}.
\]

\paragraph{Property 4.} If $\{f_n\}_{n=1}^{\infty}$ is a collection of measurable functions, and
\[
 \lim_{n\to\infty}f_n(x)=f(x),
\]
then $f$ is measurable.

Since $f(x)=\limsup_{n\to\infty}f_n(x)=\liminf_{n\to\infty}f_n(x)$, this property is a consequence of Property 3.

\paragraph{Property 5.} If $f$ and $g$ are measurable, then
\begin{enumerate}
\item[(i)] the integer powers $f^k$, $k\ge1$, are measurable;
\item[(ii)] $f+g$ and $fg$ are measurable if both $f$ and $g$ are finite-valued.
\end{enumerate}
For (i) we simply note that if $k$ is odd, then
\[
 \{f^k>a\}=\{f>a^{1/k}\},
\]
and if $k$ is even and $a\ge0$, then
\[
 \{f^k>a\}=\{f>a^{1/k}\}\cup\{f<-a^{1/k}\}.
\]
For (ii), we first see that $f+g$ is measurable because
\[
 \{f+g>a\}
 =\bigcup_{r\in\mathbb{Q}}\{f>a-r\}\cap\{g>r\},
\]
with $\mathbb{Q}$ denoting the rationals.

Finally, $fg$ is measurable because of the previous results and the fact that
\[
 fg=\frac14\bigl[(f+g)^2-(f-g)^2\bigr].
\]

We shall say that two functions $f$ and $g$ defined on a set $E$ are equal almost everywhere, and write
\[
 f(x)=g(x)\quad\text{a.e. }x\in E,
\]
if the set $\{x\in E:f(x)\ne g(x)\}$ has measure zero. We sometimes abbreviate this by saying that $f=g$ a.e. More generally, a property or statement is said to hold almost everywhere (a.e.) if it is true except on a set of measure zero.

One sees easily that if $f$ is measurable and $f=g$ a.e., then $g$ is measurable. This follows at once from the fact that $\{f<a\}$ and $\{g<a\}$ differ by a set of measure zero. Moreover, all the properties above can be relaxed to conditions holding almost everywhere. For instance, if $\{f_n\}_{n=1}^{\infty}$ is a collection of measurable functions, and
\[
 \lim_{n\to\infty}f_n(x)=f(x)\quad\text{a.e.},
\]
then $f$ is measurable.

Note that if $f$ and $g$ are defined almost everywhere on a measurable subset $E\subset\mathbb{R}^d$, then the functions $f+g$ and $fg$ can only be defined on the intersection of the domains of $f$ and $g$. Since the union of two sets of measure zero has again measure zero, $f+g$ is defined almost everywhere on $E$. We summarize this discussion as follows.

\paragraph{Property 6.} Suppose $f$ is measurable, and $f(x)=g(x)$ for a.e. $x$. Then $g$ is measurable.

In this light, Property 5(ii) also holds when $f$ and $g$ are finite-valued almost everywhere.

\subsection{Approximation by simple functions or step functions}

The theorems in this section are all of the same nature and provide further insight in the structure of measurable functions. We begin by approximating pointwise, non-negative measurable functions by simple functions.

\begin{theorem}
Suppose $f$ is a non-negative measurable function on $\mathbb{R}^d$. Then there exists an increasing sequence of non-negative simple functions $\{\varphi_k\}_{k=1}^{\infty}$ that converges pointwise to $f$, namely,
\[
 \varphi_k(x)\le\varphi_{k+1}(x)
 \qquad\text{and}\qquad
 \lim_{k\to\infty}\varphi_k(x)=f(x),
\]
for all $x$.
\end{theorem}

\begin{proof}
We begin first with a truncation. For $N\ge1$, let $Q_N$ denote the cube centered at the origin and of side length $N$. Then we define
\[
 F_N(x)=
 \begin{cases}
 f(x),&\text{if }x\in Q_N\text{ and }f(x)\le N,\\
 N,&\text{if }x\in Q_N\text{ and }f(x)>N,\\
 0,&\text{otherwise}.
 \end{cases}
\]
Then, $F_N(x)\to f(x)$ as $N$ tends to infinity for all $x$. Now, we partition the range of $F_N$, namely $[0,N]$, as follows. For fixed $N,M\ge1$, we define
\[
 E_{\ell,M}
 =\left\{x\in Q_N:\frac{\ell}{M}<F_N(x)\le\frac{\ell+1}{M}\right\},
 \qquad 0\le\ell<NM.
\]
Then we may form
\[
 F_{N,M}(x)=\sum_\ell\frac{\ell}{M}\chi_{E_{\ell,M}}(x).
\]
Each $F_{N,M}$ is a simple function that satisfies
\[
 0\le F_N(x)-F_{N,M}(x)\le\frac1M
\]
for all $x$. If we now choose $N=M=2^k$ with $k\ge1$ integral, and let $\varphi_k=F_{2^k,2^k}$, then we see that
\[
 0\le F_{2^k}(x)-\varphi_k(x)\le\frac1{2^k}
\]
for all $x$, $\{\varphi_k\}$ is increasing, and this sequence satisfies all the desired properties.
\end{proof}

Note that the result holds for non-negative functions that are extended-valued, if the limit $+\infty$ is allowed. We now drop the assumption that $f$ is non-negative, and also allow the extended limit $-\infty$.

\begin{theorem}
Suppose $f$ is measurable on $\mathbb{R}^d$. Then there exists a sequence of simple functions $\{\varphi_k\}_{k=1}^{\infty}$ that satisfies
\[
 |\varphi_k(x)|\le |\varphi_{k+1}(x)|
 \qquad\text{and}\qquad
 \lim_{k\to\infty}\varphi_k(x)=f(x),
\]
for all $x$. In particular, we have $|\varphi_k(x)|\le|f(x)|$ for all $x$ and $k$.
\end{theorem}

\begin{proof}
We use the following decomposition of the function $f$:
\[
 f(x)=f^+(x)-f^-(x),
\]
where
\[
 f^+(x)=\max(f(x),0)
 \qquad\text{and}\qquad
 f^-(x)=\max(-f(x),0).
\]
Since both $f^+$ and $f^-$ are non-negative, the previous theorem yields two increasing sequences of non-negative simple functions $\{\varphi_k^{(1)}(x)\}_{k=1}^{\infty}$ and $\{\varphi_k^{(2)}(x)\}_{k=1}^{\infty}$ which converge pointwise to $f^+$ and $f^-$, respectively. Then, if we let
\[
 \varphi_k(x)=\varphi_k^{(1)}(x)-\varphi_k^{(2)}(x),
\]
we see that $\varphi_k(x)$ converges to $f(x)$ for all $x$. Finally, the sequence $\{|\varphi_k|\}$ is increasing because the definition of $f^+$, $f^-$ and the properties of $\varphi_k^{(1)}$ and $\varphi_k^{(2)}$ imply that
\[
 |\varphi_k(x)|=\varphi_k^{(1)}(x)+\varphi_k^{(2)}(x).
\]
\end{proof}

We may now go one step further, and approximate by step functions. Here, in general, the convergence may hold only almost everywhere.

\begin{theorem}
Suppose $f$ is measurable on $\mathbb{R}^d$. Then there exists a sequence of step functions $\{\psi_k\}_{k=1}^{\infty}$ that converges pointwise to $f(x)$ for almost every $x$.
\end{theorem}

\begin{proof}
By the previous result, it suffices to show that if $E$ is a measurable set with finite measure, then $f=\chi_E$ can be approximated by step functions. To this end, we recall part (iv) of Theorem 3.4, which states that for every $\epsilon$ there exist cubes $Q_1,\ldots,Q_N$ such that
\[
 m\left(E\triangle\bigcup_{j=1}^NQ_j\right)\le\epsilon.
\]
By considering the grid formed by extending the sides of these cubes, we see that there exist almost disjoint rectangles $\widetilde R_1,\ldots,\widetilde R_M$ such that
\[
 \bigcup_{j=1}^NQ_j=\bigcup_{j=1}^M\widetilde R_j.
\]
By taking rectangles $R_j$ contained in $\widetilde R_j$, and slightly smaller in size, we find a collection of disjoint rectangles that satisfy
\[
 m\left(E\triangle\bigcup_{j=1}^M R_j\right)\le2\epsilon.
\]
Therefore
\[
 f(x)=\sum_{j=1}^M\chi_{R_j}(x),
\]
except possibly on a set of measure $\le2\epsilon$. Consequently, for every $k\ge1$, there exists a step function $\psi_k(x)$ such that if
\[
 E_k=\{x:f(x)\ne\psi_k(x)\},
\]
then $m(E_k)\le2^{-k}$. If we let
\[
 F_K=\bigcup_{j=K+1}^{\infty}E_j
 \qquad\text{and}\qquad
 F=\bigcap_{K=1}^{\infty}F_K,
\]
then $m(F)=0$ since $m(F_K)\le2^{-K}$, and $\psi_k(x)\to f(x)$ for all $x$ in the complement of $F$, which is the desired result.
\end{proof}

\subsection{Littlewood's three principles}

Although the notions of measurable sets and measurable functions represent new tools, we should not overlook their relation to the older concepts they replaced. Littlewood aptly summarized these connections in the form of three principles that provide a useful intuitive guide in the initial study of the theory.
\begin{enumerate}
\item[(i)] Every set is nearly a finite union of intervals.
\item[(ii)] Every function is nearly continuous.
\item[(iii)] Every convergent sequence is nearly uniformly convergent.
\end{enumerate}
The sets and functions referred to above are of course assumed to be measurable. The catch is in the word ``nearly,'' which has to be understood appropriately in each context. A precise version of the first principle appears in part (iv) of Theorem 3.4. An exact formulation of the third principle is given in the following important result.

\begin{theorem}[Egorov]
Suppose $\{f_k\}_{k=1}^{\infty}$ is a sequence of measurable functions defined on a measurable set $E$ with $m(E)<\infty$, and assume that $f_k\to f$ a.e. on $E$. Given $\epsilon>0$, we can find a closed set $A_\epsilon\subset E$ such that
\[
 m(E-A_\epsilon)\le\epsilon
\]
and $f_k\to f$ uniformly on $A_\epsilon$.
\end{theorem}

\begin{proof}
We may assume without loss of generality that $f_k(x)\to f(x)$ for every $x\in E$. For each pair of non-negative integers $n$ and $k$, let
\[
 E_k^n=\{x\in E:|f_j(x)-f(x)|<1/n,\text{ for all }j>k\}.
\]
Now fix $n$ and note that $E_k^n\subset E_{k+1}^n$, and $E_k^n\nearrow E$ as $k$ tends to infinity. By Corollary 3.3, we find that there exists $k_n$ such that
\[
 m(E-E_{k_n}^n)<\frac1{2^n}.
\]
By construction, we then have
\[
 |f_j(x)-f(x)|<\frac1n
 \qquad\text{whenever }j>k_n\text{ and }x\in E_{k_n}^n.
\]
We choose $N$ so that $\sum_{n=N}^{\infty}2^{-n}<\epsilon/2$, and let
\[
 \widetilde A_\epsilon=\bigcap_{n\ge N}E_{k_n}^n.
\]
We first observe that
\[
 m(E-\widetilde A_\epsilon)
 \le\sum_{n=N}^{\infty}m(E-E_{k_n}^n)
 <\frac\epsilon2.
\]
Next, if $\delta>0$, we choose $n\ge N$ such that $1/n<\delta$, and note that $x\in\widetilde A_\epsilon$ implies $x\in E_{k_n}^n$. We see therefore that $|f_j(x)-f(x)|<\delta$ whenever $j>k_n$. Hence $f_k$ converges uniformly to $f$ on $\widetilde A_\epsilon$.

Finally, using Theorem 3.4 choose a closed subset $A_\epsilon\subset\widetilde A_\epsilon$ with
\[
 m(\widetilde A_\epsilon-A_\epsilon)<\frac\epsilon2.
\]
As a result, we have $m(E-A_\epsilon)<\epsilon$ and the theorem is proved.
\end{proof}

The next theorem attests to the validity of the second of Littlewood's principles.

\begin{theorem}[Lusin]
Suppose $f$ is measurable and finite valued on $E$ with $E$ of finite measure. Then for every $\epsilon>0$ there exists a closed set $F_\epsilon$, with
\[
 F_\epsilon\subset E,
 \qquad
 m(E-F_\epsilon)\le\epsilon,
\]
and such that $f|_{F_\epsilon}$ is continuous.
\end{theorem}

By $f|_{F_\epsilon}$ we mean the restriction of $f$ to the set $F_\epsilon$. The conclusion of the theorem states that if $f$ is viewed as a function defined only on $F_\epsilon$, then $f$ is continuous. However, the theorem does not make the stronger assertion that the function $f$ defined on $E$ is continuous at the points of $F_\epsilon$.

\begin{proof}
Let $f_n$ be a sequence of step functions so that $f_n\to f$ a.e. Then we may find sets $E_n$ so that $m(E_n)<2^{-n}$ and $f_n$ is continuous outside $E_n$. By Egorov's theorem, we may find a set $A_{\epsilon/3}$ on which $f_n\to f$ uniformly and $m(E-A_{\epsilon/3})\le\epsilon/3$. Then we consider
\[
 F'=A_{\epsilon/3}-\bigcup_{n\ge N}E_n
\]
for $N$ so large that $\sum_{n\ge N}2^{-n}<\epsilon/3$. Now for every $n\ge N$ the function $f_n$ is continuous on $F'$; thus $f$ (being the uniform limit of $\{f_n\}$) is also continuous on $F'$. To finish the proof, we merely need to approximate the set $F'$ by a closed set $F_\epsilon\subset F'$ such that
\[
 m(F'-F_\epsilon)<\frac\epsilon3.
\]
\end{proof}

\section{* The Brunn--Minkowski inequality}

Since addition and multiplication by scalars are basic features of vector spaces, it is not surprising that properties of these operations arise in a fundamental way in the theory of Lebesgue measure on $\mathbb{R}^d$. We have already discussed in this connection the translation-invariance and relative dilation-invariance of Lebesgue measure. Here we come to the study of the sum of two measurable sets $A$ and $B$, defined as
\[
 A+B=\{x\in\mathbb{R}^d:x=x'+x''\text{ with }x'\in A\text{ and }x''\in B\}.
\]
This notion is of importance in a number of questions, in particular in the theory of convex sets; we shall apply it to the isoperimetric problem in Chapter 3.

In this regard the first (admittedly vague) question we can pose is whether one can establish any general estimate for the measure of $A+B$ in terms of the measures of $A$ and $B$ (assuming that these three sets are measurable). We can see easily that it is not possible to obtain an upper bound for $m(A+B)$ in terms of $m(A)$ and $m(B)$. Indeed, simple examples show that we may have $m(A)=m(B)=0$ while $m(A+B)>0$. (See Exercise 20.)

In the converse direction one might ask for a general estimate of the form
\[
 m(A+B)^\alpha\ge c_\alpha\bigl(m(A)^\alpha+m(B)^\alpha\bigr),
\]
where $\alpha$ is a positive number and the constant $c_\alpha$ is independent of $A$ and $B$. Clearly, the best one can hope for is $c_\alpha=1$. The role of the exponent $\alpha$ can be understood by considering convex sets. Such sets $A$ are defined by the property that whenever $x$ and $y$ are in $A$ then the line segment joining them,
\[
 \{xt+y(1-t):0\le t\le1\},
\]
also belongs to $A$. If we recall the definition $\lambda A=\{\lambda x:x\in A\}$ for $\lambda>0$, we note that whenever $A$ is convex, then $A+\lambda A=(1+\lambda)A$. However, $m((1+\lambda)A)=(1+\lambda)^dm(A)$, and thus the presumed inequality can hold only if
\[
 (1+\lambda)^{d\alpha}\ge1+\lambda^{d\alpha},
\]
for all $\lambda>0$. Now
\begin{equation}
 (a+b)^\gamma\ge a^\gamma+b^\gamma
 \qquad\text{if }\gamma\ge1\text{ and }a,b\ge0,
 \tag{7}
\end{equation}
while the reverse inequality holds if $0\le\gamma\le1$. (See Exercise 38.) This yields $\alpha\ge1/d$. Moreover, (7) shows that the inequality with the exponent $1/d$ implies the corresponding inequality with $\alpha\ge1/d$, and so we are naturally led to the inequality
\begin{equation}
 m(A+B)^{1/d}\ge m(A)^{1/d}+m(B)^{1/d}.
 \tag{8}
\end{equation}

Before proceeding with the proof of (8), we need to mention a technical impediment that arises. While we may assume that $A$ and $B$ are measurable, it does not necessarily follow that then $A+B$ is measurable. (See Exercise 13 in the next chapter.) However it is easily seen that this difficulty does not occur when, for example, $A$ and $B$ are closed sets, or when one of them is open. (See Exercise 19.)

With the above considerations in mind we can state the main result.

\begin{theorem}
Suppose $A$ and $B$ are measurable sets in $\mathbb{R}^d$ and their sum $A+B$ is also measurable. Then the inequality (8) holds.
\end{theorem}

\begin{proof}
Let us first check (8) when $A$ and $B$ are rectangles with side lengths $\{a_j\}_{j=1}^d$ and $\{b_j\}_{j=1}^d$, respectively. Then (8) becomes
\begin{equation}
 \left(\prod_{j=1}^d(a_j+b_j)\right)^{1/d}
 \ge
 \left(\prod_{j=1}^d a_j\right)^{1/d}
 +
 \left(\prod_{j=1}^d b_j\right)^{1/d},
 \tag{9}
\end{equation}
which by homogeneity we can reduce to the special case where $a_j+b_j=1$ for each $j$. In fact, notice that if we replace $a_j,b_j$ by $\lambda_j a_j,\lambda_j b_j$, with $\lambda_j>0$, then both sides of (9) are multiplied by $(\lambda_1\lambda_2\cdots\lambda_d)^{1/d}$. We then need only choose $\lambda_j=(a_j+b_j)^{-1}$. With this reduction, the inequality (9) is an immediate consequence of the arithmetic-geometric inequality (Exercise 39)
\[
 \frac1d\sum_{j=1}^d x_j
 \ge
 \left(\prod_{j=1}^d x_j\right)^{1/d},
 \qquad\text{for all }x_j\ge0:
\]
we add the two inequalities that result when we set $x_j=a_j$ and $x_j=b_j$, respectively.

We next turn to the case when each $A$ and $B$ are the union of finitely many rectangles whose interiors are disjoint. We shall prove (8) in this case by induction on the total number of rectangles in $A$ and $B$. We denote this number by $n$. Here it is important to note that the desired inequality is unchanged when we translate $A$ and $B$ independently. In fact, replacing $A$ by $A+h$ and $B$ by $B+h'$ replaces $A+B$ by $A+B+h+h'$, and thus the corresponding measures remain the same. We now choose a pair of disjoint rectangles $R_1$ and $R_2$ in the collection making up $A$, and we note that they can be separated by a coordinate hyperplane. Thus we may assume that for some $j$, after translation by an appropriate $h$, $R_1$ lies in
\[
 A^-=A\cap\{x_j\le0\},
\]
and $R_2$ in
\[
 A^+=A\cap\{0\le x_j\}.
\]
Observe also that both $A^+$ and $A^-$ contain at least one less rectangle than $A$ does, and $A=A^-\cup A^+$.

We next translate $B$ so that
\[
 B^-=B\cap\{x_j\le0\}
 \qquad\text{and}\qquad
 B^+=B\cap\{x_j\ge0\}
\]
satisfy
\[
 \frac{m(B^\pm)}{m(B)}=\frac{m(A^\pm)}{m(A)}.
\]
However,
\[
 A+B\supset(A^++B^+)\cup(A^-+B^-),
\]
and the union on the right is essentially disjoint, since the two parts lie in different half-spaces. Moreover, the total number of rectangles in either $A^+$ and $B^+$, or $A^-$ and $B^-$ is also less than $n$. Thus the induction hypothesis applies and
\[
\begin{aligned}
 m(A+B)
 &\ge m(A^++B^+)+m(A^-+B^-)\\
 &\ge\bigl(m(A^+)^{1/d}+m(B^+)^{1/d}\bigr)^d
   +\bigl(m(A^-)^{1/d}+m(B^-)^{1/d}\bigr)^d\\
 &=m(A^+)
 \left[1+\left(\frac{m(B)}{m(A)}\right)^{1/d}\right]^d
 +m(A^-)
 \left[1+\left(\frac{m(B)}{m(A)}\right)^{1/d}\right]^d\\
 &=\bigl(m(A)^{1/d}+m(B)^{1/d}\bigr)^d,
\end{aligned}
\]
which gives the desired inequality (8) when $A$ and $B$ are both finite unions of rectangles with disjoint interiors.

Next, this quickly implies the result when $A$ and $B$ are open sets of finite measure. Indeed, by Theorem 1.4, for any $\epsilon>0$ we can find unions of almost disjoint rectangles $A_\epsilon$ and $B_\epsilon$, such that
\[
 A_\epsilon\subset A,\qquad B_\epsilon\subset B,
\]
with
\[
 m(A)\le m(A_\epsilon)+\epsilon
 \qquad\text{and}\qquad
 m(B)\le m(B_\epsilon)+\epsilon.
\]
Since $A+B\supset A_\epsilon+B_\epsilon$, the inequality (8) for $A_\epsilon$ and $B_\epsilon$ and a passage to a limit gives the desired result.

From this, we can pass to the case where $A$ and $B$ are arbitrary compact sets, by noting first that $A+B$ is then compact, and that if we define
\[
 A_\epsilon=\{x:d(x,A)<\epsilon\},
\]
then $A_\epsilon$ are open, and $A_\epsilon\searrow A$ as $\epsilon\to0$. With similar definitions for $B_\epsilon$ and $(A+B)_\epsilon$, we observe also that
\[
 A+B\subset A_\epsilon+B_\epsilon\subset(A+B)_{2\epsilon}.
\]
Hence, letting $\epsilon\to0$, we see that (8) for $A_\epsilon$ and $B_\epsilon$ implies the desired result for $A$ and $B$. The general case, in which we assume that $A$, $B$, and $A+B$ are measurable, then follows by approximating $A$ and $B$ from inside by compact sets, as in (iii) of Theorem 3.4.
\end{proof}

% Exercises and Problems are intentionally omitted.

% Stein & Shakarchi, Real Analysis (Princeton Lectures in Analysis III)
% Chapter 2: Integration Theory
% Transcribed from the user-provided PDF.
% Exercises and Problems are intentionally omitted.
% This file is intended to be \input{} into a book/report document.
% Suggested packages in the parent document:
%   amsmath, amssymb, amsthm, mathtools
% Suggested theorem declarations (source-style statement numbering by section):
%   \newtheorem{theorem}{Theorem}[section]
%   \newtheorem{lemma}[theorem]{Lemma}
%   \newtheorem{corollary}[theorem]{Corollary}
%   \newtheorem{proposition}[theorem]{Proposition}
% To reproduce the source section labels 1,2,3,4 inside a book class, you may also use:
%   \renewcommand{\thesection}{\arabic{section}}
%
% Equation numbers (1)--(16) follow the source where useful.

\chapter{Integration Theory}

\begin{quote}
\ldots amongst the many definitions that have been successively proposed for the integral of real-valued functions of a real variable, I have retained only those which, in my opinion, are indispensable to understand the transformations undergone by the problem of integration, and to capture the relationship between the notion of area, so simple in appearance, and certain more complicated analytical definitions of the integral.

One might ask if there is sufficient interest to occupy oneself with such complications, and if it is not better to restrict oneself to the study of functions that necessitate only simple definitions\ldots. As we shall see in this course, we would then have to renounce the possibility of resolving many problems posed long ago, and which have simple statements. It is to solve these problems, and not for love of complications, that I have introduced in this book a definition of the integral more general than that of Riemann.

\hfill H. Lebesgue, 1903
\end{quote}

\section{The Lebesgue integral: basic properties and convergence theorems}

The general notion of the Lebesgue integral on $\mathbb{R}^d$ will be defined in a step-by-step fashion, proceeding successively to increasingly larger families of functions. At each stage we shall see that the integral satisfies elementary properties such as linearity and monotonicity, and we prove appropriate convergence theorems that amount to interchanging the integral with limits. At the end of the process we shall have achieved a general theory of integration that will be decisive in the study of further problems.

We proceed in four stages, by progressively integrating:
\begin{enumerate}
\item Simple functions.
\item Bounded functions supported on a set of finite measure.
\item Non-negative functions.
\item Integrable functions (the general case).
\end{enumerate}
We emphasize from the onset that all functions are assumed to be measurable. At the beginning we also consider only finite-valued functions which take on real values. Later we shall also consider extended-valued functions, and also complex-valued functions.

\subsection*{Stage one: simple functions}

Recall from the previous chapter that a simple function $\varphi$ is a finite sum
\begin{equation}
\varphi(x)=\sum_{k=1}^{N} a_k\chi_{E_k}(x),
\tag{1}
\end{equation}
where the $E_k$ are measurable sets of finite measure and the $a_k$ are constants. A complication that arises from this definition is that a simple function can be written in a multitude of ways as such finite linear combinations; for example, $0=\chi_E-\chi_E$ for any measurable set $E$ of finite measure. Fortunately, there is an unambiguous choice for the representation of a simple function, which is natural and useful in applications.

The \emph{canonical form} of $\varphi$ is the unique decomposition as in (1), where the numbers $a_k$ are distinct and non-zero, and the sets $E_k$ are disjoint.

Finding the canonical form of $\varphi$ is straightforward: since $\varphi$ can take only finitely many distinct and non-zero values, say $c_1,\ldots,c_M$, we may set $F_k=\{x:\varphi(x)=c_k\}$, and note that the sets $F_k$ are disjoint. Therefore
\[
\varphi=\sum_{k=1}^{M}c_k\chi_{F_k}
\]
is the desired canonical form of $\varphi$.

If $\varphi$ is a simple function with canonical form
\[
\varphi(x)=\sum_{k=1}^{M}c_k\chi_{F_k}(x),
\]
then we define the Lebesgue integral of $\varphi$ by
\[
\int_{\mathbb{R}^d}\varphi(x)\,dx
=\sum_{k=1}^{M}c_km(F_k).
\]
If $E$ is a measurable subset of $\mathbb{R}^d$ with finite measure, then $\varphi(x)\chi_E(x)$ is also a simple function, and we define
\[
\int_E\varphi(x)\,dx
=\int \varphi(x)\chi_E(x)\,dx.
\]
To emphasize the choice of the Lebesgue measure $m$ in the definition of the integral, one sometimes writes
\[
\int_{\mathbb{R}^d}\varphi(x)\,dm(x)
\]
for the Lebesgue integral of $\varphi$. In fact, as a matter of convenience, we shall often write $\int\varphi(x)\,dx$ or simply $\int\varphi$ for the integral of $\varphi$ over $\mathbb{R}^d$.

\begin{proposition}
The integral of simple functions defined above satisfies the following properties:
\begin{enumerate}
\item[(i)] \textbf{Independence of the representation.} If
\[
\varphi=\sum_{k=1}^{N}a_k\chi_{E_k}
\]
is any representation of $\varphi$, then
\[
\int\varphi=\sum_{k=1}^{N}a_km(E_k).
\]
\item[(ii)] \textbf{Linearity.} If $\varphi$ and $\psi$ are simple, and $a,b\in\mathbb{R}$, then
\[
\int(a\varphi+b\psi)=a\int\varphi+b\int\psi.
\]
\item[(iii)] \textbf{Additivity.} If $E$ and $F$ are disjoint subsets of $\mathbb{R}^d$ with finite measure, then
\[
\int_{E\cup F}\varphi=\int_E\varphi+\int_F\varphi.
\]
\item[(iv)] \textbf{Monotonicity.} If $\varphi\leq\psi$ are simple, then
\[
\int\varphi\leq\int\psi.
\]
\item[(v)] \textbf{Triangle inequality.} If $\varphi$ is a simple function, then so is $|\varphi|$, and
\[
\left|\int\varphi\right|\leq\int|\varphi|.
\]
\end{enumerate}
\end{proposition}

\begin{proof}
The only conclusion that is a little tricky is the first, which asserts that the integral of a simple function can be calculated by using any of its decompositions as a linear combination of characteristic functions.

Suppose that
\[
\varphi=\sum_{k=1}^{N}a_k\chi_{E_k},
\]
where we assume that the sets $E_k$ are disjoint, but we do not suppose that the numbers $a_k$ are distinct and non-zero. For each distinct non-zero value $a$ among the $\{a_k\}$ we define
\[
E'_a=\bigcup_{a_k=a}E_k.
\]
The sets $E'_a$ are disjoint, and
\[
m(E'_a)=\sum_{a_k=a}m(E_k).
\]
Then clearly $\varphi=\sum_a a\chi_{E'_a}$, where the sum is over the distinct non-zero values of $\{a_k\}$. Thus
\[
\int\varphi
=\sum_a am(E'_a)
=\sum_{k=1}^{N}a_km(E_k).
\]

Next, suppose $\varphi=\sum_{k=1}^{N}a_k\chi_{E_k}$, where we no longer assume that the $E_k$ are disjoint. Then we can ``refine'' the decomposition $\bigcup_{k=1}^{N}E_k$ by finding sets $E_1^*,E_2^*,\ldots,E_n^*$ with the property that
\[
\bigcup_{k=1}^{N}E_k=\bigcup_{j=1}^{n}E_j^*,
\]
the sets $E_j^*$ are mutually disjoint, and for each $k$, $E_k$ is the union of those $E_j^*$ that are contained in $E_k$. For each $j$, let
\[
a_j^*=\sum_{E_k\supset E_j^*}a_k.
\]
Then clearly $\varphi=\sum_{j=1}^{n}a_j^*\chi_{E_j^*}$. This is a decomposition already dealt with above because the $E_j^*$ are disjoint. Thus
\[
\int\varphi
=\sum_j a_j^*m(E_j^*)
=\sum_j\sum_{E_k\supset E_j^*}a_km(E_j^*)
=\sum_k a_km(E_k),
\]
and conclusion (i) is established.

Conclusion (ii) follows by using any representation of $\varphi$ and $\psi$, and the obvious linearity of (i).

For the additivity over sets, one must note that if $E$ and $F$ are disjoint, then
\[
\chi_{E\cup F}=\chi_E+\chi_F,
\]
and we may use the linearity of the integral to see that
\[
\int_{E\cup F}\varphi=\int_E\varphi+\int_F\varphi.
\]

If $\eta\geq0$ is a simple function, then its canonical form is everywhere non-negative, and therefore $\int\eta\geq0$ by the definition of the integral. Applying this argument to $\psi-\varphi$ gives the desired monotonicity property.

Finally, for the triangle inequality, it suffices to write $\varphi$ in its canonical form
\[
\varphi=\sum_{k=1}^{N}a_k\chi_{E_k}
\]
and observe that
\[
|\varphi|=\sum_{k=1}^{N}|a_k|\chi_{E_k}(x).
\]
Therefore,
\[
\left|\int\varphi\right|
=\left|\sum_{k=1}^{N}a_km(E_k)\right|
\leq\sum_{k=1}^{N}|a_k|m(E_k)
=\int|\varphi|.
\]
\end{proof}

Incidentally, it is worthwhile to point out the following easy fact: whenever $f$ and $g$ are a pair of simple functions that agree almost everywhere, then $\int f=\int g$. The identity of the integrals of two functions that agree almost everywhere will continue to hold for the successive definitions of the integral that follow.

\subsection*{Stage two: bounded functions supported on a set of finite measure}

The support of a measurable function $f$ is defined to be the set of all points where $f$ does not vanish,
\[
\operatorname{supp}(f)=\{x:f(x)\neq0\}.
\]
We shall also say that $f$ is supported on a set $E$, if $f(x)=0$ whenever $x\notin E$.

Since $f$ is measurable, so is the set $\operatorname{supp}(f)$. We shall next be interested in those bounded measurable functions that have $m(\operatorname{supp}(f))<\infty$.

An important result in the previous chapter (Theorem 4.2) states the following: if $f$ is a function bounded by $M$ and supported on a set $E$, then there exists a sequence $\{\varphi_n\}$ of simple functions, with each $\varphi_n$ bounded by $M$ and supported on $E$, and such that
\[
\varphi_n(x)\to f(x)\qquad\text{for all }x.
\]
The key lemma that follows allows us to define the integral for the class of bounded functions supported on sets of finite measure.

\begin{lemma}
Let $f$ be a bounded function supported on a set $E$ of finite measure. If $\{\varphi_n\}_{n=1}^{\infty}$ is any sequence of simple functions bounded by $M$, supported on $E$, and with $\varphi_n(x)\to f(x)$ for a.e. $x$, then:
\begin{enumerate}
\item[(i)] The limit $\displaystyle \lim_{n\to\infty}\int\varphi_n$ exists.
\item[(ii)] If $f=0$ a.e., then $\displaystyle \lim_{n\to\infty}\int\varphi_n=0$.
\end{enumerate}
\end{lemma}

\begin{proof}
The assertions of the lemma would be nearly obvious if we had that $\varphi_n$ converges to $f$ uniformly on $E$. Instead, we recall one of Littlewood's principles, which states that the convergence of a sequence of measurable functions is ``nearly'' uniform. The precise statement lying behind this principle is Egorov's theorem, which we proved in Chapter 1, and which we apply here.

Since the measure of $E$ is finite, given $\epsilon>0$ Egorov's theorem guarantees the existence of a (closed) measurable subset $A_\epsilon$ of $E$ such that $m(E-A_\epsilon)\leq\epsilon$, and $\varphi_n\to f$ uniformly on $A_\epsilon$. Therefore, setting $I_n=\int\varphi_n$, we have
\begin{align*}
|I_n-I_m|
&\leq\int_E|\varphi_n(x)-\varphi_m(x)|\,dx\\
&=\int_{A_\epsilon}|\varphi_n(x)-\varphi_m(x)|\,dx
 +\int_{E-A_\epsilon}|\varphi_n(x)-\varphi_m(x)|\,dx\\
&\leq\int_{A_\epsilon}|\varphi_n(x)-\varphi_m(x)|\,dx
 +2M\,m(E-A_\epsilon)\\
&\leq\int_{A_\epsilon}|\varphi_n(x)-\varphi_m(x)|\,dx+2M\epsilon.
\end{align*}
By the uniform convergence, one has, for all $x\in A_\epsilon$ and all large $n$ and $m$, the estimate $|\varphi_n(x)-\varphi_m(x)|<\epsilon$, so we deduce that
\[
|I_n-I_m|\leq m(E)\epsilon+2M\epsilon
\]
for all large $n$ and $m$. Since $\epsilon$ is arbitrary and $m(E)<\infty$, this proves that $\{I_n\}$ is a Cauchy sequence and hence converges.

For the second part, we note that if $f=0$, we may repeat the argument above to find that $|I_n|\leq m(E)\epsilon+M\epsilon$, which yields $\lim_{n\to\infty}I_n=0$.
\end{proof}

Using Lemma 1.2 we can now turn to the integration of bounded functions that are supported on sets of finite measure. For such a function $f$ we define its Lebesgue integral by
\[
\int f(x)\,dx
=\lim_{n\to\infty}\int\varphi_n(x)\,dx,
\]
where $\{\varphi_n\}$ is any sequence of simple functions satisfying: $|\varphi_n|\leq M$, each $\varphi_n$ is supported on the support of $f$, and $\varphi_n(x)\to f(x)$ for a.e. $x$. By the previous lemma, we know that this limit exists.

Next, we must show that $\int f$ is independent of the limiting sequence $\{\varphi_n\}$ used, in order for the integral to be well-defined. Suppose that $\{\psi_n\}$ is another sequence of simple functions that is bounded by $M$, supported on $\operatorname{supp}(f)$, and such that $\psi_n(x)\to f(x)$ for a.e. $x$. If $\eta_n=\varphi_n-\psi_n$, the sequence $\{\eta_n\}$ consists of simple functions bounded by $2M$, supported on a set of finite measure, and such that $\eta_n\to0$ a.e. Thus, by the second part of the lemma, $\int\eta_n\to0$. Consequently,
\[
\lim_{n\to\infty}\int\varphi_n(x)\,dx
=\lim_{n\to\infty}\int\psi_n(x)\,dx.
\]

If $E$ is a subset of $\mathbb{R}^d$ with finite measure, and $f$ is bounded with $m(\operatorname{supp}(f))<\infty$, then it is natural to define
\[
\int_E f(x)\,dx=\int f(x)\chi_E(x)\,dx.
\]
Clearly, if $f$ is itself simple, then $\int f$ as defined above coincides with the integral of simple functions studied earlier. This extension of the definition of integration also satisfies all the basic properties of the integral of simple functions.

\begin{proposition}
Suppose $f$ and $g$ are bounded functions supported on sets of finite measure. Then the following properties hold.
\begin{enumerate}
\item[(i)] \textbf{Linearity.} If $a,b\in\mathbb{R}$, then
\[
\int(af+bg)=a\int f+b\int g.
\]
\item[(ii)] \textbf{Additivity.} If $E$ and $F$ are disjoint subsets of $\mathbb{R}^d$, then
\[
\int_{E\cup F}f=\int_Ef+\int_Ff.
\]
\item[(iii)] \textbf{Monotonicity.} If $f\leq g$, then
\[
\int f\leq\int g.
\]
\item[(iv)] \textbf{Triangle inequality.} $|f|$ is also bounded, supported on a set of finite measure, and
\[
\left|\int f\right|\leq\int|f|.
\]
\end{enumerate}
\end{proposition}

All these properties follow by using approximations by simple functions, and the properties of the integral of simple functions given in Proposition 1.1.

We are now in a position to prove the first important convergence theorem.

\begin{theorem}[Bounded convergence theorem]
Suppose that $\{f_n\}$ is a sequence of measurable functions that are all bounded by $M$, are supported on a set $E$ of finite measure, and $f_n(x)\to f(x)$ a.e. as $n\to\infty$. Then $f$ is measurable, bounded, supported on $E$ for a.e. $x$, and
\[
\int|f_n-f|\to0\qquad\text{as }n\to\infty.
\]
Consequently,
\[
\int f_n\to\int f\qquad\text{as }n\to\infty.
\]
\end{theorem}

\begin{proof}
From the assumptions one sees at once that $f$ is bounded by $M$ almost everywhere and vanishes outside $E$, except possibly on a set of measure zero. Clearly, the triangle inequality for the integral implies that it suffices to prove that $\int|f_n-f|\to0$.

Given $\epsilon>0$, we may find, by Egorov's theorem, a measurable subset $A_\epsilon$ of $E$ such that $m(E-A_\epsilon)\leq\epsilon$ and $f_n\to f$ uniformly on $A_\epsilon$. Then, for all sufficiently large $n$, $|f_n(x)-f(x)|\leq\epsilon$ for all $x\in A_\epsilon$. Hence
\begin{align*}
\int|f_n(x)-f(x)|\,dx
&\leq\int_{A_\epsilon}|f_n(x)-f(x)|\,dx
 +\int_{E-A_\epsilon}|f_n(x)-f(x)|\,dx\\
&\leq\epsilon m(E)+2M\,m(E-A_\epsilon).
\end{align*}
Since $\epsilon$ is arbitrary, the proof is complete.
\end{proof}

The above convergence theorem is a statement about the interchange of an integral and a limit, since its conclusion simply says
\[
\lim_{n\to\infty}\int f_n
=\int\lim_{n\to\infty}f_n.
\]

A useful observation at this point is the following: if $f\geq0$ is bounded and supported on a set of finite measure $E$ and $\int f=0$, then $f=0$ almost everywhere. Indeed, if
\[
E_k=\{x\in E:f(x)\geq1/k\},
\]
then $k^{-1}\chi_{E_k}(x)\leq f(x)$ implies
\[
k^{-1}m(E_k)\leq\int f.
\]
Thus $m(E_k)=0$ for all $k$, and since $\{x:f(x)>0\}=\bigcup_{k=1}^{\infty}E_k$, we see that $f=0$ almost everywhere.

\subsection*{Return to Riemann integrable functions}

We shall now show that Riemann integrable functions are also Lebesgue integrable. When we combine this with the bounded convergence theorem we have just proved, we see that Lebesgue integration resolves the second problem in the Introduction.

\begin{theorem}
Suppose $f$ is Riemann integrable on the closed interval $[a,b]$. Then $f$ is measurable, and
\[
\int_{[a,b]}^{R}f(x)\,dx
=\int_{[a,b]}^{L}f(x)\,dx,
\]
where the integral on the left-hand side is the standard Riemann integral, and that on the right-hand side is the Lebesgue integral.
\end{theorem}

\begin{proof}
By definition, a Riemann integrable function is bounded, say $|f(x)|\leq M$, so we need to prove that $f$ is measurable, and then establish the equality of integrals.

By definition of Riemann integrability, we may construct two sequences of step functions $\{\varphi_k\}$ and $\{\psi_k\}$ satisfying
\[
|\varphi_k(x)|\leq M,\qquad |\psi_k(x)|\leq M,
\]
for all $x\in[a,b]$ and $k\geq1$,
\[
\varphi_1(x)\leq\varphi_2(x)\leq\cdots\leq f(x)
\leq\cdots\leq\psi_2(x)\leq\psi_1(x),
\]
and
\begin{equation}
\lim_{k\to\infty}\int_{[a,b]}^{R}\varphi_k(x)\,dx
=\lim_{k\to\infty}\int_{[a,b]}^{R}\psi_k(x)\,dx
=\int_{[a,b]}^{R}f(x)\,dx.
\tag{2}
\end{equation}
For step functions the Riemann and Lebesgue integrals agree; therefore
\begin{equation}
\int_{[a,b]}^{R}\varphi_k(x)\,dx
=\int_{[a,b]}^{L}\varphi_k(x)\,dx,
\qquad
\int_{[a,b]}^{R}\psi_k(x)\,dx
=\int_{[a,b]}^{L}\psi_k(x)\,dx.
\tag{3}
\end{equation}
If
\[
\widetilde\varphi(x)=\lim_{k\to\infty}\varphi_k(x),
\qquad
\widetilde\psi(x)=\lim_{k\to\infty}\psi_k(x),
\]
then $\widetilde\varphi\leq f\leq\widetilde\psi$, and both $\widetilde\varphi$ and $\widetilde\psi$ are measurable. The bounded convergence theorem yields
\[
\lim_{k\to\infty}\int_{[a,b]}^{L}\varphi_k(x)\,dx
=\int_{[a,b]}^{L}\widetilde\varphi(x)\,dx,
\]
and similarly for $\psi_k$. Together with (2) and (3), this gives
\[
\int_{[a,b]}^{L}\bigl(\widetilde\psi(x)-\widetilde\varphi(x)\bigr)\,dx=0.
\]
Since $\widetilde\psi-\widetilde\varphi\geq0$, the observation following the bounded convergence theorem implies $\widetilde\psi=\widetilde\varphi$ a.e. Hence $f=\widetilde\varphi=\widetilde\psi$ a.e., which proves that $f$ is measurable. Finally, since $\varphi_k\to f$ almost everywhere,
\[
\lim_{k\to\infty}\int_{[a,b]}^{L}\varphi_k(x)\,dx
=\int_{[a,b]}^{L}f(x)\,dx,
\]
and (2)--(3) give the desired equality of the Riemann and Lebesgue integrals.
\end{proof}

\subsection*{Stage three: non-negative functions}

We proceed with the integrals of functions that are measurable and non-negative but not necessarily bounded. It will be important to allow these functions to be extended-valued, that is, these functions may take on the value $+\infty$ (on a measurable set). We recall in this connection the convention that one defines the supremum of a set of positive numbers to be $+\infty$ if the set is unbounded.

In the case of such a function $f$ we define its (extended) Lebesgue integral by
\[
\int f(x)\,dx=\sup_g\int g(x)\,dx,
\]
where the supremum is taken over all measurable functions $g$ such that $0\leq g\leq f$, and where $g$ is bounded and supported on a set of finite measure.

With the above definition there are only two possible cases; the supremum is either finite or infinite. In the first case, when $\int f(x)\,dx<\infty$, we shall say that $f$ is Lebesgue integrable or simply integrable.

Clearly, if $E$ is any measurable subset of $\mathbb{R}^d$, and $f\geq0$, then $f\chi_E$ is also positive, and we define
\[
\int_E f(x)\,dx=\int f(x)\chi_E(x)\,dx.
\]
Simple examples of functions on $\mathbb{R}^d$ that are integrable (or non-integrable) are
\[
f_a(x)=
\begin{cases}
|x|^{-a},&|x|\leq1,\\
0,&|x|>1,
\end{cases}
\qquad
F_a(x)=\frac{1}{1+|x|^a},\quad x\in\mathbb{R}^d.
\]
Then $f_a$ is integrable exactly when $a<d$, while $F_a$ is integrable exactly when $a>d$.

\begin{proposition}
The integral of non-negative measurable functions enjoys the following properties:
\begin{enumerate}
\item[(i)] \textbf{Linearity.} If $f,g\geq0$, and $a,b$ are positive real numbers, then
\[
\int(af+bg)=a\int f+b\int g.
\]
\item[(ii)] \textbf{Additivity.} If $E$ and $F$ are disjoint subsets of $\mathbb{R}^d$, and $f\geq0$, then
\[
\int_{E\cup F}f=\int_Ef+\int_Ff.
\]
\item[(iii)] \textbf{Monotonicity.} If $0\leq f\leq g$, then
\[
\int f\leq\int g.
\]
\item[(iv)] If $g$ is integrable and $0\leq f\leq g$, then $f$ is integrable.
\item[(v)] If $f$ is integrable, then $f(x)<\infty$ for almost every $x$.
\item[(vi)] If $\int f=0$, then $f(x)=0$ for almost every $x$.
\end{enumerate}
\end{proposition}

\begin{proof}
Of the first four assertions, only (i) is not an immediate consequence of the definitions. Taking $a=b=1$, suppose $\varphi\leq f$ and $\psi\leq g$, where both $\varphi$ and $\psi$ are bounded and supported on sets of finite measure. Then $\varphi+\psi\leq f+g$, and $\varphi+\psi$ is also bounded and supported on a set of finite measure. Consequently
\[
\int f+\int g\leq\int(f+g).
\]
To prove the reverse inequality, suppose $\eta$ is bounded and supported on a set of finite measure, and $\eta\leq f+g$. If
\[
\eta_1(x)=\min(f(x),\eta(x)),\qquad \eta_2=\eta-\eta_1,
\]
then $\eta_1\leq f$ and $\eta_2\leq g$, and both are bounded and supported on sets of finite measure. Hence
\[
\int\eta=\int\eta_1+\int\eta_2\leq\int f+\int g.
\]
Taking the supremum over $\eta$ yields the required inequality.

To prove (v), suppose
\[
E_k=\{x:f(x)\geq k\},\qquad E_\infty=\{x:f(x)=\infty\}.
\]
Then
\[
\int f\geq\int\chi_{E_k}f\geq k\,m(E_k),
\]
hence $m(E_k)\to0$ as $k\to\infty$. Since $E_k\searrow E_\infty$, Corollary 3.3 in the previous chapter implies $m(E_\infty)=0$.

The proof of (vi) is the same as the observation following Theorem 1.4.
\end{proof}

We now turn our attention to some important convergence theorems for the class of non-negative measurable functions. Suppose $f_n\geq0$ and $f_n(x)\to f(x)$ for almost every $x$. Is it true that $\int f_n\to\int f$? The example
\[
f_n(x)=
\begin{cases}
n,&0<x<1/n,\\
0,&\text{otherwise}
\end{cases}
\]
shows that the answer is negative: $f_n(x)\to0$ for all $x$, yet $\int f_n(x)\,dx=1$ for all $n$.

\begin{lemma}[Fatou]
Suppose $\{f_n\}$ is a sequence of measurable functions with $f_n\geq0$. If $\lim_{n\to\infty}f_n(x)=f(x)$ for a.e. $x$, then
\[
\int f\leq\liminf_{n\to\infty}\int f_n.
\]
\end{lemma}

\begin{proof}
Suppose $0\leq g\leq f$, where $g$ is bounded and supported on a set $E$ of finite measure. If $g_n(x)=\min(g(x),f_n(x))$, then $g_n$ is measurable, supported on $E$, and $g_n(x)\to g(x)$ a.e. By the bounded convergence theorem,
\[
\int g_n\to\int g.
\]
By construction $g_n\leq f_n$, and therefore
\[
\int g\leq\liminf_{n\to\infty}\int f_n.
\]
Taking the supremum over all such $g$ yields the desired inequality.
\end{proof}

\begin{corollary}
Suppose $f$ is a non-negative measurable function, and $\{f_n\}$ a sequence of non-negative measurable functions with $f_n(x)\leq f(x)$ and $f_n(x)\to f(x)$ for almost every $x$. Then
\[
\lim_{n\to\infty}\int f_n=\int f.
\]
\end{corollary}

\begin{proof}
Since $f_n\leq f$ a.e., $\int f_n\leq\int f$ for all $n$, so
\[
\limsup_{n\to\infty}\int f_n\leq\int f.
\]
Together with Fatou's lemma, this proves the desired limit.
\end{proof}

In analogy with the symbols $\nearrow$ and $\searrow$ used to describe increasing and decreasing sequences of sets, we shall write $f_n\nearrow f$ whenever
\[
f_n(x)\leq f_{n+1}(x)\quad\text{a.e. for all }n\geq1,
\qquad
f_n(x)\to f(x)\quad\text{a.e.},
\]
and $f_n\searrow f$ for the analogous decreasing situation.

\begin{corollary}[Monotone convergence theorem]
Suppose $\{f_n\}$ is a sequence of non-negative measurable functions with $f_n\nearrow f$. Then
\[
\lim_{n\to\infty}\int f_n=\int f.
\]
\end{corollary}

The monotone convergence theorem has the following useful consequence.

\begin{corollary}
Consider a series $\sum_{k=1}^{\infty}a_k(x)$, where $a_k(x)\geq0$ is measurable for every $k\geq1$. Then
\[
\int\sum_{k=1}^{\infty}a_k(x)\,dx
=\sum_{k=1}^{\infty}\int a_k(x)\,dx.
\]
If $\sum_{k=1}^{\infty}\int a_k(x)\,dx$ is finite, then the series $\sum_{k=1}^{\infty}a_k(x)$ converges for a.e. $x$.
\end{corollary}

\begin{proof}
Let
\[
f_n(x)=\sum_{k=1}^{n}a_k(x),
\qquad
f(x)=\sum_{k=1}^{\infty}a_k(x).
\]
Then $f_n\nearrow f$, and
\[
\int f_n=\sum_{k=1}^{n}\int a_k(x)\,dx.
\]
The monotone convergence theorem gives the asserted identity. If the sum of the integrals is finite, the right-hand side shows that $f$ is integrable, and hence finite almost everywhere.
\end{proof}

We give two illustrations. First, the preceding corollary gives another proof of the Borel--Cantelli lemma: if $E_1,E_2,\ldots$ are measurable and $\sum m(E_k)<\infty$, let $a_k=\chi_{E_k}$. A point belongs to infinitely many $E_k$ if and only if $\sum a_k(x)=\infty$. Since
\[
\sum_{k=1}^{\infty}\int a_k(x)\,dx=\sum_{k=1}^{\infty}m(E_k)<\infty,
\]
the sum $\sum a_k(x)$ is finite almost everywhere.

For the second illustration, consider
\[
f(x)=
\begin{cases}
|x|^{-(d+1)},&x\neq0,\\
0,&x=0.
\end{cases}
\]
We prove that $f$ is integrable outside any ball $|x|\geq\epsilon$, and moreover
\[
\int_{|x|\geq\epsilon}f(x)\,dx\leq\frac{C}{\epsilon}
\]
for some constant $C>0$. Let
\[
A_k=\{x\in\mathbb{R}^d:2^k\epsilon<|x|\leq2^{k+1}\epsilon\}
\]
and define
\[
g(x)=\sum_{k=0}^{\infty}a_k(x),
\qquad
a_k(x)=\frac{1}{(2^k\epsilon)^{d+1}}\chi_{A_k}(x).
\]
Then $f\leq g$. Since $A_k$ is obtained from $A=\{1<|x|<2\}$ by dilation by $2^k\epsilon$,
\[
m(A_k)=(2^k\epsilon)^d m(A).
\]
Thus
\[
\int g
=\sum_{k=0}^{\infty}\frac{m(A_k)}{(2^k\epsilon)^{d+1}}
=m(A)\sum_{k=0}^{\infty}\frac{(2^k\epsilon)^d}{(2^k\epsilon)^{d+1}}
=\frac{C}{\epsilon},
\]
where $C=2m(A)$. The same dilation-invariance property also gives
\[
\int_{|x|\geq\epsilon}\frac{dx}{|x|^{d+1}}
=\frac{1}{\epsilon}\int_{|x|\geq1}\frac{dx}{|x|^{d+1}}.
\]

\subsection*{Stage four: general case}

If $f$ is any real-valued measurable function on $\mathbb{R}^d$, we say that $f$ is Lebesgue integrable (or just integrable) if the non-negative measurable function $|f|$ is integrable in the sense of the previous section.

If $f$ is Lebesgue integrable, define
\[
f^+(x)=\max(f(x),0),
\qquad
f^-(x)=\max(-f(x),0).
\]
Then $f=f^+-f^-$ and $f^\pm\leq|f|$, so both $f^+$ and $f^-$ are integrable. We define
\[
\int f=\int f^+-\int f^-.
\]
More generally, if $f=f_1-f_2$ with $f_1,f_2\geq0$ integrable, then
\[
\int f=\int f_1-\int f_2.
\]
Indeed, if also $f=g_1-g_2$, then $f_1+g_2=g_1+f_2$, and linearity for non-negative functions gives the equality of the two differences of integrals.

Both the integrability of $f$ and the value of its integral are unchanged if we modify $f$ arbitrarily on a set of measure zero. It is therefore useful to adopt the convention that, in the context of integration, functions may be undefined on sets of measure zero. Moreover, an integrable function is finite-valued almost everywhere. Thus, when speaking of a function $f$, we are in effect also speaking about the collection of all functions that equal $f$ almost everywhere.

\begin{proposition}
The integral of Lebesgue integrable functions is linear, additive, monotonic, and satisfies the triangle inequality.
\end{proposition}

We now gather two results which, although instructive in their own right, are also needed in the proof of the next theorem.

\begin{proposition}
Suppose $f$ is integrable on $\mathbb{R}^d$. Then for every $\epsilon>0$:
\begin{enumerate}
\item[(i)] There exists a set of finite measure $B$ (a ball, for example) such that
\[
\int_{B^c}|f|<\epsilon.
\]
\item[(ii)] There is a $\delta>0$ such that
\[
\int_E|f|<\epsilon\qquad\text{whenever }m(E)<\delta.
\]
\end{enumerate}
The last condition is known as absolute continuity.
\end{proposition}

\begin{proof}
By replacing $f$ with $|f|$ we may assume $f\geq0$.

For the first part, let $B_N$ denote the ball of radius $N$ centered at the origin and set $f_N=f\chi_{B_N}$. Then $f_N\nearrow f$, so by the monotone convergence theorem
\[
\lim_{N\to\infty}\int f_N=\int f.
\]
Thus for some large $N$,
\[
0\leq\int f-\int f\chi_{B_N}<\epsilon,
\]
which is exactly $\int_{B_N^c}f<\epsilon$.

For the second part, let
\[
E_N=\{x:f(x)\leq N\},
\qquad
f_N=f\chi_{E_N}.
\]
Again $f_N\nearrow f$, so choose $N$ such that
\[
\int(f-f_N)<\frac{\epsilon}{2}.
\]
Pick $\delta>0$ so that $N\delta<\epsilon/2$. If $m(E)<\delta$, then
\begin{align*}
\int_Ef
&=\int_E(f-f_N)+\int_Ef_N\\
&\leq\int(f-f_N)+N\,m(E)\\
&<\frac{\epsilon}{2}+\frac{\epsilon}{2}=\epsilon.
\end{align*}
\end{proof}

Intuitively, integrable functions should in some sense vanish at infinity since their integrals are finite, and the first part of the proposition attaches a precise meaning to this intuition. One should observe, however, that integrability need not guarantee pointwise vanishing as $|x|$ becomes large.

We are now ready to prove a cornerstone of the theory of Lebesgue integration, the dominated convergence theorem.

\begin{theorem}[Dominated convergence theorem]
Suppose $\{f_n\}$ is a sequence of measurable functions such that $f_n(x)\to f(x)$ a.e. as $n\to\infty$. If $|f_n(x)|\leq g(x)$, where $g$ is integrable, then
\[
\int|f_n-f|\to0\qquad\text{as }n\to\infty,
\]
and consequently
\[
\int f_n\to\int f\qquad\text{as }n\to\infty.
\]
\end{theorem}

\begin{proof}
For each $N\geq0$ let
\[
E_N=\{x:|x|\leq N,\ g(x)\leq N\}.
\]
Given $\epsilon>0$, choose $N$ so that $\int_{E_N^c}g<\epsilon$. Then the functions $f_n\chi_{E_N}$ are bounded by $N$ and supported on a set of finite measure, so by the bounded convergence theorem,
\[
\int_{E_N}|f_n-f|<\epsilon
\]
for all large $n$. Hence
\begin{align*}
\int|f_n-f|
&=\int_{E_N}|f_n-f|+\int_{E_N^c}|f_n-f|\\
&\leq\int_{E_N}|f_n-f|+2\int_{E_N^c}g\\
&\leq3\epsilon
\end{align*}
for all large $n$.
\end{proof}

\subsection*{Complex-valued functions}

If $f$ is a complex-valued function on $\mathbb{R}^d$, write
\[
f(x)=u(x)+iv(x),
\]
where $u$ and $v$ are the real and imaginary parts of $f$. The function $f$ is measurable if and only if both $u$ and $v$ are measurable. We say that $f$ is Lebesgue integrable if
\[
|f(x)|=(u(x)^2+v(x)^2)^{1/2}
\]
is integrable.

Since
\[
|u(x)|\leq|f(x)|,
\qquad
|v(x)|\leq|f(x)|,
\qquad
|f(x)|\leq|u(x)|+|v(x)|,
\]
a complex-valued function is integrable if and only if both its real and imaginary parts are integrable. We define
\[
\int f(x)\,dx
=\int u(x)\,dx+i\int v(x)\,dx.
\]

If $E$ is a measurable subset of $\mathbb{R}^d$, and $f$ is a complex-valued measurable function on $E$, we say that $f$ is Lebesgue integrable on $E$ if $f\chi_E$ is integrable on $\mathbb{R}^d$, and define
\[
\int_Ef=\int f\chi_E.
\]
The collection of all complex-valued integrable functions on a measurable subset $E$ forms a vector space over $\mathbb{C}$. Indeed, if $f$ and $g$ are integrable, then so is $f+g$, since
\[
|(f+g)(x)|\leq|f(x)|+|g(x)|,
\]
and hence
\[
\int_E|f+g|\leq\int_E|f|+\int_E|g|<\infty.
\]
Also, if $a\in\mathbb{C}$ and $f$ is integrable, then so is $af$, and the integral continues to be linear over $\mathbb{C}$.

\section{The space $L^1$ of integrable functions}

The fact that the integrable functions form a vector space is an important observation about the algebraic properties of such functions. A fundamental analytic fact is that this vector space is complete in the appropriate norm.

For any integrable function $f$ on $\mathbb{R}^d$ we define the norm of $f$ by
\[
\|f\|=\|f\|_{L^1}=\|f\|_{L^1(\mathbb{R}^d)}
=\int_{\mathbb{R}^d}|f(x)|\,dx.
\]
The collection of all integrable functions with the above norm gives a somewhat imprecise definition of the space $L^1(\mathbb{R}^d)$. We also note that $\|f\|=0$ if and only if $f=0$ almost everywhere, and this simple property of the norm reflects the practice we have already adopted not to distinguish two functions that agree almost everywhere.

With this in mind, we take the precise definition of $L^1(\mathbb{R}^d)$ to be the space of equivalence classes of integrable functions, where two functions are equivalent if they agree almost everywhere. Often, however, it is convenient to retain the terminology that an element $f\in L^1(\mathbb{R}^d)$ is an integrable function, even though it is only an equivalence class of such functions. The norm is well-defined by the choice of any integrable function in its equivalence class. Moreover, $L^1(\mathbb{R}^d)$ inherits the property that it is a vector space.

\begin{proposition}
Suppose $f$ and $g$ are two functions in $L^1(\mathbb{R}^d)$.
\begin{enumerate}
\item[(i)] $\|af\|_{L^1(\mathbb{R}^d)}=|a|\,\|f\|_{L^1(\mathbb{R}^d)}$ for all $a\in\mathbb{C}$.
\item[(ii)] $\|f+g\|_{L^1(\mathbb{R}^d)}\leq\|f\|_{L^1(\mathbb{R}^d)}+\|g\|_{L^1(\mathbb{R}^d)}$.
\item[(iii)] $\|f\|_{L^1(\mathbb{R}^d)}=0$ if and only if $f=0$ a.e.
\item[(iv)] $d(f,g)=\|f-g\|_{L^1(\mathbb{R}^d)}$ defines a metric on $L^1(\mathbb{R}^d)$.
\end{enumerate}
\end{proposition}

In (iv), we mean that $d(f,g)\geq0$ for all integrable functions $f$ and $g$, and $d(f,g)=0$ if and only if $f=g$ a.e. Also $d(f,g)=d(g,f)$, and
\[
d(f,g)\leq d(f,h)+d(h,g),
\qquad f,g,h\in L^1(\mathbb{R}^d).
\]

A space $V$ with a metric $d$ is said to be complete if for every Cauchy sequence $\{x_k\}$ in $V$---that is, $d(x_k,x_\ell)\to0$ as $k,\ell\to\infty$---there exists $x\in V$ such that
\[
d(x_k,x)\to0\qquad\text{as }k\to\infty.
\]
Our main goal of completing the space of Riemann integrable functions will be attained once we establish the next theorem.

\begin{theorem}[Riesz--Fischer]
The vector space $L^1$ is complete in its metric.
\end{theorem}

\begin{proof}
Suppose $\{f_n\}$ is a Cauchy sequence in the norm, so that $\|f_n-f_m\|\to0$ as $n,m\to\infty$. The plan is to extract a subsequence of $\{f_n\}$ that converges to a function $f$, both pointwise almost everywhere and in the norm.

Almost-everywhere convergence does not hold for general Cauchy sequences. The main point is that if the convergence in norm is rapid enough, then almost-everywhere convergence follows, and this can be achieved by passing to an appropriate subsequence.

Choose a subsequence $\{f_{n_k}\}_{k=1}^{\infty}$ such that
\[
\|f_{n_{k+1}}-f_{n_k}\|\leq2^{-k},
\qquad k\geq1.
\]
The existence of such a subsequence follows from the Cauchy property. Consider the series
\[
f(x)=f_{n_1}(x)+\sum_{k=1}^{\infty}\bigl(f_{n_{k+1}}(x)-f_{n_k}(x)\bigr)
\]
and
\[
g(x)=|f_{n_1}(x)|+\sum_{k=1}^{\infty}|f_{n_{k+1}}(x)-f_{n_k}(x)|.
\]
Then
\[
\int|f_{n_1}|+\sum_{k=1}^{\infty}\int|f_{n_{k+1}}-f_{n_k}|
\leq\int|f_{n_1}|+\sum_{k=1}^{\infty}2^{-k}<\infty.
\]
By the monotone convergence theorem, $g$ is integrable. Since $|f|\leq g$, $f$ is integrable as well. In particular, the series defining $f$ converges almost everywhere, and its partial sums are precisely the $f_{n_k}$. Thus
\[
f_{n_k}(x)\to f(x)\qquad\text{a.e. }x.
\]

To prove convergence in $L^1$, observe that $|f-f_{n_k}|\leq g$ and apply the dominated convergence theorem to obtain
\[
\|f_{n_k}-f\|_{L^1}\to0.
\]
Finally, use the fact that the original sequence is Cauchy. Given $\epsilon>0$, choose $N$ such that $\|f_n-f_m\|<\epsilon/2$ whenever $n,m>N$. Choose $n_k>N$ also satisfying $\|f_{n_k}-f\|<\epsilon/2$. Then for every $n>N$,
\[
\|f_n-f\|\leq\|f_n-f_{n_k}\|+\|f_{n_k}-f\|<\epsilon.
\]
Hence $f_n\to f$ in $L^1$.
\end{proof}

Since every sequence that converges in norm is a Cauchy sequence in that norm, the argument in the proof yields the following.

\begin{corollary}
If $\{f_n\}_{n=1}^{\infty}$ converges to $f$ in $L^1$, then there exists a subsequence $\{f_{n_k}\}_{k=1}^{\infty}$ such that
\[
f_{n_k}(x)\to f(x)\qquad\text{a.e. }x.
\]
\end{corollary}

We say that a family $\mathcal{G}$ of integrable functions is dense in $L^1$ if for any $f\in L^1$ and $\epsilon>0$, there exists $g\in\mathcal{G}$ so that $\|f-g\|_{L^1}<\epsilon$. Fortunately, we are familiar with many families that are dense in $L^1$.

\begin{theorem}
The following families of functions are dense in $L^1(\mathbb{R}^d)$:
\begin{enumerate}
\item[(i)] The simple functions.
\item[(ii)] The step functions.
\item[(iii)] The continuous functions of compact support.
\end{enumerate}
\end{theorem}

\begin{proof}
Let $f$ be an integrable function on $\mathbb{R}^d$. We may assume that $f$ is real-valued, because we may approximate its real and imaginary parts independently. Writing $f=f^+-f^-$ reduces the proof to the case $f\geq0$.

For (i), Theorem 4.1 in Chapter 1 gives a sequence $\{\varphi_k\}$ of non-negative simple functions that increase pointwise to $f$. By the dominated convergence theorem (or the monotone convergence theorem),
\[
\|f-\varphi_k\|_{L^1}\to0.
\]

For (ii), by (i) it suffices to approximate simple functions by step functions. A simple function is a finite linear combination of characteristic functions of sets of finite measure, so it is enough to approximate $\chi_E$. The proof of Theorem 4.3 in Chapter 1 provides an almost disjoint family of rectangles $\{R_j\}_{j=1}^{M}$ such that
\[
m\left(E\triangle\bigcup_{j=1}^{M}R_j\right)\leq2\epsilon.
\]
Thus if $\psi=\sum_j\chi_{R_j}$, then
\[
\|\chi_E-\psi\|_{L^1}<2\epsilon.
\]

For (iii), by (ii) it suffices to consider the characteristic function of a rectangle. In one dimension, if $f=\chi_{[a,b]}$, choose a continuous piecewise linear function $g$ satisfying
\[
g(x)=
\begin{cases}
1,&a\leq x\leq b,\\
0,&x\leq a-\epsilon\text{ or }x\geq b+\epsilon,
\end{cases}
\]
and linear on $[a-\epsilon,a]$ and $[b,b+\epsilon]$. Then $\|f-g\|_{L^1}<2\epsilon$. In $d$ dimensions, the characteristic function of a rectangle is the product of characteristic functions of intervals, and the corresponding approximation is the product of such one-dimensional continuous functions.
\end{proof}

The preceding results for $L^1(\mathbb{R}^d)$ lead immediately to an extension in which $\mathbb{R}^d$ is replaced by any fixed measurable subset $E$ of positive measure. One can define $L^1(E)$ directly, or extend a function $f$ on $E$ by setting $\widetilde f=f$ on $E$ and $\widetilde f=0$ on $E^c$, and then define
\[
\|f\|_{L^1(E)}=\|\widetilde f\|_{L^1(\mathbb{R}^d)}.
\]
The analogues of Proposition 2.1 and Theorem 2.2 then hold for $L^1(E)$.

\subsection*{Invariance properties}

If $f$ is a function defined on $\mathbb{R}^d$, the translation of $f$ by a vector $h\in\mathbb{R}^d$ is the function $f_h$ defined by
\[
f_h(x)=f(x-h).
\]
First, there is the translation-invariance of the integral: if $f$ is integrable, then so is $f_h$ and
\begin{equation}
\int_{\mathbb{R}^d}f(x-h)\,dx
=\int_{\mathbb{R}^d}f(x)\,dx.
\tag{4}
\end{equation}
For $f=\chi_E$, this is exactly the translation-invariance of Lebesgue measure. By linearity it follows for simple functions; by monotone convergence it follows for non-negative functions; and then for general integrable functions by decomposing into positive and negative (or real and imaginary) parts. In particular,
\[
\|f_h\|=\|f\|.
\]

Using the relative invariance of Lebesgue measure under dilations and reflections, one obtains similarly
\begin{equation}
\delta^d\int_{\mathbb{R}^d}f(\delta x)\,dx
=\int_{\mathbb{R}^d}f(x)\,dx,
\qquad
\int_{\mathbb{R}^d}f(-x)\,dx
=\int_{\mathbb{R}^d}f(x)\,dx,
\tag{5}
\end{equation}
for $\delta>0$.

We record two useful consequences.

First, suppose $f$ and $g$ are measurable functions on $\mathbb{R}^d$ so that, for some fixed $x\in\mathbb{R}^d$, the function $y\mapsto f(x-y)g(y)$ is integrable. Then $y\mapsto f(y)g(x-y)$ is also integrable and
\begin{equation}
\int_{\mathbb{R}^d}f(x-y)g(y)\,dy
=\int_{\mathbb{R}^d}f(y)g(x-y)\,dy.
\tag{6}
\end{equation}
This follows by the change of variables $y\mapsto x-y$, a combination of a translation and a reflection. The integral on the left is denoted by $(f*g)(x)$ and is called the convolution of $f$ and $g$; thus (6) asserts the commutativity of convolution.

Second, using (5), for all $\epsilon>0$,
\begin{equation}
\int_{|x|\geq\epsilon}\frac{dx}{|x|^a}
=\epsilon^{-a+d}\int_{|x|\geq1}\frac{dx}{|x|^a},
\qquad a>d,
\tag{7}
\end{equation}
and
\begin{equation}
\int_{|x|\leq\epsilon}\frac{dx}{|x|^a}
=\epsilon^{-a+d}\int_{|x|\leq1}\frac{dx}{|x|^a},
\qquad a<d.
\tag{8}
\end{equation}
The integrals on the right are finite in the indicated ranges, by the argument following Corollary 1.10.

\subsection*{Translations and continuity}

We shall next examine how continuity properties of $f$ are related to the way the translations $f_h$ vary with $h$. For any given $x\in\mathbb{R}^d$, the statement $f_h(x)\to f(x)$ as $h\to0$ is the same as continuity of $f$ at $x$. A general integrable function may be discontinuous at every point, even after modification on a set of measure zero. Nevertheless, an arbitrary $f\in L^1(\mathbb{R}^d)$ enjoys an overall continuity in norm.

\begin{proposition}
Suppose $f\in L^1(\mathbb{R}^d)$. Then
\[
\|f_h-f\|_{L^1}\to0\qquad\text{as }h\to0.
\]
\end{proposition}

\begin{proof}
For any $\epsilon>0$, choose by Theorem 2.4 a continuous compactly supported function $g$ such that $\|f-g\|<\epsilon$. Then
\[
f_h-f=(g_h-g)+(f_h-g_h)-(f-g).
\]
Translation-invariance gives
\[
\|f_h-g_h\|=\|f-g\|<\epsilon.
\]
Since $g$ is continuous and compactly supported,
\[
\|g_h-g\|
=\int_{\mathbb{R}^d}|g(x-h)-g(x)|\,dx\to0
\]
as $h\to0$. Therefore, for $|h|$ sufficiently small,
\[
\|f_h-f\|<3\epsilon.
\]
\end{proof}

\section{Fubini's theorem}

In elementary calculus integrals of continuous functions of several variables are often calculated by iterating one-dimensional integrals. We shall now examine this important analytic device from the general point of view of Lebesgue integration in $\mathbb{R}^d$.

In general, we may write
\[
\mathbb{R}^d=\mathbb{R}^{d_1}\times\mathbb{R}^{d_2},
\qquad d=d_1+d_2,\quad d_1,d_2\geq1.
\]
A point in $\mathbb{R}^d$ then takes the form $(x,y)$, where $x\in\mathbb{R}^{d_1}$ and $y\in\mathbb{R}^{d_2}$. If $f$ is a function on $\mathbb{R}^{d_1}\times\mathbb{R}^{d_2}$, the slice corresponding to $y\in\mathbb{R}^{d_2}$ is the function $f^y$ of $x\in\mathbb{R}^{d_1}$ given by
\[
f^y(x)=f(x,y).
\]
Similarly, for fixed $x\in\mathbb{R}^{d_1}$ we write
\[
f_x(y)=f(x,y).
\]
For a set $E\subset\mathbb{R}^{d_1}\times\mathbb{R}^{d_2}$ we define its slices by
\[
E^y=\{x\in\mathbb{R}^{d_1}:(x,y)\in E\},
\qquad
E_x=\{y\in\mathbb{R}^{d_2}:(x,y)\in E\}.
\]

\begin{figure}[ht]
\centering
\fbox{\parbox{0.72\textwidth}{\centering Figure from the original text omitted in this transcription.}}
\caption{Slices $E^y$ and $E_x$ (for fixed $x$ and $y$) of a set $E$.}
\end{figure}

\subsection{Statement and proof of the theorem}

The theorem is not entirely straightforward because of measurability. Even if $f$ is measurable on $\mathbb{R}^d$, the slice $f^y$ need not be measurable on $\mathbb{R}^{d_1}$ for every $y$; similarly, if $E$ is measurable, $E^y$ need not be measurable for every $y$. An easy example arises in $\mathbb{R}^2$ by placing a one-dimensional non-measurable set on the $x$-axis: the resulting set has measure zero in $\mathbb{R}^2$, but the slice at $y=0$ is non-measurable. What saves us is that measurability holds for almost all slices.

\begin{theorem}
Suppose $f(x,y)$ is integrable on $\mathbb{R}^{d_1}\times\mathbb{R}^{d_2}$. Then for almost every $y\in\mathbb{R}^{d_2}$:
\begin{enumerate}
\item[(i)] The slice $f^y$ is integrable on $\mathbb{R}^{d_1}$.
\item[(ii)] The function
\[
y\longmapsto\int_{\mathbb{R}^{d_1}}f^y(x)\,dx
\]
is integrable on $\mathbb{R}^{d_2}$.
\end{enumerate}
Moreover,
\[
\int_{\mathbb{R}^{d_2}}
\left(\int_{\mathbb{R}^{d_1}}f(x,y)\,dx\right)dy
=\int_{\mathbb{R}^d}f.
\]
\end{theorem}

Clearly, the theorem is symmetric in $x$ and $y$. Thus for almost every $x$, the slice $f_x$ is integrable on $\mathbb{R}^{d_2}$, and
\[
\int_{\mathbb{R}^{d_1}}
\left(\int_{\mathbb{R}^{d_2}}f(x,y)\,dy\right)dx
=\int_{\mathbb{R}^d}f.
\]
In particular,
\[
\int_{\mathbb{R}^{d_2}}
\left(\int_{\mathbb{R}^{d_1}}f(x,y)\,dx\right)dy
=
\int_{\mathbb{R}^{d_1}}
\left(\int_{\mathbb{R}^{d_2}}f(x,y)\,dy\right)dx
=
\int_{\mathbb{R}^d}f.
\]

We may assume that $f$ is real-valued, since the theorem then applies to the real and imaginary parts of a complex-valued function. The proof consists of six steps. Let $\mathcal{F}$ denote the set of integrable functions on $\mathbb{R}^d$ that satisfy all three conclusions of the theorem. We prove that $L^1(\mathbb{R}^d)\subset\mathcal{F}$.

\paragraph{Step 1.} Any finite linear combination of functions in $\mathcal{F}$ also belongs to $\mathcal{F}$.

Indeed, let $\{f_k\}_{k=1}^{N}\subset\mathcal{F}$. For each $k$ there exists a set $A_k\subset\mathbb{R}^{d_2}$ of measure zero so that $f_k^y$ is integrable on $\mathbb{R}^{d_1}$ whenever $y\notin A_k$. If $A=\bigcup_{k=1}^{N}A_k$, then $A$ has measure zero, and outside $A$ the $y$-slice of any finite linear combination is measurable and integrable. Linearity of the integral gives the remaining conclusion.

\paragraph{Step 2.} Suppose $\{f_k\}$ is a sequence of measurable functions in $\mathcal{F}$ so that $f_k\nearrow f$ or $f_k\searrow f$, where $f$ is integrable on $\mathbb{R}^d$. Then $f\in\mathcal{F}$.

Replacing $f_k$ by $-f_k$ if necessary, it suffices to consider the increasing case. Replacing $f_k$ by $f_k-f_1$, we may assume $f_k\geq0$. By the monotone convergence theorem,
\begin{equation}
\lim_{k\to\infty}\int_{\mathbb{R}^d}f_k(x,y)\,dx\,dy
=\int_{\mathbb{R}^d}f(x,y)\,dx\,dy.
\tag{9}
\end{equation}
For each $k$ there exists $A_k\subset\mathbb{R}^{d_2}$ of measure zero such that $f_k^y$ is integrable whenever $y\notin A_k$. Let $A=\bigcup_{k=1}^{\infty}A_k$. For $y\notin A$, the monotone convergence theorem gives
\[
g_k(y)=\int_{\mathbb{R}^{d_1}}f_k^y(x)\,dx
\nearrow
g(y)=\int_{\mathbb{R}^{d_1}}f^y(x)\,dx.
\]
Each $g_k$ is integrable, so another application of monotone convergence gives
\begin{equation}
\int_{\mathbb{R}^{d_2}}g_k(y)\,dy
\longrightarrow
\int_{\mathbb{R}^{d_2}}g(y)\,dy.
\tag{10}
\end{equation}
Since $f_k\in\mathcal{F}$,
\[
\int_{\mathbb{R}^{d_2}}g_k(y)\,dy
=\int_{\mathbb{R}^d}f_k(x,y)\,dx\,dy.
\]
Combining this with (9) and (10),
\[
\int_{\mathbb{R}^{d_2}}g(y)\,dy
=\int_{\mathbb{R}^d}f(x,y)\,dx\,dy.
\]
The right-hand side is finite, so $g$ is integrable. Hence $g(y)<\infty$ for a.e. $y$, which implies $f^y$ is integrable for a.e. $y$, and all conclusions follow.

\paragraph{Step 3.} Any characteristic function of a set $E$ that is a $G_\delta$ and has finite measure belongs to $\mathcal{F}$.

We proceed in stages.

\begin{enumerate}
\item[(a)] First suppose $E$ is a bounded open cube in $\mathbb{R}^d$, with
\[
E=Q_1\times Q_2,
\]
where $Q_1$ and $Q_2$ are open cubes in $\mathbb{R}^{d_1}$ and $\mathbb{R}^{d_2}$, respectively. For each $y$, $\chi_E(x,y)$ is measurable in $x$ and
\[
g(y)=\int_{\mathbb{R}^{d_1}}\chi_E(x,y)\,dx
=
\begin{cases}
|Q_1|,&y\in Q_2,\\
0,&\text{otherwise}.
\end{cases}
\]
Thus $g=|Q_1|\chi_{Q_2}$ is integrable and
\[
\int_{\mathbb{R}^{d_2}}g(y)\,dy=|Q_1||Q_2|=|E|
=\int_{\mathbb{R}^d}\chi_E.
\]
Hence $\chi_E\in\mathcal{F}$.

\item[(b)] Suppose $E$ is a subset of the boundary of some closed cube. Since the boundary of a cube has measure zero in $\mathbb{R}^d$,
\[
\int_{\mathbb{R}^d}\chi_E(x,y)\,dx\,dy=0.
\]
For almost every $y$, the slice $E^y$ has measure zero in $\mathbb{R}^{d_1}$. Hence
\[
g(y)=\int_{\mathbb{R}^{d_1}}\chi_E(x,y)\,dx=0
\]
for a.e. $y$, and therefore $\chi_E\in\mathcal{F}$.

\item[(c)] Suppose $E$ is a finite union of closed cubes whose interiors are disjoint,
\[
E=\bigcup_{k=1}^{K}Q_k.
\]
If $\widetilde Q_k$ denotes the interior of $Q_k$, we may write $\chi_E$ as a linear combination of the $\chi_{\widetilde Q_k}$ and characteristic functions of subsets of the boundaries of the $Q_k$. By (a), (b), and Step 1, $\chi_E\in\mathcal{F}$.

\item[(d)] Suppose $E$ is open and of finite measure. By Theorem 1.4 in Chapter 1,
\[
E=\bigcup_{j=1}^{\infty}Q_j
\]
as a countable union of almost disjoint closed cubes. If
\[
f_k=\sum_{j=1}^{k}\chi_{Q_j},
\]
then $f_k\nearrow\chi_E$. Each $f_k\in\mathcal{F}$ by (c) and Step 1, and $\chi_E$ is integrable because $m(E)<\infty$. Thus Step 2 gives $\chi_E\in\mathcal{F}$.

\item[(e)] Finally, suppose $E$ is a $G_\delta$ of finite measure. Write
\[
E=\bigcap_{k=1}^{\infty}\widetilde O_k
\]
with each $\widetilde O_k$ open. Since $E$ has finite measure, there exists an open set $\widetilde O_0$ of finite measure with $E\subset\widetilde O_0$. Set
\[
O_k=\widetilde O_0\cap\bigcap_{j=1}^{k}\widetilde O_j.
\]
Then $O_k$ are open, of finite measure, $O_k\searrow E$, and therefore $\chi_{O_k}\searrow\chi_E$. By (d) and Step 2, $\chi_E\in\mathcal{F}$.
\end{enumerate}

\paragraph{Step 4.} If $E$ has measure zero, then $\chi_E\in\mathcal{F}$.

Since $E$ is measurable, choose a $G_\delta$ set $G$ with $E\subset G$ and $m(G)=0$. By Step 3, $\chi_G\in\mathcal{F}$, so
\[
\int_{\mathbb{R}^{d_2}}
\left(\int_{\mathbb{R}^{d_1}}\chi_G(x,y)\,dx\right)dy
=\int_{\mathbb{R}^d}\chi_G=0.
\]
Hence
\[
\int_{\mathbb{R}^{d_1}}\chi_G(x,y)\,dx=0
\]
for a.e. $y$. Thus $G^y$ has measure zero for a.e. $y$. Since $E^y\subset G^y$, the same is true for $E^y$, and therefore
\[
\int_{\mathbb{R}^{d_2}}
\left(\int_{\mathbb{R}^{d_1}}\chi_E(x,y)\,dx\right)dy
=0
=\int_{\mathbb{R}^d}\chi_E.
\]
Thus $\chi_E\in\mathcal{F}$.

\paragraph{Step 5.} If $E$ is any measurable subset of $\mathbb{R}^d$ with finite measure, then $\chi_E\in\mathcal{F}$.

There exists a $G_\delta$ set $G$ of finite measure such that $E\subset G$ and $m(G-E)=0$. Since
\[
\chi_E=\chi_G-\chi_{G-E},
\]
Steps 3, 4, and 1 imply $\chi_E\in\mathcal{F}$.

\paragraph{Step 6.} If $f$ is integrable, then $f\in\mathcal{F}$.

Write $f=f^+-f^-$. By Step 1 it is enough to assume $f\geq0$. By Theorem 4.1 in Chapter 1, there is a sequence $\{\varphi_k\}$ of simple functions with $\varphi_k\nearrow f$. Each $\varphi_k$ is a finite linear combination of characteristic functions of sets of finite measure, hence belongs to $\mathcal{F}$ by Steps 5 and 1. Step 2 then gives $f\in\mathcal{F}$ and proves Fubini's theorem.

\subsection{Applications of Fubini's theorem}

\begin{theorem}
Suppose $f(x,y)$ is a non-negative measurable function on $\mathbb{R}^{d_1}\times\mathbb{R}^{d_2}$. Then for almost every $y\in\mathbb{R}^{d_2}$:
\begin{enumerate}
\item[(i)] The slice $f^y$ is measurable on $\mathbb{R}^{d_1}$.
\item[(ii)] The function
\[
y\longmapsto\int_{\mathbb{R}^{d_1}}f^y(x)\,dx
\]
is measurable on $\mathbb{R}^{d_2}$.
\end{enumerate}
Moreover,
\begin{enumerate}
\item[(iii)]
\[
\int_{\mathbb{R}^{d_2}}
\left(\int_{\mathbb{R}^{d_1}}f(x,y)\,dx\right)dy
=\int_{\mathbb{R}^d}f(x,y)\,dx\,dy,
\]
in the extended sense.
\end{enumerate}
\end{theorem}

In practice this result is often used in conjunction with Fubini's theorem. Suppose a measurable function $f$ is given and one wants to compute $\int_{\mathbb{R}^d}f$. One first applies the present theorem to $|f|$. If the iterated integrals of $|f|$ are finite, then $\int|f|<\infty$, and the hypothesis of Fubini's theorem is verified. One may then calculate the integral of $f$ by iterated integration.

\begin{proof}
Consider the truncations
\[
f_k(x,y)=
\begin{cases}
f(x,y),& |(x,y)|<k\text{ and }f(x,y)<k,\\
0,&\text{otherwise}.
\end{cases}
\]
Each $f_k$ is integrable. By part (i) of Fubini's theorem there exists a set $E_k\subset\mathbb{R}^{d_2}$ of measure zero such that the slice $f_k^y$ is measurable for every $y\notin E_k$. Let $E=\bigcup_kE_k$. Then $m(E)=0$, and for $y\notin E$ we have $f_k^y\nearrow f^y$. Hence by monotone convergence,
\[
\int_{\mathbb{R}^{d_1}}f_k(x,y)\,dx
\nearrow
\int_{\mathbb{R}^{d_1}}f(x,y)\,dx.
\]
The left-hand functions are measurable in $y$ by Fubini's theorem, hence the limit is measurable. Another application of the monotone convergence theorem gives
\begin{equation}
\int_{\mathbb{R}^{d_2}}
\left(\int_{\mathbb{R}^{d_1}}f_k(x,y)\,dx\right)dy
\longrightarrow
\int_{\mathbb{R}^{d_2}}
\left(\int_{\mathbb{R}^{d_1}}f(x,y)\,dx\right)dy.
\tag{11}
\end{equation}
By Fubini's theorem,
\begin{equation}
\int_{\mathbb{R}^{d_2}}
\left(\int_{\mathbb{R}^{d_1}}f_k(x,y)\,dx\right)dy
=\int_{\mathbb{R}^d}f_k.
\tag{12}
\end{equation}
Finally, monotone convergence directly on $\mathbb{R}^d$ gives
\begin{equation}
\int_{\mathbb{R}^d}f_k\longrightarrow\int_{\mathbb{R}^d}f.
\tag{13}
\end{equation}
Combining (11)--(13) proves the theorem.
\end{proof}

\begin{corollary}
If $E$ is a measurable set in $\mathbb{R}^{d_1}\times\mathbb{R}^{d_2}$, then for almost every $y\in\mathbb{R}^{d_2}$ the slice
\[
E^y=\{x\in\mathbb{R}^{d_1}:(x,y)\in E\}
\]
is measurable. Moreover $m(E^y)$ is a measurable function of $y$ and
\[
m(E)=\int_{\mathbb{R}^{d_2}}m(E^y)\,dy.
\]
\end{corollary}

This follows immediately by applying the theorem to $\chi_E$. A symmetric result holds for the $x$-slices.

The converse need not hold. If $N\subset\mathbb{R}$ is non-measurable and
\[
E=[0,1]\times N\subset\mathbb{R}\times\mathbb{R},
\]
then
\[
E^y=
\begin{cases}
[0,1],&y\in N,\\
\varnothing,&y\notin N,
\end{cases}
\]
so every $E^y$ is measurable. Yet $E$ is not measurable, because if it were, the symmetric slice statement would imply that $E_x=N$ is measurable for almost every $x\in[0,1]$.

A more striking example is a non-measurable $E\subset[0,1]\times[0,1]$ for which all slices are measurable with $m(E^y)=0$ and $m(E_x)=1$. Its construction uses an ordering $\prec$ of the reals such that $\{x:x\prec y\}$ is countable for every $y$. One sets
\[
E=\{(x,y)\in[0,1]\times[0,1]:x\prec y\}.
\]
Then $E^y$ is countable and has measure zero, while $E_x$ is the complement of a denumerable set in $[0,1]$ and has measure one. If $E$ were measurable, this would contradict the preceding corollary.

\begin{proposition}
If $E=E_1\times E_2$ is a measurable subset of $\mathbb{R}^d$, and $m^*(E_2)>0$, then $E_1$ is measurable.
\end{proposition}

\begin{proof}
By the preceding corollary, for a.e. $y\in\mathbb{R}^{d_2}$ the slice
\[
(\chi_{E_1\times E_2})^y(x)=\chi_{E_1}(x)\chi_{E_2}(y)
\]
is measurable in $x$. We claim that there is some $y\in E_2$ for which this slice is measurable. For such a $y$, the slice equals $\chi_{E_1}$, so $E_1$ is measurable.

Let $F$ denote the set of $y$ for which the slice $E^y$ is measurable. Then $m(F^c)=0$. Since
\[
E_2=(E_2\cap F)\cup(E_2\cap F^c),
\]
we have
\[
0<m^*(E_2)
\leq m^*(E_2\cap F)+m^*(E_2\cap F^c)
=m^*(E_2\cap F).
\]
Thus $E_2\cap F\neq\varnothing$, proving the claim.
\end{proof}

\begin{lemma}
If $E_1\subset\mathbb{R}^{d_1}$ and $E_2\subset\mathbb{R}^{d_2}$, then
\[
m^*(E_1\times E_2)\leq m^*(E_1)m^*(E_2),
\]
with the understanding that if one of the sets $E_j$ has exterior measure zero, then $m^*(E_1\times E_2)=0$.
\end{lemma}

\begin{proof}
Let $\epsilon>0$. Choose cubes $\{Q_k\}_{k=1}^{\infty}$ in $\mathbb{R}^{d_1}$ and $\{Q'_\ell\}_{\ell=1}^{\infty}$ in $\mathbb{R}^{d_2}$ such that
\[
E_1\subset\bigcup_{k=1}^{\infty}Q_k,
\qquad
E_2\subset\bigcup_{\ell=1}^{\infty}Q'_\ell,
\]
and
\[
\sum_{k=1}^{\infty}|Q_k|\leq m^*(E_1)+\epsilon,
\qquad
\sum_{\ell=1}^{\infty}|Q'_\ell|\leq m^*(E_2)+\epsilon.
\]
Since
\[
E_1\times E_2\subset\bigcup_{k,\ell=1}^{\infty}(Q_k\times Q'_\ell),
\]
sub-additivity gives
\begin{align*}
m^*(E_1\times E_2)
&\leq\sum_{k,\ell=1}^{\infty}|Q_k\times Q'_\ell|\\
&=\left(\sum_{k=1}^{\infty}|Q_k|\right)
  \left(\sum_{\ell=1}^{\infty}|Q'_\ell|\right)\\
&\leq(m^*(E_1)+\epsilon)(m^*(E_2)+\epsilon).
\end{align*}
If neither exterior measure is zero, let $\epsilon\to0$.

If, for example, $m^*(E_1)=0$, let
\[
E_2^j=E_2\cap\{y\in\mathbb{R}^{d_2}:|y|\leq j\}.
\]
The preceding argument gives $m^*(E_1\times E_2^j)=0$ for each $j$. Since $E_1\times E_2^j\nearrow E_1\times E_2$, countable sub-additivity yields $m^*(E_1\times E_2)=0$.
\end{proof}

\begin{proposition}
Suppose $E_1$ and $E_2$ are measurable subsets of $\mathbb{R}^{d_1}$ and $\mathbb{R}^{d_2}$, respectively. Then $E=E_1\times E_2$ is measurable in $\mathbb{R}^d$. Moreover,
\[
m(E)=m(E_1)m(E_2),
\]
with the understanding that if one of the sets $E_j$ has measure zero, then $m(E)=0$.
\end{proposition}

\begin{proof}
It suffices to prove measurability; the formula for the measure then follows from the slice formula. For $j=1,2$, choose a $G_\delta$ set $G_j\supset E_j$ such that $m^*(G_j-E_j)=0$. Then $G=G_1\times G_2$ is measurable, and
\[
(G_1\times G_2)-(E_1\times E_2)
\subset
((G_1-E_1)\times G_2)\cup(G_1\times(G_2-E_2)).
\]
The lemma implies $m^*(G-E)=0$. Hence $E$ is measurable.
\end{proof}

\begin{corollary}
Suppose $f$ is a measurable function on $\mathbb{R}^{d_1}$. Then the function
\[
\widetilde f(x,y)=f(x)
\]
is measurable on $\mathbb{R}^{d_1}\times\mathbb{R}^{d_2}$.
\end{corollary}

\begin{proof}
Assume $f$ is real-valued. If $a\in\mathbb{R}$ and
\[
E_1=\{x\in\mathbb{R}^{d_1}:f(x)<a\},
\]
then $E_1$ is measurable. Since
\[
\{(x,y):\widetilde f(x,y)<a\}=E_1\times\mathbb{R}^{d_2},
\]
the previous proposition shows that this inverse image is measurable for every $a$.
\end{proof}

We now return to the interpretation of the integral as the ``area'' under the graph of a function.

\begin{corollary}
Suppose $f(x)$ is a non-negative function on $\mathbb{R}^d$, and let
\[
A=\{(x,y)\in\mathbb{R}^d\times\mathbb{R}:0\leq y\leq f(x)\}.
\]
Then:
\begin{enumerate}
\item[(i)] $f$ is measurable on $\mathbb{R}^d$ if and only if $A$ is measurable in $\mathbb{R}^{d+1}$.
\item[(ii)] If these conditions hold, then
\[
\int_{\mathbb{R}^d}f(x)\,dx=m(A).
\]
\end{enumerate}
\end{corollary}

\begin{proof}
If $f$ is measurable, the preceding corollary implies that
\[
F(x,y)=y-f(x)
\]
is measurable on $\mathbb{R}^{d+1}$, and therefore
\[
A=\{y\geq0\}\cap\{F\leq0\}
\]
is measurable.

Conversely, suppose $A$ is measurable. For each $x$, the slice
\[
A_x=\{y\in\mathbb{R}:(x,y)\in A\}=[0,f(x)].
\]
The slice formula gives the measurability of $m(A_x)=f(x)$ and
\[
m(A)
=\int\chi_A(x,y)\,dx\,dy
=\int_{\mathbb{R}^d}m(A_x)\,dx
=\int_{\mathbb{R}^d}f(x)\,dx.
\]
\end{proof}

We conclude this section with a useful result.

\begin{proposition}
If $f$ is a measurable function on $\mathbb{R}^d$, then the function
\[
\widetilde f(x,y)=f(x-y)
\]
is measurable on $\mathbb{R}^d\times\mathbb{R}^d$.
\end{proposition}

\begin{proof}
For $a\in\mathbb{R}$, let
\[
E=\{z\in\mathbb{R}^d:f(z)<a\}.
\]
It suffices to prove that
\[
\widetilde E=\{(x,y):x-y\in E\}
\]
is measurable whenever $E$ is measurable.

If $O$ is open, then $\widetilde O$ is open. Hence if $E$ is a $G_\delta$, then so is $\widetilde E$. We next show that sets of measure zero lift to sets of measure zero. Let
\[
\widetilde E_k=\widetilde E\cap B_k,
\qquad B_k=\{(x,y):|y|<k\}.
\]
For an open set $O\subset\mathbb{R}^d$,
\[
\chi_{\widetilde O\cap B_k}(x,y)=\chi_O(x-y)\chi_{\{|y|<k\}}(y),
\]
so by Fubini and translation-invariance,
\begin{align*}
m(\widetilde O\cap B_k)
&=\int\chi_O(x-y)\chi_{\{|y|<k\}}(y)\,dy\,dx\\
&=m(O)m(\{|y|<k\}).
\end{align*}
If $m(E)=0$, choose open $O_n\supset E$ with $m(O_n)\to0$. Then $\widetilde E_k\subset\widetilde O_n\cap B_k$ and hence $m(\widetilde E_k)=0$ for each fixed $k$. Therefore $m(\widetilde E)=0$.

Finally, any measurable set can be written as the difference of a $G_\delta$ set and a set of measure zero. The two cases above prove the result.
\end{proof}

\section{A Fourier inversion formula${}^*$}

The question of the inversion of the Fourier transform encompasses in effect the problem at the origin of Fourier analysis. This issue involves establishing the validity of the inversion formula for a function $f$ in terms of its Fourier transform $\widehat f$, that is,
\begin{equation}
\widehat f(\xi)=\int_{\mathbb{R}^d}f(x)e^{-2\pi i x\cdot\xi}\,dx,
\tag{14}
\end{equation}
and
\begin{equation}
f(x)=\int_{\mathbb{R}^d}\widehat f(\xi)e^{2\pi i x\cdot\xi}\,d\xi.
\tag{15}
\end{equation}

We have already encountered this problem in Book I in the rudimentary case when both $f$ and $\widehat f$ were continuous and had rapid (or moderate) decrease at infinity. In Book II we also considered the question in the one-dimensional setting, seen from the viewpoint of complex analysis. The most elegant and useful formulations of Fourier inversion are in terms of the $L^2$ theory, or in its greatest generality stated in the language of distributions. We shall take up these matters systematically later. It will, nevertheless, be enlightening to see what our present knowledge teaches us about the problem. We present a variant of the inversion formula appropriate for $L^1$, one that is both simple and adequate in many circumstances.

To begin with, we need to know what can be said about the Fourier transform of an arbitrary function in $L^1(\mathbb{R}^d)$.

\begin{proposition}
Suppose $f\in L^1(\mathbb{R}^d)$. Then $\widehat f$ defined by (14) is continuous and bounded on $\mathbb{R}^d$.
\end{proposition}

Indeed, since
\[
|f(x)e^{-2\pi i x\cdot\xi}|=|f(x)|,
\]
the integral representing $\widehat f$ converges for each $\xi$, and
\[
\sup_{\xi\in\mathbb{R}^d}|\widehat f(\xi)|
\leq\int_{\mathbb{R}^d}|f(x)|\,dx
=\|f\|.
\]
To verify continuity, note that for every $x$,
\[
f(x)e^{-2\pi i x\cdot\xi}
\longrightarrow
f(x)e^{-2\pi i x\cdot\xi_0}
\qquad\text{as }\xi\to\xi_0,
\]
and apply the dominated convergence theorem.

One can assert a little more than boundedness: $\widehat f(\xi)\to0$ as $|\xi|\to\infty$, but in general not much more can be said about the decrease at infinity. Consequently, for general $f\in L^1(\mathbb{R}^d)$ the function $\widehat f$ need not belong to $L^1(\mathbb{R}^d)$, and formula (15) becomes problematical. The following theorem evades this difficulty.

\begin{theorem}
Suppose $f\in L^1(\mathbb{R}^d)$ and assume also that $\widehat f\in L^1(\mathbb{R}^d)$. Then the inversion formula (15) holds for almost every $x$.
\end{theorem}

An immediate corollary is the uniqueness of the Fourier transform on $L^1$.

\begin{corollary}
Suppose $\widehat f(\xi)=0$ for all $\xi$. Then $f=0$ a.e.
\end{corollary}

The proof of the theorem requires only that we adapt the earlier arguments carried out for Schwartz functions in Chapter 5 of Book I to the present context. We begin with the ``multiplication formula.''

\begin{lemma}
Suppose $f$ and $g$ belong to $L^1(\mathbb{R}^d)$. Then
\[
\int_{\mathbb{R}^d}\widehat f(\xi)g(\xi)\,d\xi
=\int_{\mathbb{R}^d}f(y)\widehat g(y)\,dy.
\]
\end{lemma}

\begin{proof}
Both integrals converge by the preceding proposition. Consider
\[
F(\xi,y)=g(\xi)f(y)e^{-2\pi i\xi\cdot y},
\qquad (\xi,y)\in\mathbb{R}^d\times\mathbb{R}^d.
\]
It is measurable by Corollary 3.7, and
\[
\int_{\mathbb{R}^d}\int_{\mathbb{R}^d}|F(\xi,y)|\,d\xi\,dy
=
\left(\int_{\mathbb{R}^d}|g(\xi)|\,d\xi\right)
\left(\int_{\mathbb{R}^d}|f(y)|\,dy\right)<\infty.
\]
Hence Fubini's theorem applies. Evaluating the double integral first in $\xi$ gives the right-hand side, while evaluating first in $y$ gives the left-hand side.
\end{proof}

Next consider the modulated Gaussian
\[
g(\xi)=e^{-\pi\delta|\xi|^2}e^{2\pi i x\cdot\xi},
\]
where $\delta>0$ and $x\in\mathbb{R}^d$ are fixed. An elementary calculation gives
\[
\widehat g(y)
=\int_{\mathbb{R}^d}e^{-\pi\delta|\xi|^2}e^{2\pi i(x-y)\cdot\xi}\,d\xi
=\delta^{-d/2}e^{-\pi|x-y|^2/\delta}.
\]
We abbreviate the last expression as $K_\delta(x-y)$. The kernel $K_\delta$ satisfies
\[
\int_{\mathbb{R}^d}K_\delta(y)\,dy=1,
\]
and for every $\eta>0$,
\[
\int_{|y|>\eta}K_\delta(y)\,dy\to0
\qquad\text{as }\delta\to0.
\]
Applying the lemma gives
\begin{equation}
\int_{\mathbb{R}^d}
\widehat f(\xi)e^{-\pi\delta|\xi|^2}e^{2\pi i x\cdot\xi}\,d\xi
=
\int_{\mathbb{R}^d}f(y)K_\delta(x-y)\,dy.
\tag{16}
\end{equation}

Since $\widehat f\in L^1(\mathbb{R}^d)$, the dominated convergence theorem shows that the left-hand side of (16) converges, for each $x$, to
\[
\int_{\mathbb{R}^d}\widehat f(\xi)e^{2\pi i x\cdot\xi}\,d\xi
\]
as $\delta\to0$.

For the right-hand side, make the successive changes of variables $y\mapsto y+x$ and $y\mapsto-y$, using translation- and reflection-invariance. The right-hand side becomes
\[
\int_{\mathbb{R}^d}f(x-y)K_\delta(y)\,dy.
\]
We prove that this converges to $f$ in $L^1$ as $\delta\to0$. Set
\[
\Delta_\delta(x)
=\int_{\mathbb{R}^d}f(x-y)K_\delta(y)\,dy-f(x).
\]
Using $\int K_\delta=1$,
\[
\Delta_\delta(x)
=\int_{\mathbb{R}^d}(f(x-y)-f(x))K_\delta(y)\,dy,
\]
and hence
\[
|\Delta_\delta(x)|
\leq\int_{\mathbb{R}^d}|f(x-y)-f(x)|K_\delta(y)\,dy.
\]
By Fubini's theorem,
\[
\|\Delta_\delta\|
\leq\int_{\mathbb{R}^d}\|f_y-f\|K_\delta(y)\,dy,
\qquad f_y(x)=f(x-y).
\]
Given $\epsilon>0$, Proposition 2.5 gives $\eta>0$ such that
\[
\|f_y-f\|<\epsilon
\qquad\text{whenever }|y|<\eta.
\]
Therefore
\begin{align*}
\|\Delta_\delta\|
&\leq\epsilon
 +\int_{|y|>\eta}\|f_y-f\|K_\delta(y)\,dy\\
&\leq\epsilon
 +2\|f\|\int_{|y|>\eta}K_\delta(y)\,dy.
\end{align*}
The first inequality uses $\int K_\delta=1$, and the second uses
\[
\|f_y-f\|\leq\|f_y\|+\|f\|=2\|f\|.
\]
By the second property of $K_\delta$, the last expression is at most $2\epsilon$ for all sufficiently small $\delta$. Thus the right-hand side of (16) converges to $f$ in $L^1$ as $\delta\to0$.

By Corollary 2.3, a subsequence converges to $f(x)$ almost everywhere. Along the same subsequence, the left-hand side of (16) converges pointwise to the expression in (15). Hence (15) holds for almost every $x$, proving the theorem.

An immediate consequence of the theorem and the first proposition is that if $\widehat f$ were in $L^1$, then $f$ could be modified on a set of measure zero to become continuous everywhere. This is impossible for a general $f\in L^1(\mathbb{R}^d)$.

% End of Chapter 2 body. Exercises and Problems are intentionally omitted.

% Stein & Shakarchi, Real Analysis (Princeton Lectures in Analysis III)
% Chapter 3: Differentiation and Integration
% Transcribed from the user-provided PDF.
% Exercises and Problems are intentionally omitted.
% This file is intended to be \input{} into a book/report document.
% Suggested packages in the parent document:
%   amsmath, amssymb, amsthm, mathtools, graphicx
% Suggested theorem declarations:
%   \newtheorem{theorem}{Theorem}[section]
%   \newtheorem{lemma}[theorem]{Lemma}
%   \newtheorem{corollary}[theorem]{Corollary}
%   \newtheorem{proposition}[theorem]{Proposition}
%
% Figures from the source are represented by compile-safe placeholders.

\chapter{Differentiation and Integration}

\begin{quote}
\textbf{The Maximal Problem:}\\
The problem is most easily grasped when stated in the language of cricket, or any other game in which a player compiles a series of scores of which an average is recorded.

\hfill G. H. Hardy and J. E. Littlewood, 1930
\end{quote}

That differentiation and integration are inverse operations was already understood early in the study of the calculus. Here we want to reexamine this basic idea in the framework of the general theory studied in the previous chapters. Our objective is the formulation and proof of the fundamental theorem of the calculus in this setting, and the development of some of the concepts that occur. We shall try to achieve this by answering two questions, each expressing one of the ways of representing the reciprocity between differentiation and integration.

The first problem involved may be stated as follows.
\begin{itemize}
\item Suppose $f$ is integrable on $[a,b]$ and $F$ is its indefinite integral
\[
F(x)=\int_a^x f(y)\,dy.
\]
Does this imply that $F$ is differentiable (at least for almost every $x$), and that $F'=f$?
\end{itemize}
We shall see that the affirmative answer to this question depends on ideas that have broad application and are not limited to the one-dimensional situation.

For the second question we reverse the order of differentiation and integration.
\begin{itemize}
\item What conditions on a function $F$ on $[a,b]$ guarantee that $F'(x)$ exists (for a.e. $x$), that this function is integrable, and that moreover
\[
F(b)-F(a)=\int_a^b F'(x)\,dx?
\]
\end{itemize}
While this problem will be examined from a narrower perspective than the first, the issues it raises are deep and the consequences entailed are far-reaching. In particular, we shall find that this question is connected to the problem of rectifiability of curves, and as an illustration of this link, we shall establish the general isoperimetric inequality in the plane.

\section{Differentiation of the integral}

We begin with the first problem, that is, the study of differentiation of the integral. If $f$ is given on $[a,b]$ and integrable on that interval, we let
\[
F(x)=\int_a^x f(y)\,dy,\qquad a\le x\le b.
\]
To deal with $F'(x)$, we recall the definition of the derivative as the limit of the quotient
\[
\frac{F(x+h)-F(x)}{h}
\]
when $h$ tends to $0$. We note that this quotient takes the form (say in the case $h>0$)
\[
\frac1h\int_x^{x+h}f(y)\,dy=\frac1{|I|}\int_I f(y)\,dy,
\]
where $I=(x,x+h)$ and $|I|$ denotes its length. The expression is the average value of $f$ over $I$, and we might expect that, as $|I|\to0$, these averages tend to $f(x)$. Thus we ask whether
\[
\lim_{\substack{|I|\to0\\x\in I}}\frac1{|I|}\int_I f(y)\,dy=f(x)
\]
holds for suitable points $x$.

In higher dimensions we can pose a similar question, taking averages over balls $B$ containing $x$, with their volume $m(B)$ replacing $|I|$:
\[
\lim_{\substack{m(B)\to0\\x\in B}}\frac1{m(B)}\int_B f(y)\,dy=f(x)\qquad\text{for a.e. }x?
\]
The limit is taken as the volume of open balls $B$ containing $x$ goes to $0$. We shall refer to this as the \emph{averaging problem}. If $B$ has radius $r$ in $\mathbb R^d$, then $m(B)=v_dr^d$, where $v_d$ is the measure of the unit ball.

If $f$ is continuous at $x$, the limit is immediate. Given $\epsilon>0$, choose $\delta>0$ so that $|f(x)-f(y)|<\epsilon$ whenever $|x-y|<\delta$. Since
\[
f(x)-\frac1{m(B)}\int_B f(y)\,dy
=\frac1{m(B)}\int_B(f(x)-f(y))\,dy,
\]
we find, whenever $B$ has radius $<\delta/2$ and contains $x$, that
\[
\left|f(x)-\frac1{m(B)}\int_B f(y)\,dy\right|
\le \frac1{m(B)}\int_B|f(x)-f(y)|\,dy<\epsilon.
\]
The averaging problem has an affirmative answer, but its proof depends on quantitative estimates for maximal averages.

\subsection{The Hardy--Littlewood maximal function}

If $f$ is integrable on $\mathbb R^d$, we define its \emph{maximal function} $f^*$ by
\[
f^*(x)=\sup_{x\in B}\frac1{m(B)}\int_B|f(y)|\,dy,
\qquad x\in\mathbb R^d,
\]
where the supremum is taken over all balls containing $x$.

\begin{theorem}[Hardy--Littlewood maximal theorem]
Suppose $f$ is integrable on $\mathbb R^d$. Then:
\begin{enumerate}
\item[(i)] $f^*$ is measurable.
\item[(ii)] $f^*(x)<\infty$ for a.e. $x$.
\item[(iii)] $f^*$ satisfies
\begin{equation}
 m\bigl(\{x\in\mathbb R^d:f^*(x)>\alpha\}\bigr)
 \le \frac{A}{\alpha}\|f\|_{L^1(\mathbb R^d)},
 \qquad \alpha>0,
 \tag{1}
\end{equation}
where $A=3^d$ and
\[
\|f\|_{L^1(\mathbb R^d)}=\int_{\mathbb R^d}|f(x)|\,dx.
\]
\end{enumerate}
\end{theorem}

An inequality of the type (1) is called a \emph{weak-type inequality}. It is weaker than the corresponding $L^1$ estimate, as can be seen from Tchebychev's inequality
\[
m(\{x:|g(x)|>\alpha\})\le \frac1\alpha\|g\|_{L^1(\mathbb R^d)}.
\]
The exact value of $A$ is unimportant here; what matters is that it is independent of $\alpha$ and $f$.

The measurability of $f^*$ is straightforward. The set
\[
E_\alpha=\{x\in\mathbb R^d:f^*(x)>\alpha\}
\]
is open: if $x\in E_\alpha$, there is a ball $B$ containing $x$ such that
\[
\frac1{m(B)}\int_B|f(y)|\,dy>\alpha,
\]
and every point sufficiently close to $x$ also lies in $B$.

Part (ii) follows from (iii), since
\[
\{x:f^*(x)=\infty\}\subset\{x:f^*(x)>\alpha\}
\]
for every $\alpha>0$, and the right-hand side has measure tending to $0$ as $\alpha\to\infty$.

The proof of (iii) relies on an elementary Vitali-type covering argument.

\begin{lemma}
Suppose $\mathcal B=\{B_1,B_2,\ldots,B_N\}$ is a finite collection of open balls in $\mathbb R^d$. Then there exists a disjoint subcollection
$B_{i_1},\ldots,B_{i_k}$ such that
\[
m\left(\bigcup_{\ell=1}^N B_\ell\right)
\le 3^d\sum_{j=1}^k m(B_{i_j}).
\]
\end{lemma}

\begin{proof}
Suppose $B$ and $B'$ intersect and the radius of $B'$ is not greater than the radius of $B$. Then $B'$ is contained in the ball $\widetilde B$ concentric with $B$ and with three times its radius.

Pick a ball $B_{i_1}$ of maximal radius and delete it together with every ball that intersects it. Repeat the procedure with the remaining collection. After at most $N$ steps we obtain disjoint balls $B_{i_1},\ldots,B_{i_k}$. Every deleted ball is contained in the triple dilation of the selected ball responsible for its deletion. Hence, writing $\widetilde B_{i_j}$ for the concentric triple dilation,
\[
\bigcup_{\ell=1}^N B_\ell\subset \bigcup_{j=1}^k\widetilde B_{i_j}.
\]
Therefore
\[
m\left(\bigcup_{\ell=1}^N B_\ell\right)
\le \sum_{j=1}^k m(\widetilde B_{i_j})
=3^d\sum_{j=1}^k m(B_{i_j}).
\]
\end{proof}

\begin{figure}[ht]
\centering
\fbox{\parbox{0.72\textwidth}{\centering Source Figure 1 omitted: intersecting balls $B,B'$ and the concentric triple dilation $\widetilde B$.}}
\caption{The balls $B$ and $\widetilde B$.}
\end{figure}

\begin{proof}[Proof of Theorem 1.1(iii)]
Let $E_\alpha=\{x:f^*(x)>\alpha\}$. For every $x\in E_\alpha$ choose a ball $B_x$ containing $x$ and satisfying
\[
\frac1{m(B_x)}\int_{B_x}|f(y)|\,dy>\alpha,
\]
so
\begin{equation}
m(B_x)<\frac1\alpha\int_{B_x}|f(y)|\,dy.
\tag{2}
\end{equation}
Fix a compact subset $K\subset E_\alpha$. A finite subcollection $B_1,\ldots,B_N$ covers $K$. By the lemma there are disjoint $B_{i_1},\ldots,B_{i_k}$ with
\begin{equation}
m\left(\bigcup_{\ell=1}^N B_\ell\right)
\le3^d\sum_{j=1}^k m(B_{i_j}).
\tag{3}
\end{equation}
Using (2), disjointness, and (3),
\begin{align*}
m(K)
&\le3^d\sum_{j=1}^k m(B_{i_j})\\
&\le \frac{3^d}{\alpha}\sum_{j=1}^k\int_{B_{i_j}}|f(y)|\,dy\\
&\le \frac{3^d}{\alpha}\int_{\mathbb R^d}|f(y)|\,dy.
\end{align*}
Taking the supremum over compact $K\subset E_\alpha$ proves (1).
\end{proof}

\subsection{The Lebesgue differentiation theorem}

The maximal estimate leads to the solution of the averaging problem.

\begin{theorem}[Lebesgue differentiation theorem]
If $f$ is integrable on $\mathbb R^d$, then
\begin{equation}
\lim_{\substack{m(B)\to0\\x\in B}}
\frac1{m(B)}\int_B f(y)\,dy=f(x)
\qquad\text{for a.e. }x.
\tag{4}
\end{equation}
\end{theorem}

\begin{proof}
For $\alpha>0$ let
\[
E_\alpha=
\left\{x:\limsup_{\substack{m(B)\to0\\x\in B}}
\left|\frac1{m(B)}\int_Bf(y)\,dy-f(x)\right|>2\alpha\right\}.
\]
It suffices to show $m(E_\alpha)=0$ for every $\alpha>0$, because then $\bigcup_{n\ge1}E_{1/n}$ has measure zero.

Fix $\epsilon>0$. By the approximation theorem in Chapter 2, choose a continuous compactly supported $g$ such that
\[
\|f-g\|_{L^1(\mathbb R^d)}<\epsilon.
\]
Continuity gives
\[
\lim_{\substack{m(B)\to0\\x\in B}}\frac1{m(B)}\int_Bg(y)\,dy=g(x).
\]
Writing
\begin{align*}
\frac1{m(B)}\int_Bf(y)\,dy-f(x)
&=\frac1{m(B)}\int_B(f(y)-g(y))\,dy\\
&\quad+\frac1{m(B)}\int_Bg(y)\,dy-g(x)+g(x)-f(x),
\end{align*}
we obtain
\[
\limsup_{\substack{m(B)\to0\\x\in B}}
\left|\frac1{m(B)}\int_Bf(y)\,dy-f(x)\right|
\le (f-g)^*(x)+|g(x)-f(x)|.
\]
Let
\[
F_\alpha=\{x:(f-g)^*(x)>\alpha\},\qquad
G_\alpha=\{x:|f(x)-g(x)|>\alpha\}.
\]
Then $E_\alpha\subset F_\alpha\cup G_\alpha$. Tchebychev's inequality and the weak-type bound yield
\[
m(G_\alpha)\le\frac1\alpha\|f-g\|_1,
\qquad
m(F_\alpha)\le\frac A\alpha\|f-g\|_1.
\]
Hence
\[
m(E_\alpha)\le\frac{A+1}{\alpha}\epsilon.
\]
Since $\epsilon$ is arbitrary, $m(E_\alpha)=0$.
\end{proof}

Applying the theorem to $|f|$ gives $f^*(x)\ge|f(x)|$ for a.e. $x$.

A measurable function $f$ on $\mathbb R^d$ is \emph{locally integrable} if $f\chi_B$ is integrable for every ball $B$. We write $f\in L^1_{\mathrm{loc}}(\mathbb R^d)$.

\begin{theorem}
If $f\in L^1_{\mathrm{loc}}(\mathbb R^d)$, then
\[
\lim_{\substack{m(B)\to0\\x\in B}}
\frac1{m(B)}\int_Bf(y)\,dy=f(x)
\qquad\text{for a.e. }x.
\]
\end{theorem}

If $E$ is measurable, a point $x\in\mathbb R^d$ is a \emph{point of Lebesgue density} of $E$ if
\[
\lim_{\substack{m(B)\to0\\x\in B}}
\frac{m(B\cap E)}{m(B)}=1.
\]

\begin{corollary}
Suppose $E$ is a measurable subset of $\mathbb R^d$. Then:
\begin{enumerate}
\item[(i)] almost every $x\in E$ is a point of density of $E$;
\item[(ii)] almost every $x\notin E$ is not a point of density of $E$.
\end{enumerate}
\end{corollary}

If $f$ is locally integrable, the \emph{Lebesgue set} of $f$ consists of all points $x$ for which $f(x)$ is finite and
\[
\lim_{\substack{m(B)\to0\\x\in B}}
\frac1{m(B)}\int_B|f(y)-f(x)|\,dy=0.
\]
Every point of continuity belongs to the Lebesgue set.

\begin{corollary}
If $f$ is locally integrable on $\mathbb R^d$, then almost every point belongs to the Lebesgue set of $f$.
\end{corollary}

\begin{proof}
For every rational $r$, apply the differentiation theorem to $|f-r|$. There is a null set $E_r$ such that, for $x\notin E_r$,
\[
\lim_{\substack{m(B)\to0\\x\in B}}
\frac1{m(B)}\int_B|f(y)-r|\,dy=|f(x)-r|.
\]
Let $E=\bigcup_{r\in\mathbb Q}E_r$, still a null set. If $x\notin E$ and $f(x)$ is finite, choose $r\in\mathbb Q$ with $|f(x)-r|<\epsilon$. Then
\[
\frac1{m(B)}\int_B|f(y)-f(x)|\,dy
\le \frac1{m(B)}\int_B|f(y)-r|\,dy+|f(x)-r|.
\]
Taking the limsup gives at most $2\epsilon$, and $\epsilon$ is arbitrary.
\end{proof}

\paragraph{Remark.}
Elements of $L^1(\mathbb R^d)$ are equivalence classes modulo equality almost everywhere. The set of points where the averages in (4) converge to a limit is independent of the representative, because equivalent functions have the same integrals over every ball. The Lebesgue set itself, however, depends on the representative chosen.

A collection $\{U_\alpha\}$ is said to \emph{shrink regularly to $x$} (or to have bounded eccentricity at $x$) if there is a constant $c>0$ such that for every $U_\alpha$ there is a ball $B$ with
\[
x\in B,\qquad U_\alpha\subset B,\qquad m(U_\alpha)\ge c\,m(B).
\]
Open cubes containing $x$ shrink regularly to $x$, whereas the family of all rectangles in $\mathbb R^d$, $d\ge2$, does not.

\begin{corollary}
Suppose $f$ is locally integrable on $\mathbb R^d$. If $\{U_\alpha\}$ shrinks regularly to $x$, then
\[
\lim_{\substack{m(U_\alpha)\to0\\x\in U_\alpha}}
\frac1{m(U_\alpha)}\int_{U_\alpha}f(y)\,dy=f(x)
\]
for every point $x$ in the Lebesgue set of $f$.
\end{corollary}

Indeed, if $U_\alpha\subset B$, $x\in B$, and $m(U_\alpha)\ge c m(B)$, then
\[
\frac1{m(U_\alpha)}\int_{U_\alpha}|f(y)-f(x)|\,dy
\le \frac1{c\,m(B)}\int_B|f(y)-f(x)|\,dy.
\]

\section{Good kernels and approximations to the identity}

We now turn to averages of functions given as convolutions,
\[
(f*K_\delta)(x)=\int_{\mathbb R^d}f(x-y)K_\delta(y)\,dy.
\]
Here $f$ is fixed while the functions $K_\delta$ vary over a family of kernels. In our initial consideration these were called \emph{good kernels} if they are integrable and satisfy, for $\delta>0$,
\begin{enumerate}
\item[(i)]
\[
\int_{\mathbb R^d}K_\delta(x)\,dx=1;
\]
\item[(ii)]
\[
\int_{\mathbb R^d}|K_\delta(x)|\,dx\le A;
\]
\item[(iii)] for every $\eta>0$,
\[
\int_{|x|\ge\eta}|K_\delta(x)|\,dx\longrightarrow0
\qquad\text{as }\delta\to0.
\]
\end{enumerate}
Here $A$ is independent of $\delta$.

Whenever $f$ is bounded, good kernels satisfy $(f*K_\delta)(x)\to f(x)$ at every point of continuity of $f$. To obtain the same conclusion at every point in the Lebesgue set of $f$, we strengthen the assumptions. We say that $\{K_\delta\}_{\delta>0}$ is an \emph{approximation to the identity} if the kernels are integrable, satisfy (i), and in place of (ii) and (iii) satisfy
\begin{enumerate}
\item[(ii$'$)]
\[
|K_\delta(x)|\le A\delta^{-d}
\qquad\text{for all }\delta>0;
\]
\item[(iii$'$)]
\[
|K_\delta(x)|\le A\frac{\delta}{|x|^{d+1}}
\qquad\text{for all }\delta>0\text{ and }x\in\mathbb R^d.
\]
\end{enumerate}
These conditions imply those for good kernels. Indeed, one uses
\begin{equation}
\int_{|x|\ge\epsilon}\frac{dx}{|x|^{d+1}}\le\frac C\epsilon,
\tag{5}
\end{equation}
for some $C>0$. Splitting at $|x|=\delta$ gives
\begin{align*}
\int_{\mathbb R^d}|K_\delta(x)|\,dx
&=\int_{|x|<\delta}|K_\delta(x)|\,dx
 +\int_{|x|\ge\delta}|K_\delta(x)|\,dx\\
&\le A\int_{|x|<\delta}\frac{dx}{\delta^d}
 +A\delta\int_{|x|\ge\delta}\frac{dx}{|x|^{d+1}}\\
&\le A'+A''<\infty.
\end{align*}
Also, for every $\eta>0$,
\[
\int_{|x|\ge\eta}|K_\delta(x)|\,dx
\le A\delta\int_{|x|\ge\eta}\frac{dx}{|x|^{d+1}}
\le A'\frac\delta\eta\to0.
\]

The term ``approximation to the identity'' comes from the fact that $f\mapsto f*K_\delta$ converges to the identity map as $\delta\to0$ in several senses. Heuristically, a family concentrated in $|x|<\delta$ and of total mass one tends to the unit mass at the origin, or Dirac delta ``function'' $D$, formally described by
\[
D(x)=
\begin{cases}
\infty,&x=0,\\
0,&x\ne0,
\end{cases}
\qquad\text{and}\qquad
\int D(x)\,dx=1.
\]
Thus one writes loosely $K_\delta\to D$ and $(f*D)(x)=f(x)$. This heuristic can be made precise either through Lebesgue--Stieltjes measures (Chapter 6) or through distributions (Book IV).

\begin{figure}[ht]
\centering
\fbox{\parbox{0.72\textwidth}{\centering Source Figure 2 omitted: a typical approximation to the identity, supported on $[-\delta,\delta]$ with height $1/(2\delta)$.}}
\caption{An approximation to the identity.}
\end{figure}

We list several examples.

\paragraph{Example 1.}
Suppose $\varphi$ is non-negative and bounded on $\mathbb R^d$, supported on the unit ball, and
\[
\int_{\mathbb R^d}\varphi=1.
\]
Then
\[
K_\delta(x)=\delta^{-d}\varphi(\delta^{-1}x)
\]
defines an approximation to the identity.

\paragraph{Example 2.}
The Poisson kernel for the upper half-plane is
\[
P_y(x)=\frac1\pi\frac{y}{x^2+y^2},\qquad x\in\mathbb R,
\]
with parameter $\delta=y>0$.

\paragraph{Example 3.}
The heat kernel in $\mathbb R^d$ is
\[
H_t(x)=\frac1{(4\pi t)^{d/2}}e^{-|x|^2/(4t)},
\]
with $t>0$ and $\delta=t^{1/2}$.

\paragraph{Example 4.}
The Poisson kernel for the disc, regarded as a function on $\mathbb R$ supported on $[-\pi,\pi]$, is
\[
\frac1{2\pi}P_r(x)=
\begin{cases}
\displaystyle \frac1{2\pi}\frac{1-r^2}{1-2r\cos x+r^2},& |x|\le\pi,\\[6pt]
0,&|x|>\pi,
\end{cases}
\]
where $0<r<1$ and $\delta=1-r$.

\paragraph{Example 5.}
The Fej\'er kernel, similarly extended by zero outside $[-\pi,\pi]$, is
\[
\frac1{2\pi}F_N(x)=
\begin{cases}
\displaystyle \frac1{2\pi N}\frac{\sin^2(Nx/2)}{\sin^2(x/2)},&|x|\le\pi,\\[6pt]
0,&|x|>\pi,
\end{cases}
\]
with $\delta=1/N$.

\begin{theorem}
If $\{K_\delta\}_{\delta>0}$ is an approximation to the identity and $f$ is integrable on $\mathbb R^d$, then
\[
(f*K_\delta)(x)\longrightarrow f(x)\qquad(\delta\to0)
\]
for every $x$ in the Lebesgue set of $f$. In particular, the limit holds for a.e. $x$.
\end{theorem}

Since $\int K_\delta=1$,
\[
(f*K_\delta)(x)-f(x)
=\int [f(x-y)-f(x)]K_\delta(y)\,dy,
\]
and hence
\[
|(f*K_\delta)(x)-f(x)|
\le\int |f(x-y)-f(x)|\,|K_\delta(y)|\,dy.
\]
The proof uses the following lemma.

\begin{lemma}
Suppose $f$ is integrable on $\mathbb R^d$, and $x$ belongs to the Lebesgue set of $f$. Let
\[
A(r)=\frac1{r^d}\int_{|y|\le r}|f(x-y)-f(x)|\,dy,
\qquad r>0.
\]
Then $A(r)$ is continuous for $r>0$, $A(r)\to0$ as $r\to0$, and $A$ is bounded on $(0,\infty)$.
\end{lemma}

\begin{proof}
Continuity follows from the absolute continuity of the integral. Since $x$ is in the Lebesgue set and a ball of radius $r$ has measure $v_dr^d$, we have $A(r)\to0$ as $r\to0$. Continuity then gives boundedness for $0<r\le1$. For $r>1$,
\begin{align*}
A(r)
&\le \frac1{r^d}\int_{|y|\le r}|f(x-y)|\,dy
 +\frac1{r^d}\int_{|y|\le r}|f(x)|\,dy\\
&\le r^{-d}\|f\|_{L^1(\mathbb R^d)}+v_d|f(x)|,
\end{align*}
which proves the assertion.
\end{proof}

\begin{proof}[Proof of Theorem 2.1]
Decompose $\mathbb R^d$ into the ball $|y|\le\delta$ and dyadic annuli:
\[
\int |f(x-y)-f(x)|\,|K_\delta(y)|\,dy
=
\int_{|y|\le\delta}+\sum_{k=0}^\infty
\int_{2^k\delta<|y|\le2^{k+1}\delta}.
\]
Using (ii$'$), the first term is bounded by $cA(\delta)$. Using (iii$'$), the $k$th annular term is bounded by
\begin{align*}
&\frac{c\delta}{(2^k\delta)^{d+1}}
 \int_{|y|\le2^{k+1}\delta}|f(x-y)-f(x)|\,dy\\
&\hspace{3cm}\le c'2^{-k}A(2^{k+1}\delta).
\end{align*}
Therefore
\[
|(f*K_\delta)(x)-f(x)|
\le cA(\delta)+c'\sum_{k=0}^\infty2^{-k}A(2^{k+1}\delta).
\]
Given $\epsilon>0$, choose $N$ so large that $\sum_{k\ge N}2^{-k}<\epsilon$. By taking $\delta$ sufficiently small, $A(2^k\delta)<\epsilon/N$ for $0\le k\le N$. Since $A$ is bounded, the tail is $O(\epsilon)$, and the result follows.
\end{proof}

\begin{theorem}
Suppose $f$ is integrable on $\mathbb R^d$ and $\{K_\delta\}_{\delta>0}$ is an approximation to the identity. Then, for every $\delta>0$, the convolution
\[
(f*K_\delta)(x)=\int_{\mathbb R^d}f(x-y)K_\delta(y)\,dy
\]
is integrable, and
\[
\|(f*K_\delta)-f\|_{L^1(\mathbb R^d)}\longrightarrow0
\qquad\text{as }\delta\to0.
\]
\end{theorem}

The proof is the same argument used for the Gaussian kernel in the Fourier inversion discussion of Chapter 2, and is not repeated in the source.

\section{Differentiability of functions}

We now take up the second question raised at the beginning of this chapter, that of finding a broad condition on functions $F$ that guarantees the identity
\begin{equation}
F(b)-F(a)=\int_a^b F'(x)\,dx.
\tag{6}
\end{equation}
Two phenomena make a general formulation problematic. First, continuous functions need not be differentiable. Second, even if $F'(x)$ exists everywhere, the derivative need not be Lebesgue integrable. We therefore approach the problem through functions of bounded variation, a class closely related to rectifiable curves.

\subsection{Functions of bounded variation}

Let $\gamma$ be a parametrized curve in the plane,
\[
z(t)=(x(t),y(t)),\qquad a\le t\le b,
\]
where $x$ and $y$ are continuous. The curve is \emph{rectifiable} if there is $M<\infty$ such that for every partition
$a=t_0<t_1<\cdots<t_N=b$,
\begin{equation}
\sum_{j=1}^N|z(t_j)-z(t_{j-1})|\le M.
\tag{7}
\end{equation}
The length is
\[
L(\gamma)=\sup_{a=t_0<\cdots<t_N=b}
\sum_{j=1}^N|z(t_j)-z(t_{j-1})|.
\]

Suppose $F$ is a complex-valued function on $[a,b]$. For a partition $a=t_0<\cdots<t_N=b$, its variation on that partition is
\[
\sum_{j=1}^N|F(t_j)-F(t_{j-1})|.
\]
The function $F$ is said to be of \emph{bounded variation} if these sums are bounded independently of the partition. A refinement of a partition can only increase the corresponding variation.

\begin{figure}[ht]
\centering
\fbox{\parbox{0.72\textwidth}{\centering Source Figure 3 omitted: a rectifiable curve and polygonal approximations.}}
\caption{Approximation of a rectifiable curve by polygonal lines.}
\end{figure}

\begin{theorem}
A curve parametrized by $(x(t),y(t))$, $a\le t\le b$, is rectifiable if and only if both $x(t)$ and $y(t)$ are of bounded variation.
\end{theorem}

The proof is immediate from $F(t)=x(t)+iy(t)$ and the inequalities
\[
|a+ib|\le |a|+|b|\le2|a+ib|,
\qquad a,b\in\mathbb R.
\]

A real-valued function $F$ is \emph{increasing} if $F(t_1)\le F(t_2)$ whenever $t_1\le t_2$, and strictly increasing if the inequality is strict whenever $t_1<t_2$.

\paragraph{Example 1.}
A bounded monotone real-valued function is of bounded variation. If $F$ is increasing and bounded by $M$, then
\[
\sum_{j=1}^N|F(t_j)-F(t_{j-1})|
=F(b)-F(a)\le2M.
\]

\paragraph{Example 2.}
If $F$ is differentiable everywhere and $F'$ is bounded, then $F$ is of bounded variation. Indeed, if $|F'|\le M$, the mean value theorem gives
\[
|F(x)-F(y)|\le M|x-y|,
\]
and consequently the total variation is at most $M(b-a)$.

\paragraph{Example 3.}
Let
\[
F(x)=
\begin{cases}
x^a\sin(x^{-b}),&0<x\le1,\\
0,&x=0.
\end{cases}
\]
Then $F$ is of bounded variation on $[0,1]$ if and only if $a>b$.

\begin{figure}[ht]
\centering
\fbox{\parbox{0.72\textwidth}{\centering Source Figure 4 omitted: graphs of $x^a\sin(x^{-b})$ in the cases $a>b$, $a=b$, and $a<b$.}}
\caption{Graphs of $x^a\sin(x^{-b})$ for different values of $a$ and $b$.}
\end{figure}

The \emph{total variation} of $F$ on $[a,x]$ is
\[
T_F(a,x)=\sup_P\sum_{j=1}^N|F(t_j)-F(t_{j-1})|,
\]
where the supremum is over all partitions of $[a,x]$. For real-valued $F$, define the positive and negative variations by
\[
P_F(a,x)=\sup_P\sum_{(+)}\bigl(F(t_j)-F(t_{j-1})\bigr),
\]
where the sum is over those $j$ for which $F(t_j)\ge F(t_{j-1})$, and
\[
N_F(a,x)=\sup_P\sum_{(-)}-\bigl(F(t_j)-F(t_{j-1})\bigr),
\]
where the sum is over those $j$ for which $F(t_j)\le F(t_{j-1})$.

\begin{lemma}
Suppose $F$ is real-valued and of bounded variation on $[a,b]$. Then for every $a\le x\le b$,
\[
F(x)-F(a)=P_F(a,x)-N_F(a,x),
\]
and
\[
T_F(a,x)=P_F(a,x)+N_F(a,x).
\]
\end{lemma}

\begin{proof}
Given $\epsilon>0$, choose partitions which approximate $P_F$ and $N_F$ to within $\epsilon$, and pass to a common refinement. On that refinement,
\[
F(x)-F(a)
=\sum_{(+)}(F(t_j)-F(t_{j-1}))
-\sum_{(-)}-(F(t_j)-F(t_{j-1})).
\]
Letting $\epsilon\to0$ gives the first identity. Also, on every partition,
\[
\sum_j|F(t_j)-F(t_{j-1})|
=\sum_{(+)}(F(t_j)-F(t_{j-1}))
 +\sum_{(-)}-(F(t_j)-F(t_{j-1})),
\]
and the second identity follows by taking suprema and again using common refinements.
\end{proof}

\begin{theorem}
A real-valued function $F$ on $[a,b]$ is of bounded variation if and only if $F$ is the difference of two increasing bounded functions.
\end{theorem}

\begin{proof}
If $F=F_1-F_2$ with $F_1,F_2$ bounded and increasing, then $F$ is plainly of bounded variation. Conversely, if $F$ has bounded variation, set
\[
F_1(x)=P_F(a,x)+F(a),\qquad F_2(x)=N_F(a,x).
\]
Both are increasing and bounded, and the preceding lemma gives $F=F_1-F_2$.
\end{proof}

A complex-valued function of bounded variation is therefore a complex linear combination of four bounded increasing functions.

For a rectifiable curve $z(t)$, let $L(A,B)$ denote the length of the segment corresponding to $A\le t\le B$. Then
\[
L(A,C)+L(C,B)=L(A,B)\qquad(A\le C\le B).
\tag{8}
\]
Moreover, when $z$ is continuous, the length function $L(A,B)$ is continuous in both endpoints. Equivalently: if a function of bounded variation is continuous, then so is its total variation.

The next theorem is central to the theory of differentiation.

\begin{theorem}
If $F$ is of bounded variation on $[a,b]$, then $F$ is differentiable almost everywhere.
\end{theorem}

By the preceding decomposition theorem it is enough to treat increasing functions. We first assume in addition that $F$ is continuous; discontinuities are handled later by jump functions.

\begin{lemma}[Riesz rising sun lemma]
Suppose $G$ is real-valued and continuous on $\mathbb R$. Let $E$ be the set of points $x$ such that
\[
G(x+h)>G(x)
\]
for some $h=h_x>0$. If $E$ is non-empty, it is open, hence a countable disjoint union $E=\bigcup_k(a_k,b_k)$. If $(a_k,b_k)$ is finite, then
\[
G(b_k)-G(a_k)=0.
\]
\end{lemma}

\begin{proof}
The set $E$ is open by continuity. Let $(a_k,b_k)$ be a finite component. Since $a_k\notin E$, we cannot have $G(b_k)>G(a_k)$. Suppose $G(b_k)<G(a_k)$. By continuity choose $c\in(a_k,b_k)$, farthest to the right, such that
\[
G(c)=\frac{G(a_k)+G(b_k)}2.
\]
Because $c\in E$, there is $d>c$ with $G(d)>G(c)$. Since $b_k\notin E$, necessarily $d<b_k$. Continuity then gives $c'>d$ with $c'<b_k$ and $G(c')=G(c)$, contradicting the maximal choice of $c$. Hence $G(a_k)=G(b_k)$.
\end{proof}

\begin{figure}[ht]
\centering
\fbox{\parbox{0.72\textwidth}{\centering Source Figure 5 omitted: the ``rising sun'' picture, with shaded portions of the graph.}}
\caption{The rising sun lemma.}
\end{figure}

\begin{corollary}
Suppose $G$ is real-valued and continuous on $[a,b]$. Let $E$ be the set of $x\in(a,b)$ for which $G(x+h)>G(x)$ for some $h>0$. Then $E$ is empty or open; in the latter case it is a disjoint union of intervals $(a_k,b_k)$ and
\[
G(a_k)=G(b_k),
\]
except possibly when $a_k=a$, in which case only $G(a_k)\le G(b_k)$ is guaranteed.
\end{corollary}

For the proof of Theorem 3.4 define
\[
\Delta_h(F)(x)=\frac{F(x+h)-F(x)}{h}.
\]
The four Dini numbers at $x$ are
\begin{align*}
D^+(F)(x)&=\limsup_{\substack{h\to0\\h>0}}\Delta_h(F)(x),&
D_+(F)(x)&=\liminf_{\substack{h\to0\\h>0}}\Delta_h(F)(x),\\
D^-(F)(x)&=\limsup_{\substack{h\to0\\h<0}}\Delta_h(F)(x),&
D_-(F)(x)&=\liminf_{\substack{h\to0\\h<0}}\Delta_h(F)(x).
\end{align*}
Clearly $D_+\le D^+$ and $D_-\le D^-$. It is enough to show
\begin{enumerate}
\item[(i)] $D^+(F)(x)<\infty$ for a.e. $x$;
\item[(ii)] $D^+(F)(x)\le D_-(F)(x)$ for a.e. $x$.
\end{enumerate}
Applying (ii) to $-F(-x)$ gives $D^-(F)\le D_+(F)$ a.e.; consequently all four Dini numbers are equal and finite a.e.

Assume now that $F$ is increasing, bounded, and continuous. For $\gamma>0$ let
\[
E_\gamma=\{x:D^+(F)(x)>\gamma\}.
\]
This set is measurable. Apply the preceding corollary to $G(x)=F(x)-\gamma x$. We obtain a union of disjoint intervals $(a_k,b_k)$ containing $E_\gamma$ and satisfying
\[
F(b_k)-F(a_k)\ge\gamma(b_k-a_k).
\]
Therefore
\[
m(E_\gamma)
\le \sum_k(b_k-a_k)
\le \frac1\gamma\sum_k(F(b_k)-F(a_k))
\le \frac{F(b)-F(a)}\gamma.
\]
Letting $\gamma\to\infty$ proves (i).

For (ii), fix rationals $R>r$ and let
\[
E=\{x\in[a,b]:D^+(F)(x)>R\text{ and }r>D_-(F)(x)\}.
\]
Suppose $m(E)>0$. Choose an open set $O$ with
\[
E\subset O\subset(a,b),\qquad m(O)<m(E)\frac Rr.
\]
Write $O=\bigcup_n I_n$ as a disjoint union of open intervals. Applying the corollary first to $-F(-x)+rx$ on $-I_n$ and reflecting back produces disjoint intervals $(a_k,b_k)\subset I_n$ with
\[
F(b_k)-F(a_k)\le r(b_k-a_k).
\]
Applying it again, now to $F(x)-Rx$ on each $(a_k,b_k)$, gives an open set
\[
O_n=\bigcup_{k,j}(a_{k,j},b_{k,j})
\]
containing $E\cap I_n$ and satisfying
\[
F(b_{k,j})-F(a_{k,j})\ge R(b_{k,j}-a_{k,j}).
\]
Since $F$ is increasing,
\[
m(O_n)
\le \frac1R\sum_{k,j}(F(b_{k,j})-F(a_{k,j}))
\le \frac rR m(I_n).
\]
Summing over $n$ gives
\[
m(E)\le\sum_n m(O_n)\le\frac rR m(O)<m(E),
\]
a contradiction. This proves Theorem 3.4 in the continuous increasing case.

\begin{corollary}
If $F$ is increasing and continuous, then $F'$ exists almost everywhere. Moreover $F'$ is measurable, non-negative, and
\[
\int_a^bF'(x)\,dx\le F(b)-F(a).
\]
In particular, if $F$ is bounded on $\mathbb R$, then $F'$ is integrable on $\mathbb R$.
\end{corollary}

\begin{proof}
For $n\ge1$ let
\[
G_n(x)=\frac{F(x+1/n)-F(x)}{1/n}.
\]
Then $G_n(x)\to F'(x)$ for a.e. $x$, so $F'$ is measurable and non-negative. Fatou's lemma gives
\[
\int_a^bF'(x)\,dx\le\liminf_{n\to\infty}\int_a^bG_n(x)\,dx.
\]
Extend $F$ continuously to $\mathbb R$. A change of variables yields
\begin{align*}
\int_a^bG_n(x)\,dx
&=n\int_b^{b+1/n}F(x)\,dx
-n\int_a^{a+1/n}F(x)\,dx.
\end{align*}
By continuity these two terms tend to $F(b)$ and $F(a)$, respectively.
\end{proof}

The inequality cannot in general be replaced by equality for all continuous increasing functions, as shown by the Cantor--Lebesgue function.

\subsubsection*{The Cantor--Lebesgue function}

Let $C\subset[0,1]$ be the standard ternary Cantor set, $C=\bigcap_{k\ge0}C_k$. Define $F_1$ to be the continuous increasing function with $F_1(0)=0$, $F_1=1/2$ on $[1/3,2/3]$, $F_1(1)=1$, and linear on each component of $C_1$. Define $F_2$ similarly by
\[
F_2(x)=
\begin{cases}
0,&x=0,\\
1/4,&1/9\le x\le2/9,\\
1/2,&1/3\le x\le2/3,\\
3/4,&7/9\le x\le8/9,\\
1,&x=1,
\end{cases}
\]
and let it be linear on the intervals making up $C_2$. Continuing this construction gives continuous increasing functions $F_n$ satisfying
\[
|F_{n+1}(x)-F_n(x)|\le2^{-n-1}.
\]
Thus $F_n$ converges uniformly to a continuous increasing limit $F$, the \emph{Cantor--Lebesgue function}. It satisfies
\[
F(0)=0,\qquad F(1)=1,
\]
and is constant on every component of $[0,1]\setminus C$. Since $m(C)=0$,
\[
F'(x)=0\qquad\text{for a.e. }x,
\]
although $F(1)-F(0)=1$.

\begin{figure}[ht]
\centering
\fbox{\parbox{0.72\textwidth}{\centering Source Figures 6--7 omitted: construction and graph of the Cantor--Lebesgue function.}}
\caption{The Cantor--Lebesgue function.}
\end{figure}

\subsection{Absolutely continuous functions}

A function $F$ on $[a,b]$ is \emph{absolutely continuous} if for every $\epsilon>0$ there exists $\delta>0$ such that
\[
\sum_{k=1}^N|F(b_k)-F(a_k)|<\epsilon
\]
whenever the intervals $(a_k,b_k)$ are disjoint and
\[
\sum_{k=1}^N(b_k-a_k)<\delta.
\]
Several facts follow immediately from the definition.
\begin{itemize}
\item An absolutely continuous function is continuous, indeed uniformly continuous.
\item On a bounded interval, an absolutely continuous function is of bounded variation. Its total variation is itself absolutely continuous.
\item If
\[
F(x)=\int_a^x f(y)\,dy
\]
with $f$ integrable, then $F$ is absolutely continuous.
\end{itemize}
Thus absolute continuity is a necessary condition if one hopes for the full identity in (6).

\begin{theorem}
If $F$ is absolutely continuous on $[a,b]$, then $F'(x)$ exists almost everywhere. Moreover, if $F'(x)=0$ for a.e. $x$, then $F$ is constant.
\end{theorem}

The existence of $F'$ a.e. follows because an absolutely continuous function is the difference of two continuous increasing functions. To prove the second assertion we need a stronger covering result.

A collection $\mathcal B$ of balls is a \emph{Vitali covering} of a set $E$ if for every $x\in E$ and every $\eta>0$ there is a ball $B\in\mathcal B$ with $x\in B$ and $m(B)<\eta$.

\begin{lemma}[Vitali covering lemma]
Suppose $E$ has finite measure and $\mathcal B$ is a Vitali covering of $E$. For every $\delta>0$ there exist finitely many pairwise disjoint balls $B_1,\ldots,B_N\in\mathcal B$ such that
\[
\sum_{i=1}^Nm(B_i)\ge m(E)-\delta.
\]
\end{lemma}

\begin{proof}
Let $\gamma=3^{-d}$. We may assume $0<\delta<m(E)$. Choose a compact set $E_0\subset E$ with $m(E_0)\ge\delta$. By compactness, finitely many balls from the Vitali covering cover $E_0$. Lemma 1.2 supplies a disjoint subcollection $B_1,\ldots,B_{N_1}$ such that
\[
\sum_{i=1}^{N_1}m(B_i)\ge\gamma\delta.
\]
If this sum is already at least $m(E)-\delta$, we stop. Otherwise remove the selected balls from $E$ and repeat the argument on the remainder, choosing new balls disjoint from all previously selected ones. At the $k$th stage, unless the process has already stopped, the sum of the measures of selected balls is at least $k\gamma\delta$. Hence the process must stop after finitely many stages, certainly once
\[
k\ge\frac{m(E)-\delta}{\gamma\delta}.
\]
The resulting finite pairwise disjoint family has the desired total measure.
\end{proof}

\begin{corollary}
The balls in the preceding lemma can be chosen so that
\[
m\left(E-\bigcup_{i=1}^NB_i\right)<2\delta.
\]
\end{corollary}

\begin{proof}
Choose an open set $O\supset E$ with $m(O-E)<\delta$, and restrict the Vitali covering to balls contained in $O$. Then
\[
\left(E-\bigcup_{i=1}^NB_i\right)\cup\bigcup_{i=1}^NB_i\subset O
\]
is a disjoint union, and therefore
\begin{align*}
m\left(E-\bigcup_{i=1}^NB_i\right)
&\le m(O)-m\left(\bigcup_{i=1}^NB_i\right)\\
&\le m(E)+\delta-(m(E)-\delta)=2\delta.
\end{align*}
\end{proof}

\begin{proof}[Completion of the proof of Theorem 3.8]
It is enough to prove $F(b)=F(a)$; then the same argument applies to every subinterval. Let
\[
E=\{x\in(a,b):F'(x)\text{ exists and }F'(x)=0\}.
\]
By assumption $m(E)=b-a$.

Fix $\epsilon>0$. For each $x\in E$ and every $\eta>0$, choose an interval $I_x=(a_x,b_x)\subset[a,b]$ containing $x$, of length $<\eta$, and satisfying
\[
|F(b_x)-F(a_x)|\le\epsilon(b_x-a_x).
\]
These intervals form a Vitali covering of $E$. For $\delta>0$ select finitely many disjoint intervals $I_i=(a_i,b_i)$ such that
\begin{equation}
\sum_{i=1}^Nm(I_i)\ge(b-a)-\delta.
\tag{9}
\end{equation}
Then
\[
\sum_{i=1}^N|F(b_i)-F(a_i)|\le\epsilon(b-a).
\]
The complement of their union in $[a,b]$ is a finite union of closed intervals $[\alpha_k,\beta_k]$ with total length at most $\delta$. By absolute continuity, choosing $\delta$ small enough gives
\[
\sum_k|F(\beta_k)-F(\alpha_k)|\le\epsilon.
\]
Thus
\[
|F(b)-F(a)|
\le\epsilon(b-a)+\epsilon.
\]
Since $\epsilon$ is arbitrary, $F(b)=F(a)$.
\end{proof}

The preceding results culminate in the fundamental theorem of calculus for Lebesgue integration.

\begin{theorem}
Suppose $F$ is absolutely continuous on $[a,b]$. Then $F'$ exists almost everywhere and is integrable. Moreover,
\[
F(x)-F(a)=\int_a^xF'(y)\,dy,
\qquad a\le x\le b.
\]
In particular,
\[
F(b)-F(a)=\int_a^bF'(y)\,dy.
\]
Conversely, if $f$ is integrable on $[a,b]$, then there exists an absolutely continuous $F$ with $F'(x)=f(x)$ almost everywhere; indeed one may take
\[
F(x)=\int_a^x f(y)\,dy.
\]
\end{theorem}

\begin{proof}
A real-valued absolutely continuous function is the difference of two continuous increasing functions, so Corollary 3.7 shows that $F'$ is integrable. Let
\[
G(x)=\int_a^xF'(y)\,dy.
\]
Then $G$ is absolutely continuous, and hence so is $G-F$. By the Lebesgue differentiation theorem, $G'(x)=F'(x)$ for a.e. $x$, so $(F-G)'=0$ almost everywhere. Theorem 3.8 implies that $F-G$ is constant; evaluating at $a$ yields the claimed identity.

Conversely, if $F(x)=\int_a^x f(y)\,dy$, then $F$ is absolutely continuous and the Lebesgue differentiation theorem gives $F'=f$ almost everywhere.
\end{proof}

\subsection{Differentiability of jump functions}

We now remove the continuity assumption from the differentiability theorem for functions of bounded variation. It is enough to consider a bounded increasing function $F$.

For such an $F$, the one-sided limits
\[
F(x^-)=\lim_{\substack{y\to x\\y<x}}F(y),
\qquad
F(x^+)=\lim_{\substack{y\to x\\y>x}}F(y)
\]
exist and satisfy
\[
F(x^-)\le F(x)\le F(x^+).
\]
If $F(x^-)\ne F(x^+)$ we say that $F$ has a \emph{jump discontinuity} at $x$.

\begin{lemma}
A bounded increasing function $F$ on $[a,b]$ has at most countably many discontinuities.
\end{lemma}

\begin{proof}
If $F$ is discontinuous at $x$, choose a rational $r_x$ with
\[
F(x^-)<r_x<F(x^+).
\]
If $x<z$ are two discontinuities, then $F(x^+)\le F(z^-)$, and hence $r_x<r_z$. Therefore distinct discontinuities determine distinct rationals, and the set of discontinuities is countable.
\end{proof}

Let $\{x_n\}_{n\ge1}$ be the discontinuities and
\[
\alpha_n=F(x_n^+)-F(x_n^-)
\]
the jump at $x_n$. Write
\[
F(x_n)=F(x_n^-)+\theta_n\alpha_n,
\qquad0\le\theta_n\le1,
\]
and define
\[
j_n(x)=
\begin{cases}
0,&x<x_n,\\
\theta_n,&x=x_n,\\
1,&x>x_n.
\end{cases}
\]
The \emph{jump function} associated to $F$ is
\[
J_F(x)=\sum_{n=1}^\infty\alpha_nj_n(x).
\]
Since $F$ is bounded and increasing,
\[
\sum_{n=1}^\infty\alpha_n\le F(b)-F(a)<\infty,
\]
so the defining series converges absolutely and uniformly.

\begin{lemma}
If $F$ is increasing and bounded on $[a,b]$, then:
\begin{enumerate}
\item[(i)] $J$ is discontinuous precisely at the points $x_n$, and the jump of $J$ at $x_n$ equals the jump of $F$ there;
\item[(ii)] $F-J$ is increasing and continuous.
\end{enumerate}
\end{lemma}

\begin{proof}
If $x\ne x_n$ for every $n$, then each $j_n$ is continuous at $x$, and uniform convergence implies that $J$ is continuous there. At $x=x_N$, split the series into a finite sum through $N$ and a uniformly convergent tail; the tail is continuous at $x_N$, while the finite part has a jump of size $\alpha_N$.

Thus $F-J$ is continuous. If $y>x$, then
\[
J(y)-J(x)\le\sum_{x<x_n\le y}\alpha_n\le F(y)-F(x),
\]
so
\[
F(x)-J(x)\le F(y)-J(y),
\]
and $F-J$ is increasing.
\end{proof}

It remains to show that a jump function is differentiable almost everywhere.

\begin{theorem}
If $J$ is the jump function above, then $J'(x)$ exists and equals $0$ almost everywhere.
\end{theorem}

\begin{proof}
Fix $\epsilon>0$, and let $E$ be the measurable set of $x$ for which
\begin{equation}
\limsup_{h\to0}\frac{J(x+h)-J(x)}{h}>\epsilon.
\tag{10}
\end{equation}
Write $\delta=m(E)$. We show $\delta=0$.

Given $\eta>0$, choose $N$ so large that $\sum_{n>N}\alpha_n<\eta$, and put
\[
J_0(x)=\sum_{n>N}\alpha_nj_n(x).
\]
Then
\begin{equation}
J_0(b)-J_0(a)<\eta.
\tag{11}
\end{equation}
The finite sum $J-J_0$ is locally constant away from $x_1,\ldots,x_N$, so the set where (10) holds with $J_0$ in place of $J$ differs from $E$ by at most a finite set. Choose a compact set $K$ of measure at least $\delta/2$ on which
\[
\limsup_{h\to0}\frac{J_0(x+h)-J_0(x)}{h}>\epsilon.
\]
For every $x\in K$ choose an interval $(a_x,b_x)$ containing $x$ with
\[
J_0(b_x)-J_0(a_x)>\epsilon(b_x-a_x).
\]
A finite subcollection covers $K$. Lemma 1.2, in dimension one, supplies pairwise disjoint intervals $I_j=(a_j,b_j)$ such that
\[
\sum_jm(I_j)\ge\frac13m(K).
\]
Therefore
\[
J_0(b)-J_0(a)
\ge\sum_j(J_0(b_j)-J_0(a_j))
>\epsilon\sum_j(b_j-a_j)
\ge\frac\epsilon3m(K)
\ge\frac\epsilon6\delta.
\]
Together with (11), this gives $\epsilon\delta/6<\eta$. Since $\eta$ is arbitrary, $\delta=0$. Applying the argument also to the corresponding lower Dini derivative shows that $J'=0$ a.e.
\end{proof}

Combining Theorem 3.14 with the continuous case proves Theorem 3.4 in full generality.

\section{Rectifiable curves and the isoperimetric inequality}

We return to the formula
\begin{equation}
L=\int_a^b\bigl(x'(t)^2+y'(t)^2\bigr)^{1/2}\,dt
\tag{12}
\end{equation}
for the length $L$ of the curve parametrized by $(x(t),y(t))$.

Rectifiable curves are precisely those continuous curves whose coordinate functions are of bounded variation. Formula (12), however, need not hold under this hypothesis alone. For example, let
\[
x(t)=F(t),\qquad y(t)=F(t),\qquad 0\le t\le1,
\]
where $F$ is the Cantor--Lebesgue function. The curve traces the straight segment from $(0,0)$ to $(1,1)$ and therefore has length $\sqrt2$, while $x'(t)=y'(t)=0$ for almost every $t$.

The formula is restored when the coordinate functions are absolutely continuous.

\begin{theorem}
Suppose $(x(t),y(t))$ is a curve for $a\le t\le b$. If both $x$ and $y$ are absolutely continuous, then the curve is rectifiable and
\[
L=\int_a^b\bigl(x'(t)^2+y'(t)^2\bigr)^{1/2}\,dt.
\]
\end{theorem}

This follows from the next proposition, applied to $F(t)=x(t)+iy(t)$.

\begin{proposition}
Suppose $F$ is complex-valued and absolutely continuous on $[a,b]$. Then
\[
T_F(a,b)=\int_a^b|F'(t)|\,dt.
\]
\end{proposition}

\begin{proof}
For any partition $a=t_0<t_1<\cdots<t_N=b$, Theorem 3.11 gives
\begin{align*}
\sum_{j=1}^N|F(t_j)-F(t_{j-1})|
&=\sum_{j=1}^N\left|\int_{t_{j-1}}^{t_j}F'(t)\,dt\right|\\
&\le\sum_{j=1}^N\int_{t_{j-1}}^{t_j}|F'(t)|\,dt\\
&=\int_a^b|F'(t)|\,dt.
\end{align*}
Hence
\begin{equation}
T_F(a,b)\le\int_a^b|F'(t)|\,dt.
\tag{13}
\end{equation}

For the reverse inequality, fix $\epsilon>0$ and choose a step function $g$ on $[a,b]$ such that
\[
F'=g+h,
\qquad
\int_a^b|h(t)|\,dt<\epsilon.
\]
Set
\[
G(x)=\int_a^xg(t)\,dt,
\qquad
H(x)=\int_a^xh(t)\,dt.
\]
Then $F=G+H$ up to an additive constant, and
\[
T_F(a,b)\ge T_G(a,b)-T_H(a,b).
\]
By (13), $T_H(a,b)<\epsilon$. Choose a partition on whose intervals the step function $g$ is constant. Then
\begin{align*}
T_G(a,b)
&\ge\sum_j|G(t_j)-G(t_{j-1})|\\
&=\sum_j\left|\int_{t_{j-1}}^{t_j}g(t)\,dt\right|\\
&=\int_a^b|g(t)|\,dt.
\end{align*}
Since
\[
\int_a^b|g(t)|\,dt
\ge\int_a^b|F'(t)|\,dt-\epsilon,
\]
we obtain
\[
T_F(a,b)\ge\int_a^b|F'(t)|\,dt-2\epsilon.
\]
Letting $\epsilon\to0$ proves the proposition.
\end{proof}

A rectifiable curve has a natural \emph{arc-length parametrization}. If $L(A,B)$ is its length function, put
\[
s=s(t)=L(a,t),
\]
so $s$ maps $[a,b]$ continuously and increasingly onto $[0,L]$. Define
\[
\widetilde z(s)=z(t)\quad\text{when }s=s(t).
\]
This is well defined: if $s(t_1)=s(t_2)$, then the curve has zero length on $[t_1,t_2]$, so $z(t_1)=z(t_2)$. Moreover,
\[
|\widetilde z(s_1)-\widetilde z(s_2)|\le|s_1-s_2|.
\]

\begin{theorem}
Suppose $(x(t),y(t))$, $a\le t\le b$, is rectifiable and has length $L$. Let
\[
\widetilde z(s)=(\widetilde x(s),\widetilde y(s)),\qquad0\le s\le L,
\]
be the arc-length parametrization. Then $\widetilde x$ and $\widetilde y$ are absolutely continuous,
\[
|\widetilde z'(s)|=1
\qquad\text{for a.e. }s\in[0,L],
\]
and
\[
L=\int_0^L\bigl(\widetilde x'(s)^2+\widetilde y'(s)^2\bigr)^{1/2}\,ds.
\]
\end{theorem}

\begin{proof}
The Lipschitz estimate
\[
|\widetilde z(s_1)-\widetilde z(s_2)|\le|s_1-s_2|
\]
implies absolute continuity and $|\widetilde z'(s)|\le1$ a.e. The total variation of $\widetilde z$ is $L$, so Proposition 4.2 gives
\[
L=\int_0^L|\widetilde z'(s)|\,ds.
\]
Since the integrand is bounded above by $1$, equality is possible only if $|\widetilde z'(s)|=1$ a.e.
\end{proof}

\subsection*{4.1* Minkowski content of a curve}

A curve $z(t)$, $a\le t\le b$, is \emph{simple} if $t\mapsto z(t)$ is injective on $[a,b]$. It is a \emph{closed simple curve} if it is injective on $[a,b)$ and $z(a)=z(b)$. More generally, it is \emph{quasi-simple} if the parametrization is injective outside finitely many points.

\begin{figure}[ht]
\centering
\fbox{\parbox{0.72\textwidth}{\centering Source Figure 8 omitted: a quasi-simple curve.}}
\caption{A quasi-simple curve.}
\end{figure}

Let $\Gamma$ denote the point set traced by the curve. For a compact set $K\subset\mathbb R^2$, write
\[
K^\delta=\{x\in\mathbb R^2:d(x,K)<\delta\}.
\]

\begin{figure}[ht]
\centering
\fbox{\parbox{0.72\textwidth}{\centering Source Figure 9 omitted: a curve $\Gamma$ and its $\delta$-neighborhood $\Gamma^\delta$.}}
\caption{The curve $\Gamma$ and the set $\Gamma^\delta$.}
\end{figure}

We say that $K$ has \emph{Minkowski content} if the limit
\[
\lim_{\delta\to0}\frac{m(K^\delta)}{2\delta}
\]
exists, and denote it by $\mathcal M(K)$.

\begin{theorem}
Suppose $\Gamma=\{z(t):a\le t\le b\}$ is a quasi-simple curve. The Minkowski content of $\Gamma$ exists if and only if $\Gamma$ is rectifiable. In that case,
\[
\mathcal M(\Gamma)=L,
\]
where $L$ is the length of the curve.
\end{theorem}

Define the upper and lower Minkowski contents by
\[
\mathcal M^*(K)=\limsup_{\delta\to0}\frac{m(K^\delta)}{2\delta},
\qquad
\mathcal M_*(K)=\liminf_{\delta\to0}\frac{m(K^\delta)}{2\delta}.
\]
The theorem follows from two inequalities.

\begin{proposition}
Suppose $\Gamma=\{z(t):a\le t\le b\}$ is quasi-simple. If $\mathcal M_*(\Gamma)<\infty$, then the curve is rectifiable, and its length $L$ satisfies
\[
L\le\mathcal M_*(\Gamma).
\]
\end{proposition}

The proof uses the following observation.

\begin{lemma}
If $\Gamma=\{z(t):a\le t\le b\}$ is any curve, and
\[
\Delta=|z(b)-z(a)|,
\]
then
\[
m(\Gamma^\delta)\ge2\delta\Delta.
\]
\end{lemma}

\begin{proof}
Translate and rotate so that the endpoints are $(A,0)$ and $(B,0)$, with $A<B$ and $\Delta=B-A$. For each $x\in[A,B]$, continuity supplies $t$ with $x=x(t)$. The vertical slice $(\Gamma^\delta)_x$ contains the interval
\[
(y(t)-\delta,y(t)+\delta),
\]
so its one-dimensional measure is at least $2\delta$. Fubini's theorem gives
\[
m(\Gamma^\delta)
=\int_{\mathbb R}m_1((\Gamma^\delta)_x)\,dx
\ge\int_A^B2\delta\,dx
=2\delta(B-A).
\]
\end{proof}

\begin{figure}[ht]
\centering
\fbox{\parbox{0.72\textwidth}{\centering Source Figure 10 omitted: vertical slices through $\Gamma^\delta$.}}
\caption{The situation in Lemma 4.6.}
\end{figure}

\begin{proof}[Proof of Proposition 4.5]
First suppose the curve is simple. Let $P$ be a partition $a=t_0<\cdots<t_N=b$ and
\[
L_P=\sum_{j=1}^N|z(t_j)-z(t_{j-1})|.
\]
For every $\epsilon>0$, continuity allows us to choose proper closed subintervals $I_j=[a_j,b_j]\subset(t_{j-1},t_j)$ such that
\[
\sum_{j=1}^N|z(b_j)-z(a_j)|\ge L_P-\epsilon.
\]
Let $\Gamma_j$ be the corresponding pieces of the curve. Simplicity implies that the compact sets $\Gamma_j$ are disjoint, hence for sufficiently small $\delta$ their $\delta$-neighborhoods are disjoint. By Lemma 4.6,
\[
m(\Gamma^\delta)
\ge\sum_{j=1}^Nm((\Gamma_j)^\delta)
\ge2\delta\sum_{j=1}^N|z(b_j)-z(a_j)|.
\]
Therefore
\[
\frac{m(\Gamma^\delta)}{2\delta}\ge L_P-\epsilon.
\]
Passing to the lower limit and then taking the supremum over $P$ gives
\[
\mathcal M_*(\Gamma)\ge L.
\]
For a quasi-simple curve, refine the partitions to include the finitely many exceptional parameters; the same argument applies piecewise.
\end{proof}

\begin{proposition}
Suppose $\Gamma$ is rectifiable with length $L$. Then
\[
\mathcal M^*(\Gamma)\le L.
\]
\end{proposition}

\begin{proof}
Use the arc-length parametrization $z(s)$, $0\le s\le L$, so $z$ is absolutely continuous and $|z'(s)|=1$ for a.e. $s$.

Fix $0<\epsilon<1$. By Egorov's theorem there exist a measurable set $E_\epsilon$ with $m(E_\epsilon)<\epsilon$ and $r_\epsilon>0$ such that
\begin{equation}
\sup_{0<|h|<r_\epsilon}
\left|\frac{z(s+h)-z(s)}h-z'(s)\right|<\epsilon
\qquad(s\in[0,L]\setminus E_\epsilon).
\tag{14}
\end{equation}
We may also suppose $z'(s)$ exists and $|z'(s)|=1$ outside $E_\epsilon$.

Choose $0<\rho<r_\epsilon$, $\rho<1$, and partition $[0,L]$ into consecutive intervals $I_1,\ldots,I_N$ of length $\rho$ except possibly the last, so $N\le L/\rho+1$. Call $I_j$ \emph{good} if it contains a point outside $E_\epsilon$, and \emph{bad} otherwise. The union of the bad intervals has measure $<\epsilon$, hence there are at most $\epsilon/\rho+1$ of them.

Let $\Gamma_j=z(I_j)$. For a good interval $I_j=[a_j,b_j]$, choose $s_0\in I_j\setminus E_\epsilon$ and translate and rotate coordinates so that $z(s_0)=0$ and $z'(s_0)=1$. By (14), $\Gamma_j$ is contained in
\[
[a_j-s_0-\epsilon\rho,\ b_j-s_0+\epsilon\rho]
\times[-\epsilon\rho,\epsilon\rho].
\]
Consequently $(\Gamma_j)^\delta$ is contained in the enlargement
\[
[a_j-s_0-\epsilon\rho-\delta,\ b_j-s_0+\epsilon\rho+\delta]
\times[-\epsilon\rho-\delta,\epsilon\rho+\delta],
\]
which gives
\begin{equation}
m((\Gamma_j)^\delta)
\le2\delta\rho+O(\epsilon\delta\rho+\delta^2+\epsilon\rho^2).
\tag{15}
\end{equation}
For an arbitrary interval, the arc-length parametrization satisfies $|z(s)-z(s')|\le|s-s'|$, so $\Gamma_j$ lies in a disc of radius $\rho$. Hence
\begin{equation}
m((\Gamma_j)^\delta)=O(\delta^2+\rho^2).
\tag{16}
\end{equation}
Summing (15) over good intervals and (16) over bad ones gives
\[
\frac{m(\Gamma^\delta)}{2\delta}
\le L+O\left(
\rho+\epsilon+\frac\delta\rho+\frac{\epsilon\rho}\delta
+\frac{\rho^2}\delta
\right).
\]
Choose
\[
\rho=\frac\delta{\epsilon^{1/2}},
\]
and restrict to $0<\delta<\epsilon^{1/2}r_\epsilon$. Then
\[
\limsup_{\delta\to0}\frac{m(\Gamma^\delta)}{2\delta}
\le L+O(\epsilon+\epsilon^{1/2}).
\]
Letting $\epsilon\to0$ yields $\mathcal M^*(\Gamma)\le L$.
\end{proof}

Together, Propositions 4.5 and 4.7 prove Theorem 4.4.

\subsection*{4.2* Isoperimetric inequality}

The isoperimetric inequality in the plane says that, among curves of a fixed length, the circle encloses the largest area. Let $\Omega$ be a bounded open subset of $\mathbb R^2$ whose boundary $\Gamma$ is a rectifiable curve of length $\ell(\Gamma)$. The curve need not be simple.

\begin{theorem}[Isoperimetric inequality]
\[
4\pi\,m(\Omega)\le\ell(\Gamma)^2.
\]
\end{theorem}

\begin{proof}
For $\delta>0$ define
\[
\Omega^+(\delta)=\{x\in\mathbb R^2:d(x,\Omega)<\delta\},
\]
and
\[
\Omega^-(\delta)=\{x\in\mathbb R^2:d(x,\Omega^c)\ge\delta\}.
\]
Then
\[
\Omega^-(\delta)\subset\Omega\subset\Omega^+(\delta).
\]
If
\[
\Gamma^\delta=\{x:d(x,\Gamma)<\delta\},
\]
then the union
\begin{equation}
\Omega^+(\delta)=\Omega^-(\delta)\cup\Gamma^\delta
\tag{17}
\end{equation}
is disjoint.

Let $D(\delta)$ be the open disc of radius $\delta$ centered at the origin. Then
\begin{equation}
\begin{cases}
\Omega^+(\delta)\supset\Omega+D(\delta),\\
\Omega\supset\Omega^-(\delta)+D(\delta).
\end{cases}
\tag{18}
\end{equation}

\begin{figure}[ht]
\centering
\fbox{\parbox{0.72\textwidth}{\centering Source Figure 12 omitted: $\Omega^-(\delta)\subset\Omega\subset\Omega^+(\delta)$.}}
\caption{The sets $\Omega$, $\Omega^-(\delta)$ and $\Omega^+(\delta)$.}
\end{figure}

Apply the Brunn--Minkowski inequality to the first inclusion in (18):
\[
m(\Omega^+(\delta))^{1/2}
\ge m(\Omega)^{1/2}+m(D(\delta))^{1/2}.
\]
Since $m(D(\delta))=\pi\delta^2$,
\[
m(\Omega^+(\delta))
\ge m(\Omega)+2\pi^{1/2}\delta\,m(\Omega)^{1/2}.
\]
Similarly, the second inclusion yields
\[
m(\Omega)
\ge m(\Omega^-(\delta))
 +2\pi^{1/2}\delta\,m(\Omega^-(\delta))^{1/2},
\]
or
\[
-m(\Omega^-(\delta))
\ge -m(\Omega)+2\pi^{1/2}\delta\,m(\Omega^-(\delta))^{1/2}.
\]
Using (17),
\[
m(\Gamma^\delta)
=m(\Omega^+(\delta))-m(\Omega^-(\delta))
\]
and therefore
\[
m(\Gamma^\delta)
\ge2\pi^{1/2}\delta
\left(m(\Omega)^{1/2}+m(\Omega^-(\delta))^{1/2}\right).
\]
Divide by $2\delta$ and take the limsup as $\delta\to0$. Since $\Omega^-(\delta)\uparrow\Omega$,
\[
\mathcal M^*(\Gamma)
\ge2\pi^{1/2}m(\Omega)^{1/2}.
\]
By Proposition 4.7, $\ell(\Gamma)\ge\mathcal M^*(\Gamma)$, hence
\[
\ell(\Gamma)\ge2\pi^{1/2}m(\Omega)^{1/2},
\]
which is equivalent to the desired inequality.
\end{proof}

\paragraph{Remark.}
The same argument does not actually require the boundary to be a rectifiable curve. For any bounded open set $\Omega$ whose boundary is $\Gamma$,
\[
4\pi\,m(\Omega)\le\mathcal M^*(\Gamma)^2.
\]

% End of Chapter 3. Exercises and Problems omitted by request.

% Chapters 4 and 5 are intentionally omitted.
\setcounter{chapter}{5}

% Stein & Shakarchi, Real Analysis (Princeton Lectures in Analysis III)
% Chapter 6: Abstract Measure and Integration Theory
% Transcribed from the user-provided PDF.
% Exercises and Problems are intentionally omitted.
% This file is intended to be \input{} into a book/report document.
% Suggested packages in the parent document:
%   amsmath, amssymb, amsthm, mathtools, graphicx
% Suggested theorem declarations:
%   \newtheorem{theorem}{Theorem}[section]
%   \newtheorem{lemma}[theorem]{Lemma}
%   \newtheorem{corollary}[theorem]{Corollary}
%   \newtheorem{proposition}[theorem]{Proposition}
% Figures from the source are represented by compile-safe placeholders.

\chapter{Abstract Measure and Integration Theory}

\begin{quote}
What immediately suggest itself, then, is that these characteristic properties themselves be treated as the main object of investigation, by defining and dealing with abstract objects which need satisfy no other conditions than those required by the very theory to be developed.

This procedure has been made use of -- more or less consciously -- by mathematicians of every era. The geometry of Euclid and the literal algebra of the sixteenth and seventeenth centuries arose in this way. But only in more recent times has this method, called the axiomatic method, been consistently developed and carried through to its logical conclusion.

It is our intention to treat the theories of measure and integration by means of the axiomatic method just described.

\hfill C. Carath\'eodory, 1918
\end{quote}

In much of mathematics integration plays a significant role. It is used, in one form or another, when dealing with questions that arise in analysis on a variety of different spaces. While in some situations it suffices to integrate continuous or other simple functions on these spaces, the deeper study of a number of other problems requires integration based on the more refined ideas of measure theory. The development of these ideas, going beyond the setting of the Euclidean space $\mathbb R^d$, is the goal of this chapter.

The starting point is a fruitful insight of Carath\'eodory and the resulting theorems that lead to construction of measures in very general circumstances. Once this has been achieved, the deduction of the fundamental facts about integration in the general context then follows a familiar path.

We apply the abstract theory to obtain several useful results: the theory of product measures; the polar coordinate integration formula, which is a consequence of this; the construction of the Lebesgue--Stieltjes integral and its corresponding Borel measure on the real line; and the general notion of absolute continuity. Finally, we treat some of the basic limit theorems of ergodic theory. This not only gives an illustration of the abstract framework we have established, but also provides a link with the differentiation theorems studied in Chapter 3.

\section{Abstract measure spaces}

A \emph{measure space} consists of a set $X$ equipped with two fundamental objects:
\begin{enumerate}
\item[(I)] A $\sigma$-algebra $\mathcal M$ of ``measurable'' sets, which is a non-empty collection of subsets of $X$ closed under complements and countable unions and intersections.
\item[(II)] A measure $\mu:\mathcal M\to[0,\infty]$ with the following defining property: if $E_1,E_2,\ldots$ is a countable family of disjoint sets in $\mathcal M$, then
\[
\mu\left(\bigcup_{n=1}^{\infty}E_n\right)=\sum_{n=1}^{\infty}\mu(E_n).
\]
\end{enumerate}
A measure space is therefore often denoted by the triple $(X,\mathcal M,\mu)$ to emphasize its three main components. Sometimes, however, when there is no ambiguity we will abbreviate this notation by referring to the measure space as $(X,\mu)$, or simply $X$.

A feature that a measure space often enjoys is the property of being $\sigma$-finite. This means that $X$ can be written as the union of countably many measurable sets of finite measure.

At this early stage we give only two simple examples of measure spaces:
\begin{enumerate}
\item[(i)] The first is the discrete example with $X$ a countable set, $X=\{x_n\}_{n=1}^{\infty}$, $\mathcal M$ the collection of all subsets of $X$, and the measure $\mu$ determined by $\mu(x_n)=\mu_n$, with $\{\mu_n\}_{n=1}^{\infty}$ a given sequence of (extended) non-negative numbers. Note that
\[
\mu(E)=\sum_{x_n\in E}\mu_n.
\]
When $\mu_n=1$ for all $n$, we call $\mu$ the \emph{counting measure}, and also denote it by $\#$. In this case integration will amount to nothing but the summation of (absolutely) convergent series.
\item[(ii)] Here $X=\mathbb R^d$, $\mathcal M$ is the collection of Lebesgue measurable sets, and
\[
\mu(E)=\int_E f\,dx,
\]
where $f$ is a given non-negative measurable function on $\mathbb R^d$. The case $f=1$ corresponds to the Lebesgue measure. The countable additivity of $\mu$ follows from the usual additivity and limiting properties of integrals of non-negative functions proved in Chapter 2.
\end{enumerate}

The construction of measure spaces relevant for most applications requires further ideas, and to these we now turn.

\subsection{Exterior measures and Carath\'eodory's theorem}

To begin the construction of a measure and its corresponding measurable sets in the general setting requires, as in the special case of Lebesgue measure considered in Chapter 1, a prerequisite notion of ``exterior'' measure. This is defined as follows.

Let $X$ be a set. An \emph{exterior measure} (or \emph{outer measure}) $\mu^*$ on $X$ is a function $\mu^*$ from the collection of all subsets of $X$ to $[0,\infty]$ that satisfies the following properties:
\begin{enumerate}
\item[(i)] $\mu^*(\varnothing)=0$.
\item[(ii)] If $E_1\subset E_2$, then $\mu^*(E_1)\leq \mu^*(E_2)$.
\item[(iii)] If $E_1,E_2,\ldots$ is a countable family of sets, then
\[
\mu^*\left(\bigcup_{j=1}^{\infty}E_j\right)\leq\sum_{j=1}^{\infty}\mu^*(E_j).
\]
\end{enumerate}
For instance, the exterior Lebesgue measure $m^*$ in $\mathbb R^d$ defined in Chapter 1 enjoys all these properties. In fact, this example belongs to a large class of exterior measures that can be obtained using ``coverings'' by a family of special sets whose measures are taken as known. This idea is systematized by the notion of a ``premeasure'' taken up below in Section 1.3. A different type of example is the exterior $\alpha$-dimensional Hausdorff measure $m^*_{\alpha}$ defined in Chapter 7.

Given an exterior measure $\mu^*$, the problem that one faces is how to define the corresponding notion of measurable sets. In the case of Lebesgue measure in $\mathbb R^d$ such sets were characterized by their difference from open (or closed) sets, when considered in terms of $\mu^*$. For the general case, Carath\'eodory found an ingenious substitute condition. It is as follows.

A set $E$ in $X$ is \emph{Carath\'eodory measurable} or simply \emph{measurable} if one has
\begin{equation}
\mu^*(A)=\mu^*(E\cap A)+\mu^*(E^c\cap A)\qquad\text{for every }A\subset X.
\tag{1}
\end{equation}
In other words, $E$ separates any set $A$ in two parts that behave well in regard to the exterior measure $\mu^*$. For this reason, (1) is sometimes referred to as the \emph{separation condition}. One can show that in $\mathbb R^d$ with the Lebesgue exterior measure the notion of measurability (1) is equivalent to the definition of Lebesgue measurability given in Chapter 1. (See Exercise 3.)

A first observation we make is that to prove a set $E$ is measurable, it suffices to verify
\[
\mu^*(A)\geq \mu^*(E\cap A)+\mu^*(E^c\cap A)\qquad\text{for all }A\subset X,
\]
since the reverse inequality is automatically verified by the sub-additivity property (iii) of the exterior measure. We see immediately from the definition that sets of exterior measure zero are necessarily measurable.

\begin{theorem}
Given an exterior measure $\mu^*$ on a set $X$, the collection $\mathcal M$ of Carath\'eodory measurable sets forms a $\sigma$-algebra. Moreover, $\mu^*$ restricted to $\mathcal M$ is a measure.
\end{theorem}

\begin{proof}
Clearly, $\varnothing$ and $X$ belong to $\mathcal M$ and the symmetry inherent in condition (1) shows that $E^c\in\mathcal M$ whenever $E\in\mathcal M$. Thus $\mathcal M$ is non-empty and closed under complements.

Next, we prove that $\mathcal M$ is closed under finite unions of disjoint sets, and $\mu^*$ is finitely additive on $\mathcal M$. Indeed, if $E_1,E_2\in\mathcal M$, and $A$ is any subset of $X$, then
\begin{align*}
\mu^*(A)
&=\mu^*(E_2\cap A)+\mu^*(E_2^c\cap A)\\
&=\mu^*(E_1\cap E_2\cap A)+\mu^*(E_1^c\cap E_2\cap A)\\
&\quad+\mu^*(E_1\cap E_2^c\cap A)+\mu^*(E_1^c\cap E_2^c\cap A)\\
&\geq \mu^*((E_1\cup E_2)\cap A)+\mu^*((E_1\cup E_2)^c\cap A).
\end{align*}
Therefore $E_1\cup E_2\in\mathcal M$. If $E_1$ and $E_2$ are disjoint, then
\begin{align*}
\mu^*(E_1\cup E_2)
&=\mu^*(E_1\cap(E_1\cup E_2))+\mu^*(E_1^c\cap(E_1\cup E_2))\\
&=\mu^*(E_1)+\mu^*(E_2).
\end{align*}

Finally, let $E_1,E_2,\ldots$ be a countable collection of disjoint sets in $\mathcal M$, and define
\[
G_n=\bigcup_{j=1}^nE_j,
\qquad
G=\bigcup_{j=1}^{\infty}E_j.
\]
For each $n$, the set $G_n$ is a finite union of sets in $\mathcal M$, hence $G_n\in\mathcal M$. Moreover, for any $A\subset X$,
\[
\mu^*(G_n\cap A)=\sum_{j=1}^n\mu^*(E_j\cap A),
\]
by induction. Since $G^c\subset G_n^c$,
\[
\mu^*(A)=\mu^*(G_n\cap A)+\mu^*(G_n^c\cap A)
\geq \sum_{j=1}^n\mu^*(E_j\cap A)+\mu^*(G^c\cap A).
\]
Letting $n\to\infty$ gives
\[
\mu^*(A)\geq\sum_{j=1}^{\infty}\mu^*(E_j\cap A)+\mu^*(G^c\cap A)
\geq \mu^*(G\cap A)+\mu^*(G^c\cap A)
\geq \mu^*(A).
\]
Thus all inequalities are equalities, $G\in\mathcal M$, and, by taking $A=G$, $\mu^*$ is countably additive on $\mathcal M$.
\end{proof}

Our previous observation that sets of exterior measure $0$ are Carath\'eodory measurable shows that the measure space $(X,\mathcal M,\mu)$ in the theorem is \emph{complete}: whenever $F\in\mathcal M$ satisfies $\mu(F)=0$ and $E\subset F$, then $E\in\mathcal M$.

\subsection{Metric exterior measures}

If the underlying set $X$ is endowed with a ``distance function'' or ``metric,'' there is a particular class of exterior measures that is of interest in practice. The importance of these exterior measures is that they induce measures on the natural $\sigma$-algebra generated by the open sets in $X$.

A \emph{metric space} is a set $X$ equipped with a function $d:X\times X\to[0,\infty)$ that satisfies:
\begin{enumerate}
\item[(i)] $d(x,y)=0$ if and only if $x=y$.
\item[(ii)] $d(x,y)=d(y,x)$ for all $x,y\in X$.
\item[(iii)] $d(x,z)\leq d(x,y)+d(y,z)$ for all $x,y,z\in X$.
\end{enumerate}
The last property is the triangle inequality. For example, $\mathbb R^d$ with $d(x,y)=|x-y|$ is a metric space. Another example is the space of continuous functions on a compact set $K$ with
\[
d(f,g)=\sup_{x\in K}|f(x)-g(x)|.
\]

A metric space $(X,d)$ is naturally equipped with a family of open balls
\[
B_r(x)=\{y\in X:d(x,y)<r\}.
\]
A set $O\subset X$ is open if for any $x\in O$ there exists $r>0$ so that $B_r(x)\subset O$. A set is closed if its complement is open. Finally, on a metric space $X$ we define the Borel $\sigma$-algebra $\mathcal B_X$, the smallest $\sigma$-algebra of sets in $X$ that contains the open sets of $X$. Its elements are called \emph{Borel sets}.

Given two sets $A$ and $B$ in a metric space $(X,d)$, define
\[
d(A,B)=\inf\{d(x,y):x\in A,\ y\in B\}.
\]
An exterior measure $\mu^*$ on $X$ is a \emph{metric exterior measure} if
\[
\mu^*(A\cup B)=\mu^*(A)+\mu^*(B)\qquad\text{whenever }d(A,B)>0.
\]

\begin{theorem}
If $\mu^*$ is a metric exterior measure on a metric space $X$, then the Borel sets in $X$ are measurable. Hence $\mu^*$ restricted to $\mathcal B_X$ is a measure.
\end{theorem}

\begin{proof}
It suffices to prove that closed sets in $X$ are Carath\'eodory measurable. Let $F$ be closed and $A\subset X$ with $\mu^*(A)<\infty$. For $n>0$ let
\[
A_n=\{x\in F^c\cap A:d(x,F)\geq1/n\}.
\]
Then $A_n\subset A_{n+1}$ and $F^c\cap A=\bigcup_{n=1}^{\infty}A_n$. Since the distance between $F\cap A$ and $A_n$ is at least $1/n$,
\begin{equation}
\mu^*(A)\geq \mu^*((F\cap A)\cup A_n)=\mu^*(F\cap A)+\mu^*(A_n).
\tag{2}
\end{equation}
We claim that
\begin{equation}
\lim_{n\to\infty}\mu^*(A_n)=\mu^*(F^c\cap A).
\tag{3}
\end{equation}
Let $B_n=A_{n+1}\cap A_n^c$. One checks from the triangle inequality that
\[
d(B_{n+1},A_n)\geq \frac1{n(n+1)}.
\]
Hence metric additivity gives
\[
\mu^*(A_{2k+1})\geq \sum_{j=1}^{k}\mu^*(B_{2j}),
\qquad
\mu^*(A_{2k})\geq \sum_{j=1}^{k}\mu^*(B_{2j-1}).
\]
Since $\mu^*(A)<\infty$, both series $\sum\mu^*(B_{2j})$ and $\sum\mu^*(B_{2j-1})$ converge. Finally,
\[
\mu^*(A_n)\leq\mu^*(F^c\cap A)
\leq\mu^*(A_n)+\sum_{j=n+1}^{\infty}\mu^*(B_j),
\]
which proves (3). Letting $n\to\infty$ in (2) proves the separation condition for $F$.
\end{proof}

Given a metric space $X$, a measure $\mu$ defined on the Borel sets of $X$ will be referred to as a \emph{Borel measure}. Borel measures that assign a finite measure to all balls of finite radius also satisfy a useful regularity property.

\begin{proposition}
Suppose the Borel measure $\mu$ is finite on all balls in $X$ of finite radius. Then for any Borel set $E$ and any $\epsilon>0$, there are an open set $O$ and a closed set $F$ such that $E\subset O$ and $\mu(O-E)<\epsilon$, while $F\subset E$ and $\mu(E-F)<\epsilon$.
\end{proposition}

\begin{proof}
We need the following preliminary observation. Suppose $F^*=\bigcup_{k=1}^{\infty}F_k$, where the $F_k$ are closed. Then for any $\epsilon>0$ we can find a closed set $F\subset F^*$ such that $\mu(F^*-F)<\epsilon$. We may assume the $F_k$ are increasing. Fix $x_0\in X$, and let $B_n=\{x:d(x,x_0)<n\}$, with $B_0=\varnothing$. Since $\bigcup_{n=1}^{\infty}B_n=X$,
\[
F^*=\bigcup_{n=1}^{\infty}\bigl(F^*\cap(B_n-B_{n-1})\bigr).
\]
For each $n$, $F^*\cap(B_n-B_{n-1})$ is the increasing limit of $F_k\cap(B_n-B_{n-1})$, and $B_n$ has finite measure. Hence choose $N(n)$ so that
\[
\mu\bigl((F^*-F_{N(n)})\cap(B_n-B_{n-1})\bigr)<\epsilon/2^n.
\]
Set
\[
F=\bigcup_{n=1}^{\infty}\bigl(F_{N(n)}\cap(B_n-B_{n-1})\bigr).
\]
Then $\mu(F^*-F)<\epsilon$. Moreover $F\cap B_k$ is closed for every $k$, hence $F$ is closed.

Let $\mathcal C$ be the collection of all sets satisfying the two regularity conclusions of the proposition. If $E\in\mathcal C$, then $E^c\in\mathcal C$. If $E=\bigcup_{k=1}^{\infty}E_k$ with $E_k\in\mathcal C$, choose open $O_k\supset E_k$ so that $\mu(O_k-E_k)<\epsilon/2^k$; then $O=\bigcup O_k$ works from outside. Similarly choose closed $F_k\subset E_k$ with $\mu(E_k-F_k)<\epsilon/2^k$. The union $F^*=\bigcup F_k$ approximates $E$ from inside, and the preliminary observation produces a closed $F\subset F^*$ with the required error. Thus $\mathcal C$ is closed under countable unions and complements.

Finally every open set $O$ belongs to $\mathcal C$. The outer condition is immediate. For the inner condition take
\[
F_k=\{x\in B_k:d(x,O^c)\geq1/k\}.
\]
Then each $F_k$ is closed and $O=\bigcup F_k$, so the preliminary observation applies. Thus $\mathcal C$ is a $\sigma$-algebra containing the open sets, and hence all Borel sets.
\end{proof}

\subsection{The extension theorem}

As we have seen, a class of measurable sets on $X$ can be constructed once we start with a given exterior measure. However, the definition of an exterior measure usually depends on a more primitive idea of measure defined on a simpler class of sets. This is the role of a premeasure.

An \emph{algebra} in $X$ is a non-empty collection of subsets of $X$ that is closed under complements, finite unions, and finite intersections. Let $\mathcal A$ be an algebra in $X$. A \emph{premeasure} on $\mathcal A$ is a function $\mu_0:\mathcal A\to[0,\infty]$ satisfying
\begin{enumerate}
\item[(i)] $\mu_0(\varnothing)=0$.
\item[(ii)] If $E_1,E_2,\ldots$ is a countable collection of disjoint sets in $\mathcal A$ with $\bigcup_{k=1}^{\infty}E_k\in\mathcal A$, then
\[
\mu_0\left(\bigcup_{k=1}^{\infty}E_k\right)=\sum_{k=1}^{\infty}\mu_0(E_k).
\]
\end{enumerate}
In particular, $\mu_0$ is finitely additive on $\mathcal A$.

\begin{lemma}
If $\mu_0$ is a premeasure on an algebra $\mathcal A$, define $\mu^*$ on any subset $E$ of $X$ by
\[
\mu^*(E)=\inf\left\{\sum_{j=1}^{\infty}\mu_0(E_j):
E\subset\bigcup_{j=1}^{\infty}E_j,\ E_j\in\mathcal A\text{ for all }j\right\}.
\]
Then $\mu^*$ is an exterior measure on $X$ that satisfies:
\begin{enumerate}
\item[(i)] $\mu^*(E)=\mu_0(E)$ for all $E\in\mathcal A$.
\item[(ii)] All sets in $\mathcal A$ are measurable in the sense of (1).
\end{enumerate}
\end{lemma}

\begin{proof}
Proving that $\mu^*$ is an exterior measure presents no difficulty. Suppose $E\in\mathcal A$. Clearly $\mu^*(E)\leq\mu_0(E)$. Conversely, if $E\subset\bigcup_{j=1}^{\infty}E_j$ with $E_j\in\mathcal A$, set
\[
E_k'=E\cap\left(E_k-\bigcup_{j=1}^{k-1}E_j\right).
\]
Then $E_k'$ are disjoint elements of $\mathcal A$, $E_k'\subset E_k$, and $E=\bigcup_{k=1}^{\infty}E_k'$. Hence
\[
\mu_0(E)=\sum_{k=1}^{\infty}\mu_0(E_k')\leq\sum_{k=1}^{\infty}\mu_0(E_k),
\]
so $\mu_0(E)\leq\mu^*(E)$.

Finally, let $A\subset X$, $E\in\mathcal A$, and $\epsilon>0$. Choose $E_j\in\mathcal A$ with $A\subset\bigcup E_j$ and
\[
\sum_{j=1}^{\infty}\mu_0(E_j)\leq\mu^*(A)+\epsilon.
\]
Finite additivity gives
\begin{align*}
\sum_{j=1}^{\infty}\mu_0(E_j)
&=\sum_{j=1}^{\infty}\mu_0(E\cap E_j)+\sum_{j=1}^{\infty}\mu_0(E^c\cap E_j)\\
&\geq \mu^*(E\cap A)+\mu^*(E^c\cap A).
\end{align*}
Letting $\epsilon\to0$ proves the separation condition.
\end{proof}

The $\sigma$-algebra generated by an algebra $\mathcal A$ is by definition the smallest $\sigma$-algebra that contains $\mathcal A$.

\begin{theorem}
Suppose that $\mathcal A$ is an algebra of sets in $X$, $\mu_0$ a premeasure on $\mathcal A$, and $\mathcal M$ the $\sigma$-algebra generated by $\mathcal A$. Then there exists a measure $\mu$ on $\mathcal M$ that extends $\mu_0$.
\end{theorem}

One notes below that $\mu$ is the only such extension of $\mu_0$ under the assumption that $\mu$ is $\sigma$-finite.

\begin{proof}
The exterior measure $\mu^*$ induced by $\mu_0$ defines a measure $\mu$ on the $\sigma$-algebra of Carath\'eodory measurable sets. By the previous lemma, $\mu$ is therefore a measure on $\mathcal M$ extending $\mu_0$.

To prove uniqueness when $\mu$ is $\sigma$-finite, suppose $\nu$ is another measure on $\mathcal M$ coinciding with $\mu_0$ on $\mathcal A$. Let $F\in\mathcal M$ have finite measure. If $F\subset\bigcup_jE_j$ with $E_j\in\mathcal A$, then
\[
\nu(F)\leq\sum_{j=1}^{\infty}\nu(E_j)=\sum_{j=1}^{\infty}\mu_0(E_j),
\]
so $\nu(F)\leq\mu(F)$. For the reverse inequality, if $E=\bigcup E_j$, then
\[
\nu(E)=\lim_{n\to\infty}\nu\left(\bigcup_{j=1}^nE_j\right)
=\lim_{n\to\infty}\mu\left(\bigcup_{j=1}^nE_j\right)=\mu(E).
\]
Choose $E_j$ so that $\mu(E)\leq\mu(F)+\epsilon$. Since $\mu(F)<\infty$, $\mu(E-F)\leq\epsilon$, and
\[
\mu(F)\leq\mu(E)=\nu(E)=\nu(F)+\nu(E-F)
\leq\nu(F)+\mu(E-F)\leq\nu(F)+\epsilon.
\]
Thus $\mu(F)=\nu(F)$.

Finally, since $\mu$ is $\sigma$-finite, write $X=\bigcup E_j$ as a disjoint union with $E_j\in\mathcal A$ and $\mu(E_j)<\infty$. Then for every $F\in\mathcal M$,
\[
\mu(F)=\sum_j\mu(F\cap E_j)=\sum_j\nu(F\cap E_j)=\nu(F).
\]
\end{proof}

For later use define $\mathcal A_{\sigma}$ as the collection of countable unions of sets in $\mathcal A$, and $\mathcal A_{\sigma\delta}$ as the collection of countable intersections of sets in $\mathcal A_{\sigma}$.

\begin{proposition}
For any set $E$ and any $\epsilon>0$, there are sets $E_1\in\mathcal A_{\sigma}$ and $E_2\in\mathcal A_{\sigma\delta}$ such that $E\subset E_1$, $E\subset E_2$, and $\mu^*(E_1)\leq\mu^*(E)+\epsilon$, while $\mu^*(E_2)=\mu^*(E)$.
\end{proposition}

\section{Integration on a measure space}

Once we have established the basic properties of a measure space $X$, the fundamental facts about measurable functions and integration of such functions on $X$ can be deduced as in the case of the Lebesgue measure on $\mathbb R^d$. Indeed, the results in Section 4 of Chapter 1 and all of Chapter 2 go over to the general case, with proofs remaining almost word-for-word the same. For this reason we shall not repeat these arguments but limit ourselves to the bare statement of the main points. To avoid unnecessary complications we assume throughout that the measure space $(X,\mathcal M,\mu)$ is $\sigma$-finite.

\subsection*{Measurable functions}

A function $f$ on $X$ with values in the extended real numbers is measurable if
\[
f^{-1}([ -\infty,a))=\{x\in X:f(x)<a\}\in\mathcal M
\qquad\text{for all }a\in\mathbb R.
\]
With this definition, the basic properties of measurable functions obtained in the case of $\mathbb R^d$ with the Lebesgue measure continue to hold. For instance, the collection of measurable functions is closed under the basic algebraic manipulations, and pointwise limits of measurable functions are measurable.

The notion of ``almost everywhere'' is now with respect to the measure $\mu$. Thus if $f$ and $g$ are measurable functions on $X$, we write $f=g$ a.e. to say that
\[
\mu(\{x\in X:f(x)\neq g(x)\})=0.
\]

A simple function on $X$ takes the form
\[
\sum_{k=1}^N a_k\chi_{E_k},
\]
where $E_k$ are measurable sets of finite measure and $a_k$ are real numbers. Approximations by simple functions continue to hold in this abstract setting:
\begin{itemize}
\item Suppose $f$ is a non-negative measurable function on a measure space $(X,\mathcal M,\mu)$. Then there exists a sequence of simple functions $\{\varphi_k\}_{k=1}^{\infty}$ satisfying
\[
\varphi_k(x)\leq\varphi_{k+1}(x),
\qquad
\lim_{k\to\infty}\varphi_k(x)=f(x)
\quad\text{for all }x.
\]
In general, if $f$ is merely measurable, there exists a sequence of simple functions $\{\varphi_k\}$ satisfying
\[
|\varphi_k(x)|\leq|\varphi_{k+1}(x)|,
\qquad
\lim_{k\to\infty}\varphi_k(x)=f(x)
\quad\text{for all }x.
\]
\end{itemize}
The proof is obtained by the same argument as Theorems 4.1 and 4.2 in Chapter 1, with the $\sigma$-finiteness of $X$ replacing the role played there by expanding cubes.

Another important result that generalizes immediately is Egorov's theorem:
\begin{itemize}
\item Suppose $\{f_k\}_{k=1}^{\infty}$ is a sequence of measurable functions defined on a measurable set $E\subset X$ with $\mu(E)<\infty$, and $f_k\to f$ a.e. Then for every $\epsilon>0$ there is a set $A_\epsilon\subset E$ with $\mu(E-A_\epsilon)\leq\epsilon$ and such that $f_k\to f$ uniformly on $A_\epsilon$.
\end{itemize}

\subsection*{Definition and main properties of the integral}

The four-step approach to the construction of the Lebesgue integral in Chapter 2 carries over to the situation of a $\sigma$-finite measure space $(X,\mathcal M,\mu)$. For a non-negative measurable function $f$ on $X$, its integral with respect to $\mu$ is denoted by
\[
\int_X f(x)\,d\mu(x),
\]
or, when no confusion can arise, by $\int_X f\,d\mu$, $\int f\,d\mu$, or simply $\int f$. A measurable function $f$ is integrable if
\[
\int_X|f(x)|\,d\mu(x)<\infty.
\]

The elementary properties of the integral, such as linearity and monotonicity, continue to hold, as do the following basic limit theorems:
\begin{enumerate}
\item[(i)] \textbf{Fatou's lemma.} If $\{f_n\}$ is a sequence of non-negative measurable functions on $X$, then
\[
\int_X\liminf_{n\to\infty}f_n\,d\mu
\leq
\liminf_{n\to\infty}\int_X f_n\,d\mu.
\]
\item[(ii)] \textbf{Monotone convergence.} If $\{f_n\}$ is a sequence of non-negative measurable functions with $f_n\nearrow f$, then
\[
\lim_{n\to\infty}\int_X f_n\,d\mu=\int_X f\,d\mu.
\]
\item[(iii)] \textbf{Dominated convergence.} If $\{f_n\}$ is a sequence of measurable functions with $f_n\to f$ a.e., and $|f_n|\leq g$ for some integrable $g$, then
\[
\int_X|f_n-f|\,d\mu\to0,
\]
and consequently
\[
\int_X f_n\,d\mu\to\int_X f\,d\mu.
\]
\end{enumerate}

\subsection*{The spaces $L^1(X,\mu)$ and $L^2(X,\mu)$}

The equivalence classes, modulo functions that vanish almost everywhere, of integrable functions on $(X,\mathcal M,\mu)$ form a vector space equipped with the norm
\begin{equation}
\|f\|_{L^1(X,\mu)}=\int_X|f(x)|\,d\mu(x).
\tag{4}
\end{equation}
This space is denoted by $L^1(X,\mu)$.

Similarly, $L^2(X,\mu)$ consists of equivalence classes of measurable functions for which
\[
\int_X|f(x)|^2\,d\mu(x)<\infty,
\]
with norm
\begin{equation}
\|f\|_{L^2(X,\mu)}=
\left(\int_X|f(x)|^2\,d\mu(x)\right)^{1/2}.
\tag{5}
\end{equation}
There is also an inner product
\[
(f,g)=\int_X f(x)\overline{g(x)}\,d\mu(x).
\]
The results proved earlier extend to this setting and give:
\begin{itemize}
\item The space $L^1(X,\mu)$ is a complete normed vector space.
\item The space $L^2(X,\mu)$ is a (possibly non-separable) Hilbert space.
\end{itemize}

\section{Examples}

We now discuss some useful examples of the general theory.

\subsection{Product measures and a general Fubini theorem}

Our first example concerns the construction of product measures, and leads to a general form of the theorem that expresses a multiple integral as a repeated integral, extending the case of Euclidean space considered in Section 3 of Chapter 2.

Suppose $(X_1,\mathcal M_1,\mu_1)$ and $(X_2,\mathcal M_2,\mu_2)$ are a pair of measure spaces. We want to describe the product measure $\mu_1\times\mu_2$ on
\[
X=X_1\times X_2=\{(x_1,x_2):x_1\in X_1,\ x_2\in X_2\}.
\]
We assume here that the two measure spaces are each complete and $\sigma$-finite.

We begin by considering measurable rectangles: subsets of $X$ of the form $A\times B$, with $A\in\mathcal M_1$ and $B\in\mathcal M_2$. Let $\mathcal A$ denote the collection of all sets in $X$ that are finite unions of disjoint measurable rectangles. It is easy to check that $\mathcal A$ is an algebra of subsets of $X$.

On rectangles define
\[
\mu_0(A\times B)=\mu_1(A)\mu_2(B).
\]
The fact that $\mu_0$ has a unique extension to the algebra $\mathcal A$ for which it becomes a premeasure follows from the following fact: whenever a rectangle $A\times B$ is the disjoint union of a countable collection of rectangles $\{A_j\times B_j\}$,
\begin{equation}
\mu_0(A\times B)=\sum_{j=1}^{\infty}\mu_0(A_j\times B_j).
\tag{6}
\end{equation}
Indeed, if $x_1\in A$, then $B$ is the disjoint union of the $B_j$ for which $x_1\in A_j$. Hence
\[
\chi_A(x_1)\mu_2(B)=\sum_{j=1}^{\infty}\chi_{A_j}(x_1)\mu_2(B_j).
\]
Integrating in $x_1$ and using monotone convergence gives (6).

The extension theorem now gives a measure $\mu=\mu_1\times\mu_2$ on the $\sigma$-algebra $\mathcal M$ generated by the algebra $\mathcal A$. In this way we obtain the product measure space
\[
(X_1\times X_2,\mathcal M,\mu_1\times\mu_2).
\]

Given $E\in\mathcal M$, consider the slices
\[
E_{x_1}=\{x_2\in X_2:(x_1,x_2)\in E\},
\qquad
E^{x_2}=\{x_1\in X_1:(x_1,x_2)\in E\}.
\]

\begin{proposition}
If $E$ belongs to $\mathcal A_{\sigma\delta}$, then $E^{x_2}$ is $\mu_1$-measurable for every $x_2$; moreover $\mu_1(E^{x_2})$ is a $\mu_2$-measurable function. In addition,
\begin{equation}
\int_{X_2}\mu_1(E^{x_2})\,d\mu_2=(\mu_1\times\mu_2)(E).
\tag{7}
\end{equation}
\end{proposition}

\begin{proof}
All assertions are immediate when $E$ is a measurable rectangle. Suppose next that $E\in\mathcal A_\sigma$. Decompose $E$ as a countable union of disjoint rectangles $E_j$. Then for each $x_2$,
\[
E^{x_2}=\bigcup_{j=1}^{\infty}E_j^{x_2},
\]
and the slices on the right are disjoint. Applying (7) to each rectangle and using monotone convergence proves the result for $E\in\mathcal A_\sigma$.

Assume next $E\in\mathcal A_{\sigma\delta}$ and $(\mu_1\times\mu_2)(E)<\infty$. Write $E=\bigcap_{j=1}^{\infty}E_j$, where $E_j\in\mathcal A_\sigma$ and $E_{j+1}\subset E_j$. Put
\[
f_j(x_2)=\mu_1(E_j^{x_2}),
\qquad
f(x_2)=\mu_1(E^{x_2}).
\]
Then $E^{x_2}$ is the decreasing limit of measurable sets $E_j^{x_2}$, and $f_j(x_2)\to f(x_2)$. Thus $f$ is measurable, and since the $f_j$ are decreasing non-negative functions,
\[
\int_{X_2}f(x_2)\,d\mu_2
=\lim_{j\to\infty}\int_{X_2}f_j(x_2)\,d\mu_2.
\]
This proves (7) in the finite-measure case.

For the general case use $\sigma$-finiteness to choose increasing sequences $F_j\subset X_1$, $G_j\subset X_2$ with finite measures and unions equal to $X_1$ and $X_2$, respectively. Apply the finite-measure result to $E\cap(F_j\times G_j)$ and let $j\to\infty$.
\end{proof}

\begin{proposition}
If $E$ is an arbitrary measurable set in $X$, then the conclusions of the previous proposition are still valid except that we only assert that $E^{x_2}$ is $\mu_1$-measurable and $\mu_1(E^{x_2})$ is defined for almost every $x_2\in X_2$.
\end{proposition}

\begin{proof}
Suppose first that $E$ has measure zero. Proposition 1.6 gives $F\in\mathcal A_{\sigma\delta}$ with $E\subset F$ and $(\mu_1\times\mu_2)(F)=0$. By (7), $F^{x_2}$ has $\mu_1$-measure zero for almost every $x_2$. Since $E^{x_2}\subset F^{x_2}$ and the measure spaces are complete, $E^{x_2}$ is measurable and has measure zero for those $x_2$.

For a general measurable $E$, choose $F\in\mathcal A_{\sigma\delta}$ with $F\supset E$ and $F-E=Z$ of measure zero. Then $F^{x_2}-E^{x_2}=Z^{x_2}$, and the result follows from the zero-measure case.
\end{proof}

\begin{theorem}[General Fubini theorem]
In the setting above, suppose $f(x_1,x_2)$ is an integrable function on $(X_1\times X_2,\mu_1\times\mu_2)$. Then:
\begin{enumerate}
\item[(i)] For almost every $x_2\in X_2$, the slice
\[
f^{x_2}(x_1)=f(x_1,x_2)
\]
is integrable on $(X_1,\mu_1)$.
\item[(ii)] The function
\[
x_2\longmapsto\int_{X_1}f(x_1,x_2)\,d\mu_1
\]
is integrable on $X_2$.
\item[(iii)]
\[
\int_{X_2}\left(\int_{X_1}f(x_1,x_2)\,d\mu_1\right)d\mu_2
=\int_{X_1\times X_2}f\,d(\mu_1\times\mu_2).
\]
\end{enumerate}
\end{theorem}

\begin{proof}
The conclusions are stable under finite linear combinations, so it suffices to consider non-negative functions. When $f=\chi_E$ with $E$ of finite measure, the result is contained in the previous proposition. Hence it holds for simple functions, and by monotone convergence for all non-negative functions. Applying the result to positive and negative parts gives the integrable case.
\end{proof}

In general the product space constructed above is not complete. If it is completed in the natural way, the theorem remains valid, with only a minor modification of the argument.

\subsection{Integration formula for polar coordinates}

The polar coordinates of a point $x\in\mathbb R^d-\{0\}$ are the pair $(r,\gamma)$, where $0<r<\infty$ and $\gamma$ belongs to the unit sphere
\[
S^{d-1}=\{x\in\mathbb R^d:|x|=1\}.
\]
They are determined by
\begin{equation}
r=|x|,
\qquad
\gamma=\frac{x}{|x|},
\qquad
x=r\gamma.
\tag{8}
\end{equation}
Our goal is the formula
\begin{equation}
\int_{\mathbb R^d}f(x)\,dx
=\int_{S^{d-1}}\left(\int_0^{\infty}f(r\gamma)r^{d-1}\,dr\right)d\sigma(\gamma).
\tag{9}
\end{equation}

Consider the measure space $(X_1,\mathcal M_1,\mu_1)$ with $X_1=(0,\infty)$, $\mathcal M_1$ the Lebesgue measurable sets, and
\[
d\mu_1(r)=r^{d-1}\,dr.
\]
Next let $X_2=S^{d-1}$. For $E\subset S^{d-1}$ put
\[
\widetilde E=\left\{x\in\mathbb R^d:\frac{x}{|x|}\in E,\ 0<|x|<1\right\}.
\]
Declare $E$ measurable exactly when $\widetilde E$ is Lebesgue measurable in $\mathbb R^d$, and define
\[
\mu_2(E)=\sigma(E)=d\,m(\widetilde E).
\]
Both spaces are complete and $\sigma$-finite. The sphere carries the metric $d(\gamma,\gamma')=|\gamma-\gamma'|$, and open subsets of the sphere are measurable in this construction.

\begin{theorem}
Suppose $f$ is integrable on $\mathbb R^d$. Then for almost every $\gamma\in S^{d-1}$ the slice
\[
f^\gamma(r)=f(r\gamma)
\]
is integrable with respect to $r^{d-1}\,dr$. Moreover,
\[
\int_0^\infty f^\gamma(r)r^{d-1}\,dr
\]
is integrable on $S^{d-1}$, and identity (9) holds. There is a corresponding result with the order of integration of $r$ and $\gamma$ reversed.
\end{theorem}

\begin{proof}
Let $\mu=\mu_1\times\mu_2$ on $X_1\times X_2$, identified with $\mathbb R^d-\{0\}$. We first show
\begin{equation}
m(E)=\mu(E)
\tag{10}
\end{equation}
for measurable rectangles $E=E_1\times E_2$. If $E_1=(0,1)$, then $E$ is the sector $\widetilde E_2$, while $\mu_1(E_1)=1/d$, so (10) follows from the definition of $\sigma$. Relative dilation invariance then gives (10) for $E_1=(0,b)$, and by limiting arguments and subtraction for arbitrary measurable $E_1$. Thus (10) holds on the algebra generated by measurable rectangles, and by uniqueness of the extension it holds on the generated $\sigma$-algebra.

Any open set in $\mathbb R^d-\{0\}$ can be written as a countable union of rectangles $A_j\times B_j$ with $A_j$ open in $(0,\infty)$ and $B_j$ open in $S^{d-1}$, so all Borel sets belong to the product $\sigma$-algebra. Lebesgue measurable sets then follow by completion. Hence (9) holds first for characteristic functions, then for simple functions, then by monotone convergence for non-negative functions, and finally for general integrable functions.
\end{proof}

\subsection{Borel measures on $\mathbb R$ and the Lebesgue--Stieltjes integral}

The Stieltjes integral was introduced to generalize the Riemann integral $\int_a^b f(x)\,dx$ by replacing the increments $dx$ by increments $dF(x)$ for a given increasing function $F$ on $[a,b]$. We pursue this idea from the general point of view of the present chapter.

To obtain a unique correspondence between measures and increasing functions, we first normalize the latter. Recall that an increasing function $F$ can have at most countably many discontinuities. If $x_0$ is one, then the one-sided limits
\[
F(x_0^-)=\lim_{\substack{x\to x_0\\x<x_0}}F(x),
\qquad
F(x_0^+)=\lim_{\substack{x\to x_0\\x>x_0}}F(x)
\]
exist, with $F(x_0^-)\leq F(x_0)\leq F(x_0^+)$. Modify $F$ at each discontinuity, if necessary, by setting $F(x_0)=F(x_0^+)$. The resulting increasing function is right-continuous at every point; such functions are called \emph{normalized}.

\begin{theorem}
Let $F$ be an increasing normalized function on $\mathbb R$. Then there is a unique measure $\mu$, also denoted by $dF$, on the Borel sets $\mathcal B$ of $\mathbb R$ such that
\[
\mu((a,b])=F(b)-F(a),\qquad a<b.
\]
Conversely, if $\mu$ is a measure on $\mathcal B$ that is finite on bounded intervals, then the function
\[
F(x)=
\begin{cases}
\mu((0,x]),&x>0,\\
0,&x=0,\\
-\mu((x,0]),&x<0,
\end{cases}
\]
is increasing and normalized.
\end{theorem}

\begin{proof}
Define an exterior measure candidate on all subsets $E\subset\mathbb R$ by
\[
\mu^*(E)=\inf\sum_{j=1}^{\infty}\bigl(F(b_j)-F(a_j)\bigr),
\]
where the infimum is taken over all coverings $E\subset\bigcup_{j=1}^{\infty}(a_j,b_j]$.
It is easy to verify that $\mu^*$ is an exterior measure.

We first show
\[
\mu^*((a,b])=F(b)-F(a).
\]
The inequality $\leq$ is immediate. Conversely, suppose $\bigcup_j(a_j,b_j]$ covers $(a,b]$. Fix $a<a'<b$ and $\epsilon>0$. By right-continuity choose $b_j'>b_j$ so that
\[
F(b_j')\leq F(b_j)+\epsilon/2^j.
\]
The open intervals $(a_j,b_j')$ cover $[a',b]$, so by compactness a finite subcollection covers it. Hence
\[
F(b)-F(a')
\leq\sum_{j=1}^N\bigl(F(b_j')-F(a_j)\bigr)
\leq\sum_{j=1}^N\bigl(F(b_j)-F(a_j)+\epsilon/2^j\bigr)
\leq\mu^*((a,b])+\epsilon.
\]
Letting $a'\to a$ and then $\epsilon\to0$ proves the equality.

Next we show that $\mu^*$ is a metric exterior measure for the usual metric on $\mathbb R$. Let $d(E_1,E_2)\geq\delta>0$. Given a covering of $E_1\cup E_2$ by half-open intervals with total $F$-increment within $\epsilon$ of $\mu^*(E_1\cup E_2)$, subdivide so every interval has length $<\delta$. Each interval then meets at most one of the two sets. Splitting the covering accordingly gives
\[
\mu^*(E_1)+\mu^*(E_2)\leq\mu^*(E_1\cup E_2)+\epsilon.
\]
Let $\epsilon\to0$. Hence $\mu^*$ is metric, so the Borel sets are measurable and the induced measure satisfies $\mu((a,b])=F(b)-F(a)$.

For uniqueness, suppose $\nu$ is another Borel measure with the same values on all $(a,b]$. Any open interval is a disjoint countable union of such half-open intervals, so $\mu$ and $\nu$ agree on open intervals, hence on open sets. Since both are finite on bounded intervals, the regularity proposition shows that they agree on all Borel sets.

Conversely, start from a Borel measure finite on bounded intervals and define $F$ as in the statement. It is clearly increasing. If $x_0>0$, the sets $(0,x_0+1/n]$ decrease to $(0,x_0]$, and finite measure gives
\[
F(x_0+1/n)\to F(x_0).
\]
Monotonicity then implies right-continuity. The other cases are similar.
\end{proof}

\paragraph{Remarks.}
Several comments about the theorem are useful.
\begin{enumerate}
\item[(i)] Two increasing functions $F$ and $G$ give the same measure if $F-G$ is constant. The converse is also true.
\item[(ii)] The measure constructed in the proof is defined on a larger $\sigma$-algebra than the Borel sets and is complete, although its restriction to the Borel sets often suffices.
\item[(iii)] If $F$ is an increasing normalized function on $[a,b]$, extend it to $\mathbb R$ by $F(x)=F(a)$ for $x<a$ and $F(x)=F(b)$ for $x>b$. Then one often writes
\[
\int_{\mathbb R}f(x)\,d\mu(x)=\int_a^b f(x)\,dF(x).
\]
If $F$ arises from an increasing function $F_0$ on $\mathbb R$ and one wishes to account for a possible jump at $a$, one may write
\[
\int_{a^-}^b f(x)\,dF(x)=\int_a^b f(x)\,d\mu_0(x),
\]
where $\mu_0$ corresponds to $F_0$.
\item[(iv)] The definition extends to functions of bounded variation. If
\[
F=\sum_{j=1}^4\epsilon_jF_j,
\qquad \epsilon_j\in\{\pm1,\pm i\},
\]
with the $F_j$ increasing and normalized, define
\[
\int_a^b f\,dF=\sum_{j=1}^4\epsilon_j\int_a^b f\,dF_j,
\]
requiring $f$ to be integrable with respect to the measure $\sum_j\mu_j$.
\item[(v)] The values of these integrals can be calculated more directly in the following cases:
\begin{enumerate}
\item[(a)] If $F$ is absolutely continuous on $[a,b]$, then
\[
\int_a^b f(x)\,dF(x)=\int_a^b f(x)F'(x)\,dx.
\]
\item[(b)] If $F$ is a pure jump function with jumps $\{\alpha_n\}$ at points $\{x_n\}$, then, for suitable $f$,
\[
\int_a^b f(x)\,dF(x)=\sum_{n=1}^{\infty}f(x_n)\alpha_n.
\]
In particular $\mu(\{x_n\})=\alpha_n$, and $\mu(E)=0$ for sets containing none of the $x_n$.
\item[(c)] For the Heaviside function $H(x)=1$ for $x\geq0$ and $H(x)=0$ for $x<0$,
\[
\int_{-\infty}^{\infty}f(x)\,dH(x)=f(0),
\]
which is another expression for the Dirac delta function.
\end{enumerate}
\end{enumerate}

\section{Absolute continuity of measures}

The generalization of the notion of absolute continuity considered in Chapter 3 requires that we extend the ideas of a measure to encompass set functions that may be positive or negative. We describe this notion first.

\subsection{Signed measures}

Loosely speaking, a signed measure possesses all the properties of a measure, except that it may take positive or negative values. More precisely, a \emph{signed measure} $\nu$ on a $\sigma$-algebra $\mathcal M$ is a mapping that satisfies:
\begin{enumerate}
\item[(i)] $\nu$ is extended-valued in the sense that $-\infty<\nu(E)\leq\infty$ for all $E\in\mathcal M$.
\item[(ii)] If $\{E_j\}_{j=1}^{\infty}$ are disjoint subsets of $\mathcal M$, then
\[
\nu\left(\bigcup_{j=1}^{\infty}E_j\right)=\sum_{j=1}^{\infty}\nu(E_j).
\]
\end{enumerate}
For this to hold, the sum must be independent of rearrangements of terms, so if the measure of the union is finite, the series converges absolutely.

Examples arise by dropping the non-negativity assumption in
\[
\nu(E)=\int_E f\,d\mu,
\]
where $(X,\mathcal M,\mu)$ is a measure space and $f$ is measurable. To guarantee the signed-measure properties, $f$ is required to be integrable in the extended sense that $\int f^-\,d\mu<\infty$, while $\int f^+\,d\mu$ may be infinite.

Given a signed measure $\nu$, it is always possible to find a positive measure that dominates it and is the smallest such measure. Define the \emph{total variation} of $\nu$ by
\[
|\nu|(E)=\sup\sum_{j=1}^{\infty}|\nu(E_j)|,
\]
where the supremum is taken over all measurable countable partitions $E=\bigcup_{j=1}^{\infty}E_j$.

\begin{proposition}
The total variation $|\nu|$ of a signed measure $\nu$ is itself a positive measure that satisfies $\nu\leq|\nu|$.
\end{proposition}

\begin{proof}
Suppose $\{E_j\}$ is a countable collection of disjoint measurable sets and $E=\bigcup_jE_j$. It suffices to prove
\begin{equation}
\sum_j|\nu|(E_j)\leq|\nu|(E),
\qquad
|\nu|(E)\leq\sum_j|\nu|(E_j).
\tag{11}
\end{equation}
For the first inequality choose real numbers $\alpha_j<|\nu|(E_j)$. For each $j$ write $E_j=\bigcup_iF_{i,j}$ disjointly so that
\[
\alpha_j\leq\sum_i|\nu(F_{i,j})|.
\]
Since the $F_{i,j}$ together partition $E$,
\[
\sum_j\alpha_j\leq\sum_{j,i}|\nu(F_{i,j})|\leq|\nu|(E).
\]
Taking suprema gives the first inequality.

For the reverse inequality, let $\{F_k\}$ be any partition of $E$. Since $\{F_k\cap E_j\}_j$ partitions $F_k$,
\begin{align*}
\sum_k|\nu(F_k)|
&=\sum_k\left|\sum_j\nu(F_k\cap E_j)\right|\\
&\leq\sum_{k,j}|\nu(F_k\cap E_j)|\\
&=\sum_j\sum_k|\nu(F_k\cap E_j)|\\
&\leq\sum_j|\nu|(E_j).
\end{align*}
Taking the supremum over partitions $\{F_k\}$ yields the second inequality.
\end{proof}

Define the \emph{positive variation} and \emph{negative variation} of $\nu$ by
\[
\nu^+=\frac12(|\nu|+\nu),
\qquad
\nu^-=\frac12(|\nu|-\nu).
\]
Then $\nu^+$ and $\nu^-$ are positive measures, and
\[
\nu=\nu^+-\nu^- ,
\qquad
|\nu|=\nu^++\nu^-.
\]
If $\nu(E)=\infty$, then $|\nu|(E)=\infty$ and $\nu^-(E)$ is defined to be zero. A signed measure is called $\sigma$-finite if $|\nu|$ is $\sigma$-finite. Since
\[
-|\nu|\leq\nu\leq|\nu|,
\]
$\sigma$-finiteness of $\nu$ implies that of $\nu^+$ and $\nu^-$.

\subsection{Absolute continuity}

Given two measures defined on a common $\sigma$-algebra, two extreme situations may occur: the measures may be supported on separate parts of the space, or the support of one may be an essential part of the support of the other. The Lebesgue--Radon--Nikodym theorem states that, in a precise sense, the relation between arbitrary measures is a combination of these possibilities.

\subsubsection*{Mutually singular and absolutely continuous measures}

Two signed measures $\nu$ and $\mu$ on $(X,\mathcal M)$ are \emph{mutually singular} if there are disjoint measurable subsets $A$ and $B$ such that
\[
\nu(E)=\nu(A\cap E),
\qquad
\mu(E)=\mu(B\cap E)
\qquad\text{for all }E\in\mathcal M.
\]
We write $\nu\perp\mu$.

If $\nu$ is a signed measure and $\mu$ a positive measure, we say that $\nu$ is \emph{absolutely continuous} with respect to $\mu$ if
\begin{equation}
\nu(E)=0\qquad\text{whenever }E\in\mathcal M\text{ and }\mu(E)=0.
\tag{12}
\end{equation}
We write $\nu\ll\mu$.

An important example is given by integration. If $f\in L^1(X,\mu)$, or more generally is integrable in the extended sense, define
\begin{equation}
\nu(E)=\int_E f\,d\mu.
\tag{13}
\end{equation}
Then $\nu\ll\mu$. We use the shorthand $d\nu=f\,d\mu$.

In the one-dimensional case considered in Chapter 3, one had the stronger condition
\begin{equation}
\text{for each }\epsilon>0\text{ there is }\delta>0\text{ such that }\mu(E)<\delta
\Longrightarrow |\nu(E)|<\epsilon.
\tag{14}
\end{equation}

\begin{proposition}
Assertion (14) implies (12). Conversely, if $|\nu|$ is a finite measure, then (12) implies (14).
\end{proposition}

\begin{proof}
The first implication is immediate. For the converse it suffices to consider positive $\nu$, replacing $\nu$ by $|\nu|$. Suppose (14) fails for some fixed $\epsilon>0$. For each $n$ choose $E_n$ with $\mu(E_n)<2^{-n}$ but $\nu(E_n)\geq\epsilon$. Put
\[
E_n^*=\bigcup_{k\geq n}E_k,
\qquad
E^*=\bigcap_{n=1}^{\infty}E_n^*=\limsup_{n\to\infty}E_n.
\]
Then $\mu(E_n^*)\leq 2^{-(n-1)}$, so $\mu(E^*)=0$. On the other hand $\nu(E_n^*)\geq\epsilon$, and since $\nu$ is finite and $E_n^*\downarrow E^*$,
\[
\nu(E^*)=\lim_{n\to\infty}\nu(E_n^*)\geq\epsilon,
\]
contradicting absolute continuity.
\end{proof}

\begin{theorem}[Lebesgue--Radon--Nikodym]
Suppose $\mu$ is a $\sigma$-finite positive measure on $(X,\mathcal M)$ and $\nu$ a $\sigma$-finite signed measure on $\mathcal M$. Then there exist unique signed measures $\nu_a$ and $\nu_s$ on $\mathcal M$ such that
\[
\nu_a\ll\mu,
\qquad
\nu_s\perp\mu,
\qquad
\nu=\nu_a+\nu_s.
\]
In addition, $\nu_a$ takes the form $d\nu_a=f\,d\mu$; that is,
\[
\nu_a(E)=\int_E f(x)\,d\mu(x)
\]
for some extended $\mu$-integrable function $f$.
\end{theorem}

If $\nu\ll\mu$, the theorem gives $d\nu=f\,d\mu$, a generalization of the representation obtained in Chapter 3.

\begin{proof}
We start with the case when both $\nu$ and $\mu$ are positive and finite. Let $\rho=\nu+\mu$, and consider on $L^2(X,\rho)$ the linear functional
\[
\ell(\psi)=\int_X\psi(x)\,d\nu(x).
\]
It is bounded because
\begin{align*}
|\ell(\psi)|
&\leq\int_X|\psi|\,d\nu
\leq\int_X|\psi|\,d\rho\\
&\leq\rho(X)^{1/2}\left(\int_X|\psi|^2\,d\rho\right)^{1/2}.
\end{align*}
By the Riesz representation theorem there is $g\in L^2(X,\rho)$ such that
\begin{equation}
\int_X\psi\,d\nu=\int_X\psi g\,d\rho
\qquad\text{for all }\psi\in L^2(X,\rho).
\tag{15}
\end{equation}
Taking $\psi=\chi_E$ and using $\nu\leq\rho$ gives $0\leq g\leq1$ a.e.; modifying on a null set, assume this everywhere. Since $d\rho=d\nu+d\mu$, equation (15) can be rewritten as
\begin{equation}
\int_X\psi(1-g)\,d\nu=\int_X\psi g\,d\mu.
\tag{16}
\end{equation}

Set
\[
A=\{x:0\leq g(x)<1\},
\qquad
B=\{x:g(x)=1\},
\]
and define
\[
\nu_a(E)=\nu(A\cap E),
\qquad
\nu_s(E)=\nu(B\cap E).
\]
Taking $\psi=\chi_B$ in (16) gives $\mu(B)=0$, so $\nu_s\perp\mu$.

Next take $\psi=\chi_E(1+g+\cdots+g^n)$ in (16):
\begin{equation}
\int_E(1-g^{n+1})\,d\nu
=\int_E g(1+g+\cdots+g^n)\,d\mu.
\tag{17}
\end{equation}
The left-hand side tends by dominated convergence to $\nu(A\cap E)=\nu_a(E)$, while the right-hand side tends to
\[
\int_E\frac{g}{1-g}\,d\mu.
\]
Thus
\[
\nu_a(E)=\int_E f\,d\mu,
\qquad
f=\frac{g}{1-g}.
\]
Since $\nu_a(X)\leq\nu(X)<\infty$, $f\in L^1(X,\mu)$.

If $\mu$ and $\nu$ are merely $\sigma$-finite and positive, decompose $X$ into measurable pieces $E_j$ on which both measures are finite, apply the preceding construction to the restricted measures, and sum the resulting absolutely continuous and singular parts. For a signed $\nu$, apply the argument separately to $\nu^+$ and $\nu^-$.

For uniqueness, suppose also $\nu=\nu_a'+\nu_s'$ with $\nu_a'\ll\mu$ and $\nu_s'\perp\mu$. Then
\[
\nu_a-\nu_a'=\nu_s'-\nu_s.
\]
The left side is absolutely continuous with respect to $\mu$, while the right side is singular with respect to $\mu$. Hence both sides vanish.
\end{proof}

\setcounter{section}{5}
\setcounter{theorem}{0}
\section*{5* Ergodic theorems}
\addcontentsline{toc}{section}{5* Ergodic theorems}

Ergodic theory had its beginnings in certain problems in statistical mechanics studied in the late nineteenth century. Since then it has grown rapidly and has gained wide influence in a number of mathematical disciplines, in particular those related to dynamical systems and probability theory. It is not our purpose to try to give an account of this broad and fascinating theory. Rather, we restrict our presentation to some of the basic limit theorems that lie at its foundation. These theorems are most naturally formulated in the general context of abstract measure spaces, and thus for us they serve as excellent illustrations of the general framework developed in this chapter.

The setting is a $\sigma$-finite measure space $(X,\mathcal M,\mu)$ endowed with a mapping $\tau:X\to X$ such that whenever $E$ is measurable, so is $\tau^{-1}(E)$, and
\[
\mu(\tau^{-1}(E))=\mu(E).
\]
Such a mapping is called a \emph{measure-preserving transformation}. If in addition $\tau$ is a bijection and $\tau^{-1}$ is also measure-preserving, then $\tau$ is a \emph{measure-preserving isomorphism}.

If $\tau$ is measure-preserving, then $f(\tau(x))$ is measurable whenever $f$ is measurable, and integrable whenever $f$ is integrable; moreover
\begin{equation}
\int_X f(\tau(x))\,d\mu(x)=\int_X f(x)\,d\mu(x).
\tag{18}
\end{equation}
Indeed, the assertion is immediate for characteristic functions, since $\chi_E(\tau(x))=\chi_{\tau^{-1}(E)}(x)$, and then follows by the usual limiting arguments. Equivalently, for real-valued measurable $f$ and every real $\alpha$,
\[
\mu(\{x:f(x)>\alpha\})=\mu(\{x:f(\tau(x))>\alpha\}).
\]

Before proceeding further, we describe several examples.
\begin{enumerate}
\item[(i)] Let $X=\mathbb Z$ with counting measure and define $\tau(n)=n+1$. This is a measure-preserving isomorphism.
\item[(ii)] Let $X=\mathbb R^d$ with Lebesgue measure and $\tau(x)=x+h$ for fixed $h\in\mathbb R^d$. This is a measure-preserving isomorphism.
\item[(iii)] Let $X=\mathbb R/\mathbb Z$, identified with $(0,1]$ and equipped with Lebesgue measure. For any real $\alpha$, the translation $x\mapsto x+\alpha\pmod{\mathbb Z}$ is measure-preserving; geometrically it is a rotation of the circle by angle $2\pi\alpha$.
\item[(iv)] Let $X=(0,1]$ with Lebesgue measure and let $\tau(x)=2x\pmod1$, the doubling map. Any measurable $E$ has two pre-images, one in $(0,1/2]$ and one in $(1/2,1]$, each of measure $\mu(E)/2$. Thus $\tau$ is measure-preserving but is not an isomorphism because it is not injective.
\item[(v)] Let $X=[0,1)$ and define $\tau(x)=\{1/x\}$, the fractional part of $1/x$, with $\tau(0)=0$. Gauss observed that the measure
\[
d\mu(x)=\frac{dx}{1+x}
\]
is preserved by $\tau$. Every $x\in(0,1)$ has infinitely many pre-images $1/(x+k)$, $k\geq1$.
\end{enumerate}

\begin{figure}[ht]
\centering
\fbox{\parbox{0.78\textwidth}{\centering Source Figure 1: pre-images $E_1$ and $E_2$ under the doubling map.}}
\caption{Pre-images under the doubling map.}
\end{figure}

These notions abstract the idea of a dynamical system. The space $X$ represents the totality of states, a point $x\in X$ a particular state, and $\tau:X\to X$ the transformation after one unit of time. The iterates $\tau^n$ describe the evolution after $n$ units of time, and one is led to study the averages
\begin{equation}
A_n(f)(x)=\frac1n\sum_{k=0}^{n-1}f(\tau^k(x))
\tag{19}
\end{equation}
and their limits as $n\to\infty$.

\subsection*{5.1 Mean ergodic theorem}
\addcontentsline{toc}{subsection}{5.1 Mean ergodic theorem}

The first theorem dealing with the averages (19) is purely Hilbert-space in character. For the specific application, take $H=L^2(X,\mathcal M,\mu)$ and define
\begin{equation}
T(f)(x)=f(\tau(x)).
\tag{20}
\end{equation}
Then $T$ is an isometry,
\begin{equation}
\|Tf\|=\|f\|,
\tag{21}
\end{equation}
by (18) applied to $|f|^2$. If $\tau$ is a measure-preserving isomorphism, then $T$ is unitary, but this is not assumed.

Let
\[
S=\{f\in H:Tf=f\}
\]
be the closed subspace of invariant vectors, and let $P$ be the orthogonal projection onto $S$.

\begin{theorem}[Mean ergodic theorem]
Suppose $T$ is an isometry of the Hilbert space $H$, and let $P$ be the orthogonal projection on the subspace of invariant vectors of $T$. Let
\[
A_n=\frac1n(I+T+T^2+\cdots+T^{n-1}).
\]
Then for each $f\in H$, $A_n(f)$ converges to $P(f)$ in norm as $n\to\infty$.
\end{theorem}

Together with $S$, consider
\[
S^*=\{f\in H:T^*f=f\},
\qquad
S_1=\{f\in H:f=g-Tg,\ g\in H\},
\]
and let $\overline{S_1}$ denote the closure of $S_1$.

\begin{lemma}
The following relations hold:
\begin{enumerate}
\item[(i)] $S=S^*$.
\item[(ii)] The orthogonal complement of $\overline{S_1}$ is $S$.
\end{enumerate}
\end{lemma}

\begin{proof}
Since $T$ is an isometry, $(Tf,Tg)=(f,g)$, hence $T^*T=I$. If $Tf=f$, then $T^*f=T^*Tf=f$, so $S\subset S^*$. Conversely, if $T^*f=f$, then
\[
(Tf,f)=(f,T^*f)=\|f\|^2,
\]
while $\|Tf\|=\|f\|$. Equality in Cauchy--Schwarz implies $Tf=cf$, and the displayed identity forces $c=1$. Thus $S^*\subset S$.

Next, $f\perp S_1$ exactly when
\[
(f,g-Tg)=0\qquad\text{for all }g\in H,
\]
that is, when $(f-T^*f,g)=0$ for all $g$, hence $f=T^*f$. By (i) this means $f\in S$.
\end{proof}

\begin{proof}[Proof of Theorem 5.1]
Write $f=f_0+f_1$, where $f_0\in S$ and $f_1\in\overline{S_1}$. Fix $\epsilon>0$ and choose $f_1'\in S_1$ with $\|f_1-f_1'\|<\epsilon$. Then
\begin{equation}
A_n(f)=A_n(f_0)+A_n(f_1')+A_n(f_1-f_1').
\tag{22}
\end{equation}
Since $Tf_0=f_0$,
\[
A_n(f_0)=f_0=P(f).
\]
Write $f_1'=g-Tg$. Then
\[
A_n(f_1')
=\frac1n\sum_{k=0}^{n-1}T^k(I-T)g
=\frac1n(g-T^ng),
\]
which tends to $0$ in norm because $T$ is an isometry. Finally,
\[
\|A_n(f_1-f_1')\|
\leq\frac1n\sum_{k=0}^{n-1}\|T^k(f_1-f_1')\|
<\epsilon.
\]
Thus $\limsup_{n\to\infty}\|A_n(f)-P(f)\|\leq\epsilon$, proving the theorem.
\end{proof}

\subsection*{5.2 Maximal ergodic theorem}
\addcontentsline{toc}{subsection}{5.2 Maximal ergodic theorem}

We now turn to almost everywhere convergence of the averages. As in the differentiation theorems of Chapter 3, the key is a maximal function, here defined by
\begin{equation}
f^*(x)=\sup_{1\leq m<\infty}\frac1m\sum_{k=0}^{m-1}|f(\tau^k(x))|.
\tag{23}
\end{equation}

\begin{theorem}[Maximal ergodic theorem]
Whenever $f\in L^1(X,\mu)$, the maximal function $f^*(x)$ is finite for almost every $x$. Moreover, there is a universal constant $A$ such that
\begin{equation}
\mu(\{x:f^*(x)>\alpha\})\leq\frac{A}{\alpha}\|f\|_{L^1(X,\mu)}
\qquad\text{for all }\alpha>0.
\tag{24}
\end{equation}
\end{theorem}

The proof below deduces this from the one-dimensional Hardy--Littlewood maximal theorem and gives $A=6$. The exact value is not relevant. Since $f^*$ is the supremum of countably many measurable functions, it is measurable. We may assume $f\geq0$.

\paragraph{Step 1. The case $X=\mathbb Z$ and $\tau(n)=n+1$.}
For a function $f$ on $\mathbb Z$, define its extension $\widetilde f$ to $\mathbb R$ by
\[
\widetilde f(x)=f(n),\qquad n\leq x<n+1.
\]
For $E\subset\mathbb Z$, let $\widetilde E=\bigcup_{n\in E}[n,n+1)$. Then
\[
m(\widetilde E)=\#(E),
\qquad
\|\widetilde f\|_{L^1(\mathbb R)}=\|f\|_{L^1(\mathbb Z)}.
\]
Moreover
\[
\sum_{k=0}^{m-1}f(n+k)=\int_0^m\widetilde f(n+t)\,dt.
\]
If $x\in[n,n+1)$, then
\[
\int_0^m\widetilde f(n+t)\,dt
\leq\int_{-1}^{m}\widetilde f(x+t)\,dt,
\]
so
\begin{equation}
f^*(n)\leq2(\widetilde f)^*(x)
\qquad\text{whenever }x\in[n,n+1).
\tag{25}
\end{equation}
The one-dimensional maximal theorem therefore gives
\begin{equation}
\#\{n:f^*(n)>\alpha\}
\leq\frac6\alpha\|f\|_{L^1(\mathbb Z)}.
\tag{26}
\end{equation}

\paragraph{Step 2. The general case.}
For a positive integer $N$, define the truncated maximal function
\[
f_N^*(x)=\sup_{1\leq m\leq N}\frac1m\sum_{k=0}^{m-1}f(\tau^k(x)).
\]
It suffices to prove
\begin{equation}
\mu\{x:f_N^*(x)>\alpha\}\leq\frac{A}{\alpha}\|f\|_{L^1(X,\mu)}
\tag{27}
\end{equation}
with $A$ independent of $N$.

Temporarily write $M$ for the discrete maximal operator on $\mathbb Z$. Define $F$ on $X\times\mathbb Z$ by
\[
F(x,n)=
\begin{cases}
f(\tau^n(x)),&n\geq0,\\
0,&n<0.
\end{cases}
\]
Then
\[
A_m(f)(\tau^n(x))=\frac1m\sum_{k=0}^{m-1}F(x,n+k).
\]
Fix a large positive integer $a$, put $b=a+N$, and truncate $F$ by $F_b(x,n)=F(x,n)$ for $n<b$ and $0$ otherwise. For $m\leq N$ and $n<a$,
\[
A_m(f)(\tau^n(x))=\frac1m\sum_{k=0}^{m-1}F_b(x,n+k),
\]
hence
\begin{equation}
f_N^*(\tau^n(x))\leq M(F_b)(x,n)
\qquad(n<a).
\tag{28}
\end{equation}
Let $E_\alpha=\{x:f_N^*(x)>\alpha\}$. Measure preservation gives
\[
\mu\{x:f_N^*(\tau^n(x))>\alpha\}=\mu(E_\alpha).
\]
On $X\times\mathbb Z$, (28) and (26) imply
\begin{align*}
a\mu(E_\alpha)
&\leq\int_X\#\{n:M(F_b)(x,n)>\alpha\}\,d\mu(x)\\
&\leq\frac{A}{\alpha}\int_X\sum_{n=0}^{b-1}f(\tau^n(x))\,d\mu(x)\\
&=\frac{A}{\alpha}b\|f\|_{L^1(X)}.
\end{align*}
Thus
\[
\mu(E_\alpha)\leq\frac{A}{\alpha}\left(1+\frac Na\right)\|f\|_{L^1(X)},
\]
and letting $a\to\infty$ gives (27). Finally $N\to\infty$ proves the theorem.

\subsection*{5.3 Pointwise ergodic theorem}
\addcontentsline{toc}{subsection}{5.3 Pointwise ergodic theorem}

The pointwise, or individual, ergodic theorem combines the ideas of the preceding two theorems. Here it is convenient to assume the measure space is finite and normalized so that $\mu(X)=1$.

\begin{theorem}[Pointwise ergodic theorem]
Suppose $f$ is integrable over $X$. Then for almost every $x\in X$ the averages
\[
A_m(f)(x)=\frac1m\sum_{k=0}^{m-1}f(\tau^k(x))
\]
converge to a limit as $m\to\infty$.
\end{theorem}

\begin{corollary}
If this limit is denoted by $P'(f)$, then
\[
\int_X|P'(f)(x)|\,d\mu(x)
\leq\int_X|f(x)|\,d\mu(x).
\]
Moreover $P'(f)=P(f)$ whenever $f\in L^2(X,\mu)$.
\end{corollary}

\begin{proof}
Since $\mu(X)=1$, we have $L^2(X,\mu)\subset L^1(X,\mu)$ and $\|f\|_{L^1}\leq\|f\|_{L^2}$. Moreover $L^2$ is dense in $L^1$: if $f\in L^1$, put $f_n=f$ on $\{|f|\leq n\}$ and $f_n=0$ elsewhere; then $f_n\in L^2$ and $\|f-f_n\|_{L^1}\to0$.

Starting with $f\in L^1$ and $\epsilon>0$, write
\[
f=F+H,
\qquad
\|H\|_{L^1}<\epsilon,
\]
where
\[
F=F_0+(I-T)G,
\qquad
F_0,G\in L^2,
\qquad
TF_0=F_0.
\]
To obtain this, first write $f=f'+h'$ with $f'\in L^2$ and $\|h'\|_{L^1}<\epsilon/2$. Since $S$ and $\overline{S_1}$ are orthogonal complements in $L^2$, approximate the $\overline{S_1}$ component of $f'$ by an element $(I-T)G\in S_1$ with $L^2$ error $<\epsilon/2$. Because $\|h\|_{L^1}\leq\|h\|_{L^2}$ on a probability space, the required decomposition follows.

Now
\[
A_m(F)=F_0+\frac1m(I-T^m)G.
\]
The term $m^{-1}G$ tends to zero, and $m^{-1}T^mG$ tends to zero almost everywhere. Indeed,
\[
\int_X\sum_{m=1}^{\infty}\frac{|T^mG|^2}{m^2}\,d\mu
=\sum_{m=1}^{\infty}\frac{\|T^mG\|_2^2}{m^2}
=\|G\|_2^2\sum_{m=1}^{\infty}\frac1{m^2}<\infty.
\]
Thus $A_m(F)(x)$ converges for almost every $x$.

For $\alpha>0$, let
\[
E_\alpha=
\left\{x:\lim_{N\to\infty}\sup_{n,m\geq N}|A_n(f)(x)-A_m(f)(x)|>\alpha\right\}.
\]
Since the averages of $F$ converge a.e., almost every point of $E_\alpha$ lies in the set where
\[
2\sup_m|A_m(H)(x)|>\alpha.
\]
By the maximal ergodic theorem,
\[
\mu(E_\alpha)\leq\frac{2A}{\alpha}\|H\|_{L^1}
\leq\frac{2A\epsilon}{\alpha}.
\]
Since $\epsilon$ is arbitrary, $\mu(E_\alpha)=0$. Hence $A_m(f)(x)$ is Cauchy for almost every $x$, proving the theorem.

If $f\in L^2$, the mean ergodic theorem gives $A_m(f)\to P(f)$ in $L^2$, so a subsequence converges a.e. to $P(f)$; hence $P'(f)=P(f)$. For arbitrary integrable $f$,
\[
\int_X|A_m(f)|\,d\mu
\leq\frac1m\sum_{k=0}^{m-1}\int_X|f(\tau^k(x))|\,d\mu(x)
=\int_X|f(x)|\,d\mu(x).
\]
Fatou's lemma then gives
\[
\int_X|P'(f)|\,d\mu\leq\int_X|f|\,d\mu.
\]
\end{proof}

The theorem and corollary remain valid without the finite-measure assumption; the necessary modifications are outlined in the exercises of the source.

\subsection*{5.4 Ergodic measure-preserving transformations}
\addcontentsline{toc}{subsection}{5.4 Ergodic measure-preserving transformations}

A measure-preserving transformation $\tau$ of $X$ is \emph{ergodic} if whenever $E$ is measurable and invariant in the sense that $E$ and $\tau^{-1}(E)$ differ by a null set, then either $E$ or $E^c$ has measure zero.

Equivalently, a measurable function $f$ is called invariant if $f(x)=f(\tau(x))$ for almost every $x$, and $\tau$ is ergodic exactly when every invariant function is equivalent to a constant. Indeed, if $f$ is real-valued and invariant, each set
\[
E_a=\{x:f(x)>a\}
\]
is invariant. Ergodicity forces $\mu(E_a)=0$ or $\mu(E_a^c)=0$ for every $a$, which implies $f$ is constant a.e. The converse follows by applying this statement to characteristic functions.

Assume again that $\mu(X)=1$.

\begin{corollary}
Suppose $\tau$ is an ergodic measure-preserving transformation. For any integrable function $f$,
\[
\frac1m\sum_{k=0}^{m-1}f(\tau^k(x))
\longrightarrow
\int_Xf\,d\mu
\]
for almost every $x\in X$ as $m\to\infty$.
\end{corollary}

The result says that the ``time average'' of $f$ equals its ``space average.''

\begin{proof}
For $f\in L^2$, the mean ergodic theorem gives convergence to $P(f)$, where $P$ is the orthogonal projection onto invariant vectors. Under ergodicity the invariant vectors are exactly the constants, hence
\[
P(f)=\mathbf1(f,\mathbf1)=\left(\int_Xf\,d\mu\right)\mathbf1.
\]
For $f\in L^1$, write $f=g+h$ with $g\in L^2$ and $\|h\|_{L^1}<\epsilon$. Then $P'(g)=P(g)$ and, by the preceding corollary, $\|P'(h)\|_{L^1}\leq\|h\|_{L^1}$. Hence
\[
\left\|P'(f)-\int_Xf\,d\mu\right\|_{L^1}<2\epsilon.
\]
Since $\epsilon$ is arbitrary, $P'(f)$ is the constant $\int_Xf\,d\mu$.
\end{proof}

\subsubsection*{a) Rotations of the circle}

On $\mathbb R/\mathbb Z$ with induced Lebesgue measure, let
\[
\tau(x)=x+\alpha\pmod1.
\]
Then:
\begin{itemize}
\item $\tau$ is ergodic if and only if $\alpha$ is irrational.
\end{itemize}
If $\alpha$ is irrational, the equidistribution theorem gives
\begin{equation}
\frac1n\sum_{k=0}^{n-1}f(x+k\alpha)
\longrightarrow\int_0^1f(x)\,dx
\tag{29}
\end{equation}
for every $x$ whenever $f$ is continuous and periodic. The argument first verifies this for $f(x)=e^{2\pi imx}$, $m\in\mathbb Z$, then for trigonometric polynomials, and finally uses uniform approximation of continuous periodic functions.

If $P$ is the projection onto invariant $L^2$ functions, (29) and the mean ergodic theorem show that $P$ maps continuous periodic functions to constants. Since these functions are dense in $L^2$, all invariant $L^2$ functions are constants, so $\tau$ is ergodic.

If $\alpha=p/q$ is rational, choose $E_0\subset(0,1/q)$ with $0<m(E_0)<1/q$ and let
\[
E=\bigcup_{r=0}^{q-1}(E_0+r/q).
\]
Then $E$ is invariant under $x\mapsto x+p/q$, while $0<m(E)<1$, so $\tau$ is not ergodic.

Property (29) is stronger than ergodicity: it implies that Lebesgue measure is \emph{uniquely ergodic} for this transformation. This means that any Borel probability measure $\nu$ preserved by $\tau$ must equal Lebesgue measure. Indeed, applying the mean ergodic theorem in $L^2(X,\nu)$ and using (29) shows that
\[
\int_0^1f(x)\,dx=\int_Xf\,d\nu
\]
for every continuous periodic $f$, and a limiting argument yields equality of the measures.

\subsubsection*{b) The doubling mapping}

Consider $\tau(x)=2x\pmod1$ on $(0,1]$ with Lebesgue measure. This transformation is ergodic and satisfies the stronger property of \emph{mixing}. A measure-preserving transformation is mixing if for all measurable $E,F$,
\begin{equation}
\mu(\tau^{-n}(E)\cap F)\longrightarrow\mu(E)\mu(F)
\qquad(n\to\infty).
\tag{30}
\end{equation}
In probabilistic language, $\tau^{-n}(E)$ and $F$ become asymptotically independent.

A condition equivalent to (30) is
\begin{equation}
(T^nf,g)\longrightarrow(f,\mathbf1)(\mathbf1,g)
\qquad(n\to\infty),
\tag{31}
\end{equation}
for all $f,g\in L^2(X,\mu)$. One implication is immediate by taking characteristic functions; the converse follows by approximation.

Mixing implies ergodicity. Indeed,
\[
(A_n(f),g)=\frac1n\sum_{k=0}^{n-1}(T^kf,g)
\longrightarrow(f,\mathbf1)(\mathbf1,g).
\]
Thus the mean-ergodic projection $P$ is the orthogonal projection onto constants.

The doubling map is mixing. If
\[
f(x)=e^{2\pi imx},
\qquad
g(x)=e^{2\pi ikx},
\]
then $(f,\mathbf1)(\mathbf1,g)=0$ unless $m=k=0$. On the other hand
\[
(T^nf,g)=\int_0^1e^{2\pi im2^nx}e^{-2\pi ikx}\,dx,
\]
which vanishes for sufficiently large $n$ unless $m=k=0$, in which case it equals $1$. Thus (31) holds for exponentials, hence by linearity for trigonometric polynomials, and by $L^2$ approximation for arbitrary $f,g\in L^2$.

Irrational rotations, although ergodic, are not mixing: if $f=g=e^{2\pi imx}$ with $m\neq0$, then
\[
(T^nf,g)=e^{2\pi inm\alpha}(f,g)=e^{2\pi inm\alpha},
\]
while $(f,\mathbf1)=(\mathbf1,g)=0$.

Finally, the doubling map is not uniquely ergodic. Besides Lebesgue measure, the probability measure concentrated at $1$ is also preserved by $\tau$.

\setcounter{section}{6}
\setcounter{theorem}{0}
\section*{6* Appendix: the spectral theorem}
\addcontentsline{toc}{section}{6* Appendix: the spectral theorem}

The purpose of this appendix is to present an outline of the proof of the spectral theorem for bounded symmetric operators on a Hilbert space. Details that are not central to the proof of the theorem are left to the reader. The theorem provides an application of the ideas related to Lebesgue--Stieltjes integrals treated in this chapter.

\subsection*{6.1 Statement of the theorem}
\addcontentsline{toc}{subsection}{6.1 Statement of the theorem}

A basic notion is that of a \emph{spectral resolution} (or \emph{spectral family}) on a Hilbert space $H$. This is a function $\lambda\mapsto E(\lambda)$ from $\mathbb R$ to orthogonal projections on $H$ satisfying:
\begin{enumerate}
\item[(i)] $E(\lambda)$ is increasing in the sense that $\|E(\lambda)f\|$ is an increasing function of $\lambda$ for every $f\in H$.
\item[(ii)] There is an interval $[a,b]$ such that $E(\lambda)=0$ if $\lambda<a$, and $E(\lambda)=I$ if $\lambda\geq b$.
\item[(iii)] $E(\lambda)$ is right-continuous, that is, for every $\lambda$,
\[
\lim_{\substack{\mu\to\lambda\\\mu>\lambda}}E(\mu)f=E(\lambda)f
\qquad\text{for every }f\in H.
\]
\end{enumerate}
Property (i) is equivalent, for $\mu>\lambda$, to each of: the range of $E(\mu)$ contains the range of $E(\lambda)$; $E(\mu)E(\lambda)=E(\lambda)$; and $E(\mu)-E(\lambda)$ is an orthogonal projection.

Given a spectral resolution and $f\in H$, the function
\[
\lambda\longmapsto(E(\lambda)f,f)=\|E(\lambda)f\|^2
\]
is increasing. By the polarization identity, for every $f,g\in H$ the function
\[
F(\lambda)=(E(\lambda)f,g)
\]
is of bounded variation and right-continuous.

\begin{theorem}[Spectral theorem]
Suppose $T$ is a bounded symmetric operator on a Hilbert space $H$. Then there exists a spectral resolution $\{E(\lambda)\}$ such that
\[
T=\int_{a^-}^{b}\lambda\,dE(\lambda)
\]
in the sense that for every $f,g\in H$,
\begin{equation}
(Tf,g)
=\int_{a^-}^{b}\lambda\,d(E(\lambda)f,g)
=\int_{a^-}^{b}\lambda\,dF(\lambda).
\tag{32}
\end{equation}
The integral on the right is a Lebesgue--Stieltjes integral.
\end{theorem}

For a compact symmetric operator with orthonormal basis of eigenvectors $\{\varphi_k\}$ and eigenvalues $\{\lambda_k\}$, the spectral resolution is
\[
E(\lambda)f\sim\sum_{\lambda_k\leq\lambda}(f,\varphi_k)\varphi_k.
\]
Then
\[
\|E(\lambda)f\|^2=\sum_{\lambda_k\leq\lambda}|(f,\varphi_k)|^2,
\]
and $F(\lambda)=(E(\lambda)f,g)$ is a pure jump function.

\subsection*{6.2 Positive operators}
\addcontentsline{toc}{subsection}{6.2 Positive operators}

We say that $T$ is \emph{positive}, written $T\geq0$, if $T$ is symmetric and $(Tf,f)\geq0$ for every $f\in H$. We write $T_1\geq T_2$ if $T_1-T_2\geq0$. For two orthogonal projections, $E_2\geq E_1$ if and only if $\|E_2f\|\geq\|E_1f\|$ for all $f$, equivalently if the range of $E_2$ contains that of $E_1$. If $S$ is symmetric, then $S^2$ is positive.

For a symmetric $T$, let
\begin{equation}
a=\min_{\|f\|\leq1}(Tf,f),
\qquad
b=\max_{\|f\|\leq1}(Tf,f).
\tag{33}
\end{equation}

\begin{proposition}
Suppose $T$ is symmetric. Then $\|T\|\leq M$ if and only if
\[
-MI\leq T\leq MI.
\]
As a result,
\[
\|T\|=\max(|a|,|b|).
\]
\end{proposition}

\begin{proposition}
Suppose $T$ is positive. Then there exists a symmetric operator $S=T^{1/2}$ such that $S^2=T$, and $S$ commutes with every operator that commutes with $T$.
\end{proposition}

\begin{proof}
After multiplying by a positive scalar, assume $\|T\|\leq1$. Consider the binomial expansion
\[
(1-t)^{1/2}=\sum_{k=0}^{\infty}b_kt^k,
\qquad |t|<1.
\]
The coefficients are real and $\sum|b_k|<\infty$. In fact $b_0=1$, $b_1=-1/2$, $b_2=-1/8$, and for $k\geq2$,
\[
b_k=-\frac{1}{2}\frac{1}{2}\cdots\frac{k-3/2}{k!},
\]
so $b_k=O(k^{-3/2})$.

Let
\[
s_n(t)=\sum_{k=0}^{n}b_kt^k.
\]
Then
\begin{equation}
s_n(t)^2-(1-t)=\sum_{k=0}^{2n}c_k^{(n)}t^k,
\tag{34}
\end{equation}
where
\[
\sum_{k=0}^{2n}|c_k^{(n)}|\to0
\qquad(n\to\infty).
\]
Indeed, if $r_n(t)=\sum_{k=n+1}^{\infty}b_kt^k$, then
\[
s_n(t)^2-(1-t)=r_n(t)^2-2(1-t)^{1/2}r_n(t),
\]
and comparison of coefficients gives the asserted estimate.

Set $T_1=I-T$. Then $0\leq T_1\leq I$ and $\|T_1\|\leq1$. Define
\[
S_n=s_n(T_1)=\sum_{k=0}^{n}b_kT_1^k.
\]
Since $\sum|b_k|<\infty$, $S_n$ converges in operator norm to an operator $S$. Each $S_n$ is symmetric, hence so is $S$. By (34),
\[
S_n^2-T=\sum_{k=0}^{2n}c_k^{(n)}T_1^k,
\]
and therefore $\|S_n^2-T\|\to0$. Thus $S^2=T$. Finally, any operator commuting with $T$ commutes with every polynomial in $T$, hence with all $S_n$ and with $S$.
\end{proof}

\begin{proposition}
If $T_1$ and $T_2$ are positive operators that commute, then $T_1T_2$ is positive.
\end{proposition}

Indeed, if $S=T_1^{1/2}$, then $S$ commutes with $T_2$ and
\[
T_1T_2=ST_2S.
\]
Hence
\[
(T_1T_2f,f)=(T_2Sf,Sf)\geq0.
\]

\begin{proposition}
Suppose $T$ is symmetric and $a,b$ are given by (33). If $p(t)$ is a real polynomial that is non-negative on $[a,b]$, then $p(T)$ is positive.
\end{proposition}

The proof factors $p$ into real linear factors corresponding to roots outside $[a,b]$, even-order factors for roots inside the interval, and quadratic factors corresponding to non-real conjugate roots. Each corresponding operator factor is positive, and all commute, so the preceding proposition applies.

\begin{corollary}
If $p(t)$ is a real polynomial, then
\[
\|p(T)\|\leq\sup_{t\in[a,b]}|p(t)|.
\]
\end{corollary}

\begin{proposition}
Suppose $\{T_n\}$ is a sequence of positive operators with $T_n\geq T_{n+1}$ for all $n$. Then there is a positive operator $T$ such that
\[
T_nf\to Tf
\qquad\text{for every }f\in H.
\]
\end{proposition}

\begin{proof}
For fixed $f$, the sequence $(T_nf,f)$ is decreasing and convergent. We use the estimate: if $S\geq0$ and $\|S\|\leq M$, then
\begin{equation}
\|Sf\|^2\leq(Sf,f)^{1/2}M^{3/2}\|f\|.
\tag{35}
\end{equation}
Indeed, positivity of
\[
(S(tI+S)f,(tI+S)f)
=t^2(Sf,f)+2t(Sf,Sf)+(S^2f,Sf)
\]
for all real $t$ implies its discriminant is non-positive, hence
\[
\|Sf\|^4\leq(Sf,f)(S^2f,Sf),
\]
from which (35) follows.

Apply this to $S=T_n-T_m$ with $n\leq m$. The norms $\|T_n-T_m\|$ are uniformly bounded, while $((T_n-T_m)f,f)\to0$. Therefore $\|T_nf-T_mf\|\to0$, so $T_nf$ converges to a vector $Tf$. The limit operator $T$ is positive.
\end{proof}

\subsection*{6.3 Proof of the theorem}
\addcontentsline{toc}{subsection}{6.3 Proof of the theorem}

Starting with a symmetric operator $T$ and the interval $[a,b]$ from (33), we associate to suitable functions $\Phi$ on $[a,b]$ a symmetric operator $\Phi(T)$.

If $\Phi(t)=\sum_{k=0}^{n}c_kt^k$ is a real polynomial, define
\[
\Phi(T)=\sum_{k=0}^{n}c_kT^k.
\]
This association is a homomorphism:
\[
(\Phi_1+\Phi_2)(T)=\Phi_1(T)+\Phi_2(T),
\qquad
(\Phi_1\Phi_2)(T)=\Phi_1(T)\Phi_2(T).
\]
Moreover $\Phi(T)$ is symmetric.

Every real-valued continuous function $\Phi$ on $[a,b]$ can be uniformly approximated by polynomials $p_n$. The preceding corollary shows that $p_n(T)$ converges in operator norm to a limit independent of the approximating sequence; define this limit to be $\Phi(T)$. It is symmetric. If $\Phi\geq0$ on $[a,b]$, then $\Phi(T)\geq0$.

Finally, define $\Phi(T)$ when
\[
\Phi(t)=\lim_{n\to\infty}\Phi_n(t)
\]
for a decreasing sequence of positive continuous functions $\Phi_n$. Proposition 6.7 guarantees that $\Phi_n(T)$ has a strong limit. This limit is independent of the particular decreasing approximation: if $\Phi_n'$ is another such sequence, then for fixed $k$ and $\epsilon>0$, eventually
\[
\Phi_n'(t)\leq\Phi_k(t)+\epsilon,
\]
so
\[
\Phi_n'(T)\leq\Phi_k(T)+\epsilon I.
\]
Passing first $n\to\infty$, then $k\to\infty$, then $\epsilon\to0$, and reversing the roles of the two sequences, proves equality of the limits. Also $\Phi_1\leq\Phi_2$ implies $\Phi_1(T)\leq\Phi_2(T)$.

For each real $\lambda$ define
\[
\varphi_\lambda(t)=
\begin{cases}
1,&t\leq\lambda,\\
0,&\lambda<t.
\end{cases}
\]
This is the decreasing pointwise limit of continuous functions $\varphi_{\lambda,n}$ that equal $1$ for $t\leq\lambda$, equal $0$ for $t\geq\lambda+1/n$, and are linear in between. Define
\[
E(\lambda)=\varphi_\lambda(T).
\]
If $\lambda_1\leq\lambda_2$, then
\[
E(\lambda_1)E(\lambda_2)=E(\lambda_1).
\]
In particular $E(\lambda)^2=E(\lambda)$, and since $E(\lambda)$ is symmetric it is an orthogonal projection. Moreover
\[
\|E(\lambda_1)f\|
=\|E(\lambda_1)E(\lambda_2)f\|
\leq\|E(\lambda_2)f\|,
\]
so the family is increasing. Clearly $E(\lambda)=0$ for $\lambda<a$, and $E(\lambda)=I$ for $\lambda\geq b$.

The family is right-continuous. Fix $f\in H$ and $\epsilon>0$. Choose $n$ so that
\[
\|E(\lambda)f-\varphi_{\lambda,n}(T)f\|<\epsilon.
\]
For $\mu\to\lambda$, $\varphi_{\mu,n}\to\varphi_{\lambda,n}$ uniformly, so the corresponding operators converge in norm. Combining this with the monotonicity relations for $E(\mu)$ gives
\[
E(\mu)f\to E(\lambda)f
\qquad(\mu\downarrow\lambda).
\]
Thus $\{E(\lambda)\}$ is a spectral resolution.

It remains to verify (32). Let
\[
a=\lambda_0<\lambda_1<\cdots<\lambda_k=b
\]
be a partition of $[a,b]$ of mesh less than $\delta$. Pointwise on $[a,b]$,
\[
 t\leq
\sum_{j=1}^k\lambda_j\bigl(\varphi_{\lambda_j}(t)-\varphi_{\lambda_{j-1}}(t)\bigr)
+\lambda_0\varphi_{\lambda_0}(t)
\leq t+\delta.
\]
Applying the functional calculus gives
\[
T\leq
\sum_{j=1}^k\lambda_j\bigl(E(\lambda_j)-E(\lambda_{j-1})\bigr)
+\lambda_0E(\lambda_0)
\leq T+\delta I.
\]
Hence the operator norm difference is at most $\delta$, and therefore
\[
\left|
(Tf,f)-
\sum_{j=1}^k\lambda_j\int_{(\lambda_{j-1},\lambda_j]}d(E(\lambda)f,f)
-\lambda_0(E(\lambda_0)f,f)
\right|
\leq\delta\|f\|^2.
\]
Letting the mesh tend to zero gives
\[
(Tf,f)=\int_{a^-}^{b}\lambda\,d(E(\lambda)f,f),
\]
and polarization yields (32).

A similar argument shows that if $\Phi$ is continuous on $[a,b]$, then
\begin{equation}
(\Phi(T)f,g)=\int_{a^-}^{b}\Phi(\lambda)\,d(E(\lambda)f,g).
\tag{36}
\end{equation}
The same representation extends to complex-valued continuous $\Phi$ by separating real and imaginary parts, and to $\Phi$ obtained as decreasing pointwise limits of continuous functions.

\subsection*{6.4 Spectrum}
\addcontentsline{toc}{subsection}{6.4 Spectrum}

A bounded operator $S$ on $H$ is \emph{invertible} if it is a bijection and its inverse $S^{-1}$ is bounded. The \emph{spectrum} of $S$, denoted by $\sigma(S)$, is the set of complex numbers $z$ for which $S-zI$ is not invertible.

\begin{proposition}
If $T$ is symmetric, then $\sigma(T)$ is a closed subset of the interval $[a,b]$ given by (33).
\end{proposition}

\begin{proof}
If $z\notin[a,b]$, the function
\[
\Phi(t)=(t-z)^{-1}
\]
is continuous on $[a,b]$ and
\[
\Phi(T)(T-zI)=(T-zI)\Phi(T)=I,
\]
so $T-zI$ is invertible. Thus $\sigma(T)\subset[a,b]$.

To see that the spectrum is closed, suppose $T_0=T-\lambda_0I$ is invertible. Then
\[
T_0-\epsilon I=T_0(I-\epsilon T_0^{-1}).
\]
For $|\epsilon|<\|T_0^{-1}\|^{-1}$ the Neumann series converges, and the inverse of $I-\epsilon T_0^{-1}$ is
\[
\sum_{n=0}^{\infty}\epsilon^n(T_0^{-1})^n.
\]
Consequently $(T_0-\epsilon I)^{-1}$ exists, and its norm is bounded by
\begin{equation}
\frac{\|T_0^{-1}\|}{1-|\epsilon|\|T_0^{-1}\|}.
\tag{37}
\end{equation}
Thus the complement of the spectrum is open.
\end{proof}

\begin{proposition}
For each $f\in H$, the Lebesgue--Stieltjes measure corresponding to
\[
F(\lambda)=(E(\lambda)f,f)
\]
is supported on $\sigma(T)$.
\end{proposition}

Equivalently, $F(\lambda)$ is constant on each open interval of the complement of $\sigma(T)$.

\begin{proof}
Let $J$ be an open interval in the complement of $\sigma(T)$, let $x_0\in J$, and let $J_0$ be the subinterval centered at $x_0$ of length $2\epsilon$, with
\[
\epsilon<\|(T-x_0I)^{-1}\|^{-1}.
\]
If $z$ has non-zero imaginary part, then $(T-zI)^{-1}=\Phi_z(T)$ with $\Phi_z(t)=(t-z)^{-1}$. Therefore
\[
(T-zI)^{-1}(T-\overline z I)^{-1}
\]
corresponds to
\[
\Psi_z(t)=\frac1{|t-z|^2}.
\]
By the estimate (37) and the spectral representation (36), there is a constant $A'$ such that
\[
\int\frac{dF(\lambda)}{|\lambda-z|^2}\leq A'
\]
whenever $|x_0-z|<\epsilon$ and $\operatorname{Im}z\neq0$. Passing to real $x\in J_0$ gives the same estimate. Integrating in $x\in J_0$ and using
\[
\int_{J_0}\frac{dx}{|\lambda-x|^2}=\infty
\qquad\text{for every }\lambda\in J_0,
\]
forces
\[
\int_{J_0}dF(\lambda)=0.
\]
Thus $F$ is constant on $J_0$. Since $x_0$ was arbitrary in $J$, $F$ is constant throughout $J$.
\end{proof}

\end{document}
